\documentclass[12pt]{article}
\usepackage{latexsym, amssymb}

\usepackage{psfig}
\usepackage{graphics}
\usepackage{graphicx}


\parskip=2pt plus 1pt minus 1pt
%\parskip=3pt plus 1pt minus 1pt
\topmargin -0.5in
\textheight 9.5in
\oddsidemargin 0.25in
\evensidemargin 0.25in
\textwidth 6.25in

\newcommand{\spp}{\medskip}
\newcommand{\sppp}{\medskip\medskip}


\parindent0pt

\pagestyle{myheadings}
\markboth{\hbox to \textwidth{\hfill\huge\thepage}}{}

\newcommand{\comment}[1]{}
\newcommand{\psp}{\;}
\newcommand{\pspp}{\;\;}
\newcommand{\nsp}{}         %{\!}

\newcommand{\forconf}[1]{}
\newcommand{\forpaper}[1]{#1}
\newcommand{\forthesis}[1]{}


\def\restriction{\,|}  %I had an error - unknown

\newcommand{\R}{{\mathbb R}}  %ams bold
\newcommand{\X}{{\mathbb X}}  %ams bold

%\newcommand{\R}{{\bf R}}  
%\newcommand{\X}{{\bf X}}  

\newcommand{\maxint}{[0,\tmax)}
\newcommand{\tmax}{t_{{\rm max}}}
\newcommand{\states}{{\X}}
\newcommand{\controls}{{\mathbb U}} 
\newcommand{\maxints}{[0,\smax)}
\newcommand{\smax}{s_{{\rm max}}}
\newcommand{\ioss}{IOSS}  %{{\sc ioss}}  %%% CHANGED AT COPY EDIT--ADG %%%
\newcommand{\ios}{IOS}  %{{\sc ios}}	 %%% CHANGED AT COPY EDIT--ADG %%%
\newcommand{\newf}{\widehat{f}}
%%%%%%%%%%%%%%%%%%%

\newcommand{\io}[1]{@l_{#1}}
\newcommand{\dbio}[1]{\tilde @l_{#1}}
\newcommand{\rref}[1]{{\rm (\ref{#1})}}
\newcommand{\res}[1]{\!\restriction_{[#1]}}
\newcommand{\restr}[1]{\!\!\restriction_{#1}}
\newcommand{\kk}{!K}
\newcommand{\kl}{!K!L}
\newcommand{\ki}{{\cal K_\infty}}
\newcommand{\lem}{!L^m_{!8,e}}
\newcommand{\lep}{!L^p_{!8,e}}
\newcommand{\Rpos}{\R_{!>0}}
\newcommand{\mhat}{\widehat{@m}}
\newcommand{\xhat}{\widehat{@x}}
\newcommand{\uhat}{\widehat{u}}
\newcommand{\yhat}{\widehat{y}}
\newcommand{\htheta}{\hat@q}
\newcommand{\norm}[1]{\left\Vert#1\right\Vert}
\newcommand{\normi}[1]{\norm{#1}_{!8}}
\newcommand{\abs}[1]{\left\vert #1 \right\vert}

\newcommand{\gapply}{\odot}
\newcommand{\bapply}{\odot}
\newcommand{\traj}[2]{#2!.#1}

\newcommand{\oc}{ONAC}
\newcommand{\ocword}{ouput null asymptotically controllable under disturbances}
\newcommand{\oss}{OSS} %%{{\sc oss}}   %%% CHANGED AT COPY EDIT--ADG %%%
\newcommand{\iss}{ISS}  %{{\sc iss}}    %%% CHANGED AT COPY EDIT--ADG %%%
\newcommand{\ve}{@e}
\newcommand{\vf}{@j}
\newcommand{\ga}{{@g}}
\newcommand{\ts}{@l}
\newcommand{\I}{!I}
\newcommand{\J}{{\cal J}}
\newcommand{\D}{!D}
\newcommand{\E}{!E}
\newcommand{\T}{!T}
\newcommand{\U}{!U}
\newcommand{\V}{!V}
\newcommand{\K}{!K}
\newcommand{\M}{!M}
\newcommand{\B}{!B}
\newcommand{\G}{{\cal G}}
\newcommand{\GO}{{\cal G}_1}
\newcommand{\GT}{{\cal G}_2}
\newcommand{\A}{!A}
\newcommand{\N}{{\mathbb N}}
\newcommand{\edo}{\end{document}}
\renewcommand{\tilde}{\widetilde}

\overfullrule=0pt %no dark lines if overfull

\newtheorem{theorem}{Theorem}
\newtheorem{itlemma}{Lemma}[section] %number by section (set in \em by default)
%same numbering counter as itlemmas:
\newtheorem{itproposition}[itlemma]{Proposition}
\newtheorem{itcorollary}[itlemma]{Corollary}
\newtheorem{itremark}[itlemma]{Remark}
\newtheorem{itdefinition}[itlemma]{Definition}
\newtheorem{itexample}[itlemma]{Example}
\newenvironment{lemma}{\begin{itlemma}\rm}{\end{itlemma}} %no-italics
\newenvironment{remark}{\begin{itremark}\rm}{\end{itremark}} %no-italics
\newenvironment{corollary}{\begin{itcorollary}\rm}{\end{itcorollary}}
\newenvironment{proposition}{\begin{itproposition}\rm}{\end{itproposition}}
\newenvironment{definition}{\begin{itdefinition}\rm}{\end{itdefinition}}
\newenvironment{example}{\begin{itexample}\rm}{\end{itexample}}

\def\bi{\begin{itemize}}
\def\ei{\end{itemize}}
\def\ben{\begin{enumerate}}
\def\een{\end{enumerate}}
\def \beq {\begin{eqnarray}}
\def \eeq {\end{eqnarray}}
\def \beqn {\begin{eqnarray*}}
\def \eeqn {\end{eqnarray*}}
\def\st{\, | \,}
\newcommand{\twoif}[4]{
\left\{ \begin{array}{ll}#1&#2\\#3&#4\end{array}\right.
}
\newcommand{\rf}[1]{[\ref{#1}]} %bibliographical references
\newcommand{\text}[1]{\hbox{\rm \ #1\ \/}}


\newcommand{\be}[1]{\begin{equation}\label{#1}}
\newcommand{\ee}{\end{equation}}
\newcommand{\beo}[1]{\[}  % optional equation - delete numbers unless needed
\newcommand{\eeo}{\]}     % optional equation

%\newcommand{\bl}[1]{\begin{lemma}\label{#1}}
%\newcommand{\el}{\end{lemma}}
%\newcommand{\els}{\end{lemma}}

\newcommand{\bll}[1]{\begin{lemma}\label{#1}}
%\newcommand{\el}{\end{lemma}}
\newcommand{\ells}{\end{lemma}}
\newcommand{\elll}{\end{lemma}}

\newcommand{\bl}[1]{\noindent{\bf Lemma:}}
\newcommand{\el}{}
\newcommand{\els}{}

%\newcommand{\br}[1]{\begin{remark}\label{#1}}
%\newcommand{\er}{\end{remark}}
%\newcommand{\ers}{\end{remark}}

\newcommand{\br}[1]{\noindent{\bf Remark:}}
\newcommand{\er}{}
\newcommand{\ers}{}


\newcommand{\bt}[1]{\begin{theorem}\label{#1}}

%\newcommand{\bd}[1]{\begin{definition}\label{#1}}
%\newcommand{\eds}{\end{definition}}
%\newcommand{\ed}{\end{definition}}
\newcommand{\bd}[1]{{\noindent{\bf Definition:}}}
\newcommand{\eds}{}
\newcommand{\ed}{}

%\newcommand{\bp}[1]{\begin{proposition}\label{#1}}
%\newcommand{\ep}{\end{proposition}}
%\newcommand{\eps}{\end{proposition}}
\newcommand{\bp}[1]{{\noindent{\bf Proposition:}}}
\newcommand{\ep}{}
\newcommand{\eps}{}

%\newcommand{\bpr}{\begin{proof}}
%\newcommand{\epr}{\qquad \end{proof}}
\newcommand{\bpr}{{\noindent{\bf Proof: }}}
\newcommand{\epr}{}
\newcommand{\eprr}{\qed}

\newcommand{\bc}[1]{\begin{corollary}\label{#1}}

\newcommand{\bfact}[1]{\begin{fact}\label{#1}}
\newcommand{\bex}[1]{\begin{example}\label{#1}}
\newcommand{\bem}[1]{\begin{example}\label{#1}}  %Yes, 2 different ones...
\newcommand{\ec}{\end{corollary}}
\newcommand{\efact}{\end{fact}}
\newcommand{\eex}{\end{example}}
\newcommand{\eem}{\end{example}}
\newcommand{\ele}{\end{lemmaex}}
\newcommand{\et}{\end{theorem}}


%%%%%%%%%%%%%%%%%%%%%%%%%%%%%%%%%%%%%%%%%%%%%%%%%%%%%%%%%%%%%%%%%%%%%%%%
%%  These are just for statements, no "qed" box is put in.
%%  (Also for when I have to finish something after a displayed equation
%%        or itemized.)
\newcommand{\ecs}{\end{corollary}}
\newcommand{\eers}{\end{exercise}}
\newcommand{\eexs}{\end{example}}
\newcommand{\eems}{\end{example}}
\newcommand{\eles}{\end{lemmaex}}
\newcommand{\ets}{\end{theorem}}

\newcommand{\qed}{\hfill \halmos} %put !#at right margin

\newcommand{\mybox}{\hfill \mbox{$\Box$}} %put !#at right margin (white square)

\newenvironment{proof}{\noindent {\em Proof}.\ }{\hspace*{\fill}$\halmos$\medskip}
\newcommand{\halmos}{\rule{1ex}{1.4ex}}
\renewcommand{\thefootnote}{\fnsymbol{footnote}}


\newcommand{\optbreak}[1]{{{#1} \nonumber \\ && {#1}}}
\newcommand{\ulevel}{{@g_1}}
\newcommand{\ylevel}{{@g_2}}
\newcommand{\ydissip}{{@g}}
\newcommand{\xdissip}{{@a_3}}
\newcommand{\ygdissip}{{@s_2}}
\newcommand{\ugdissip}{{@s_1}}
\newcommand{\xyrel}{{@c_2}}
\newcommand{\xurel}{{@c_1}}
\newcommand{\om}{{[-1, 1]^m}}
\newcommand{\fhat}{\widehat{f}}
\newcommand{\tslow}[2]{{@s_{#1,#2}}}
\newcommand{\That}{\hat{T}}
\newcommand{\precomp}{{@X}}
\newcommand{\ioslow}{{@q}}
\newcommand{\norma}[2]{{\norm{{#1}|_{#2}}}}
\newcommand{\Z}{{\mathbb Z}}  %ams bold
\newcommand{\sm}{{\vf}}
\newcommand{\qixdissip}{\tilde{@a}_3}
\newcommand{\floor}{{\underline@a}}
\newcommand{\pot}{{\bar@a}} 
\newcommand{\newprop}{{GASMO}}
\newcommand{\klfun}{{@l}}
\newcommand{\newklfun}{{\hat{\klfun}}}
\newcommand{\ix}{{@c}}
\newcommand{\iy}{{@g}}
\newcommand{{\rd}}{{\rm d}}
\newcommand{\rv}{{\rm v}}
\newcommand{\iini}{{@k}}
\newcommand{\hbound}{{\vartheta}}
\newcommand{\mix}{{@f}}
\newcommand{\tf}{{\tilde f}}
\newcommand{\relaxed}{!S}
\newcommand{\rpm}{!P}
\newcommand{\sqb}[1]{\left[ #1 \right]}
\newcommand{\cbrace}[1]{\left\{ #1 \right\}}
\newcommand{\rb}[1]{\left( #1 \right)}
\newcommand{\prosub}{{\partial_P}}
\newcommand{\prosuper}{{\partial^P}}
\newcommand{\abso}[1]{{\abs{#1}}}
\newcommand{\sign}{{\mbox{sign}}}
\newcommand{\mi}{{\setminus}}
\newcommand{\vr}{{\varrho}}
\newcommand{\ra}{{\to}}
\newcommand{\any}{{\forall}}
\newcommand{\deq}{{:=}}
\newcommand{\sub}{{\subseteq}}
\newcommand{\inter}{{\mbox{int}}}
\newcommand{\si}[1]{{\left\lfloor #1 \right\rfloor}}
\newcommand{\ds}{\displaystyle}
\newcommand{\dist}{\mbox{dist}}
\newcommand{\pcow}{!W}


\newcommand{\dfun}{{\bf d}}
\newcommand{\ddfun}{{\bf w}}  % {{\bf d}^d}
\newcommand{\dd} {w}          % {{d^d}}
\newcommand{\du}  {{d_u}}
\newcommand{\dufun}  {{\bf d}_u}
\newcommand{\uc}{{\bf u}}
\newcommand{\vc}{{\bf v}}
\newcommand{\yc}{{\bf y}}

\newcommand{\omd}{{[-1, 1]^{m_w}}} %{{[-1, 1]^{m_d}}}
\newcommand{\omu}{\controls_1}
\newcommand{\Omegad}{{@G}}   %{{@W^d}}
\newcommand{\Omegau}{{[-1, 1]^{m_u}}}  %   {{@W^u}}
\newcommand{\xlevel}{{@a_x}}
\newcommand{\iu}{{@g_1}}
\newcommand{\UintIOSS}{U(iI)OSS}
\newcommand{\nobserver}{@S_{n.o}}
\newcommand{\uobs}{\hat @g_1}
\newcommand{\yobs}{\hat @g_2}
\newcommand{\outputs}{!Y}

\newcommand{\intEO}{{\E_1}}
\newcommand{\precompone}{@X_1}

\newcommand{\Lyapunovlike}{Lyapunov-like}

\newcommand{\myboxineq}{\eqno\Box} %put !#at margin (white square) in eqn



\newcommand{\Arr}{$\leadsto\;$}
\newcommand{\xo}{x^{\scriptscriptstyle 0}}
\newcommand{\xtnou}{x(t,\xo)}
\newcommand{\tzz}{\tlim 0}
\newcommand{\tlim}{\raisebox{-1.5ex}{$\stackrel{\longrightarrow}%
{{\scriptstyle t->!8}}$}}
\newcommand{\afunction}{}
\newcommand{\function}{}
\newcommand{\functions}{}

\newcommand{\Therefore}{\therefore}  %in case want to put words back

\newcommand{\includepicx}[2]{\includegraphics[scale=#1]{FIGS-04NONLINEAR/#2}}

\begin{document}

{\Large

``Input-Output-to-State Stability''\\

Mikhail Krichman, Eduardo D. Sontag, Yuan Wang

SIAM J. Control and Optimization {\bf 39}(2001): 1874--1928\\

Lyapunov characterizations of input-output-to-state stability (IOSS)\\


(a type of ``well-posed detectability'')\\

equivalence IOSS \& existence of smooth IOSS-Lyapunov function\\

Corollaries:\\

\bi
\item
existence of ``norm-estimators''
\item
characterize IOSS in terms of ``relative stability''
\item
or via finite-energy estimates
\item
characterization of ISS (special case)
\ei

\spp

proof uses many interesting techniques \& concepts


\newpage

\section{Introduction} 

{\em observers: estimate, on the basis of external information 

provided by past input and output signals, 

the internal state $x(t)$ at time $t$}

\begin{center}
\setlength{\unitlength}{0.012500in}%
\begin{picture}(310,65)(340,660)
\thinlines
\put(400,660){\framebox(100,40){}}
\put(340,680){\vector(1,0){60}}
\put(500,680){\vector(1,0){60}}
\put(370,680){\line(0,1){35}}
\put(370,715){\vector(1,0){190}}
\put(560,670){\framebox(50,55){}}
\put(610,695){\vector(1,0){30}}
\put(350,690){$u$}
\put(515,690){$y$}
\put(445,678){$x$}
\put(650,695){$\hat x$}
\end{picture}
\end{center}

``algorithm'' (dynamical systems) that converges to an estimate

$\hat x(t)$ of the state $x(t)$ of the system of interest, 

using information $\{u(s), s!<t\}$ (set of past input values)

and $\{y(s), s!<t\}$ (set of past output measurements)\\


State estimation central to control theory: 

signal processing applications (Kalman filters); 

stabilization based on partial information (observers)

well understood for linear systems, still poorly developed in general\\


when interested in stabilization to an equilibrium, 

e.g. zero state $x=0$ in Euclidean space, 

weaker type of estimate sometimes enough; {\em norm-estimators:}

require only bound $\hat x(t)$ on {\em norm}
$\abs{x(t)}$ of the state $x(t)$, cf.:

{\sc L. Praly and Y. Wang}, {\em Stabilization in spite of matched unmodelled
  dynamics and an equivalent definition of input-to-state stability},
  MCSS 9 (1996) pp. 1--33

\spp

{\sc Z.-P. Jiang and L. Praly}, {\em Preliminary results about robust
  Lagrange stability in adaptive nonlinear regulation}, IJC 6 (1992), pp. 285--307

\newpage

to understand issues, start by considering special case:

when the external data ($u$ and $y$) vanish identically,

obvious estimate (assuming $x=0$ equilibrium for unforced system 

and output is zero when $x=0$) is $\hat x(t)==0$

if this estimate bounds norm of true state as $t->!8$, $\therefore$ $x(t)->0$

i.e. necessary property for the possibility of norm-estimation:

origin must be a globally asymptotically stable state wrt

``subsystem'' consisting of those states for which

$u==0$ produces output $y==0$: system is {\em zero-detectable}

\spp

[ for {\em linear systems}, zero-detectability $==$ detectability

(two trajectories produce $=$ output $=>$ they approach each other)

but not in general (``incremental'' notion $!=$)]

\spp

but: zero-detectability not ``well-posed'' enough:

ask also: \ inputs and outputs small $=>$ states small, and

inputs and outputs converge to zero as $t->!8$ $=>$ states do too

\begin{center}
%\setlength{\unitlength}{0.012500in}%
\setlength{\unitlength}{0.02in}%
\begin{picture}(220,45)(340,660)
\thinlines
\put(400,660){\framebox(100,40){}}
\put(340,680){\vector( 1, 0){ 60}}
\put(500,680){\vector( 1, 0){ 60}}
\put(350,690){$u->0$}
\put(420,678){$=>\;\; x->0$}
\put(515,690){$y->0$}
\end{picture}
\end{center}

also natural: asymptotic bounds on states as function of i/o bounds

and quantify ``overshoot'' (transient behavior)

\spp

system $@.x=f(x,u)$ with measurement (output) map $y=h(x)$

is {\em input-output-to-state stable\/} (\ioss)

if $\exists$
$@b!e\kl$ and $\ulevel, \ylevel !e\ki$ s.t. 
the estimate
\[
 \abs{x(t)} !<\max \left\{ @b(\abs{x(0)}, t), 
\ulevel\left(\norma{u}{[0, t]}\right), 
\ylevel\left(\norma{y}{[0, t]}\right)\right\}
\]
holds for any initial state $x(0)$ and any input $u(!.)$,

where $x(t) = x(t,x(0),u(!.))$ and $y(t)=h(x(t))$

(precise defs later)

\newpage

terminology \ioss:

``stability from the input/output data to the state''\\

term introduced in 

{\sc E. D. Sontag and Y. Wang}, {\em Output-to-state
 stability and detectability of nonlinear systems}, SCL 29 (1997), pp. 279--290

\spp

combines ``strong'' observability (TAC 89) and ISS;

called simply ``detectability'' in

{\sc E. D. Sontag}, {\em Some connections between stabilization and
  factorization}, in CDC89, pp. 990--995

(stated there in i/o terms, applied to controller parameterization)

\spp

and called ``strong unboundedness observability'' in

{\sc Z.-P. Jiang, A. Teel, and L. Praly}, {\em Small-gain theorem for ISS
  systems and applications}, MCSS 7(1994), pp. 95--120

(more precisely, allowing additive nonnegative constant in rhs)

and related to other ISS-like formalism for observers, e.g.:

\spp

{\sc D. J. Pan, Z. Z. Han, and Z. J. Zhang}, {\em Bounded-input-bounded-output
  stabilization of nonlinear systems using state detectors}, SCL 21 (1993), pp. 189--198

{\sc W. M. Lu}, {\em A class of globally stabilizing controllers for nonlinear
  systems}, SCL 25 (1995), pp. 13--19

\newpage


will show \ioss\ iff $\exists$ norm-estimator, using nec \& suf. 

characterization of \ioss\ by smooth dissipation functions: 

\spp

$\exists$ proper (radially unbounded) positive definite

smooth function $V$ of states (``storage function'' [Willems])

s.t. {\em dissipation inequality\/}:
\beo{diss-intro}
\frac{d}{dt}V(x(t))\;!<\;
-@s_1(\abs{x(t)}) +  @s_2(\abs{y(t)}) + @s_3(\abs{u(t)})
\eeo

\spp

holds along all trajectories, with the functions $@s_i$ $!e\ki$

\spp

this provides an ``infinitesimal'' description of \ioss

and a norm-observer is easily built from $V$

\spp

earlier several independent suggestions that

one should {\em define} ``detectability'' in dissipation terms, e.g. in:

{\sc W. M. Lu}, {\em A state-space
  approach to parameterization of stabilizing controllers for nonlinear
  systems}, IEEE TAC 40 (1995), pp. 1576--1588

\spp

detectability defined by the requirement that $\exists$
differentiable storage function $V$ as above, but with
special choice $@s_2(r):=r^2$ and no inputs

or weaken the dissipation inequality, to merely require:
\[
 x!=0 \;=>\; \frac{d}{dt}V(x(t))< @s_2(\abs{y(t)})
\]
(again, with no inputs), as in:

{\sc A. S. Morse}, {\em Control using logic-based switching}, in Trends in
  Control: A European Perspective, A. Isidori, ed., Springer-Verlag, London, 1995,
  pp. 69--114

\spp

(no ``margin'' of stability $-@s_1(\abs{x(t)})$)

\newpage

note: asking that along all trajectories there holds
\[
\frac{d}{dt}V(x(t))\;!<\;
-@s_1(\abs{x(t)}) +  @s_2(\abs{y(t)}) + @s_3(\abs{u(t)})
\]
means $V$ satisfies partial differential inequality (PDI):
\[
\max_{u\in\R^m} \left\{\nabla V(x) !.f(x,u) + @s_1(\abs{x}) 
-  @s_2(\abs{h(x)}) -  @s_3(\abs{u}) \right\} \;\leq\; 0
\]
which is same as Hamilton--Jacobi inequality
\[
g_0(x) + \frac{1}{4}\sum_{i=1}^m (\nabla V(x)!.g_i(x))^2 
  + @s_1(\abs{x}) - @s_2(\abs{h(x)}) \;\leq\; 0 \,.  
\]
in special case of quadratic input ``cost'' $@s_3(r)=r^2$ and\\
systems $@.x = f(x,u)$ affine in controls
\[
@.x \;=\; g_0(x) \,+\, \sum_{i=1}^m \,u_i \,g_i(x) 
\] 
(just replace RHS in (PDI) by max value,
obtained at $u_i = (1/2)\nabla V(x)!.g_i(x)$)

so converse result amounts to providing nec \& sufft conds

for $\exists$ smooth (and proper and positive definite) solution $V$ to PDI

\sppp

remark: $\exists$  lower semicontinuous $V$ $=>$ $\exists$ $!C^!8$
(possibly $!=$ $@s_i$'s)

(``solution'' interpreted in a weak sense, e.g. viscosity):

Mikhail Krichman, Eduardo D. Sontag

{\em Characterizations of detectability notions in terms of discontinuous
dissipation functions},
IJC 75 (2002): 882--900

\newpage


key preliminary step in the construction of $V$

(just as it was for the analogous result for ISS)

is characterization of \ioss\ property in robustness terms

by means of a ``small gain'' argument:

\spp

\ioss\ equivalent to the existence of ``robustness margin'' $@r!e\ki$

i.e. every system obtained under a feedback law $@D$

(even dynamic and/or time-varying) 

for which $\abs{@D(t)}!<@r(\abs{x(t)})$ for all $t$

must be \oss\ (i.e., is \ioss\ as a system with no inputs)

\begin{center}
\setlength{\unitlength}{0.012500in}%
\begin{picture}(135,90)(190,570)
\thinlines
\put(220,620){\framebox(80,40){}}
\put(240,570){\framebox(40,30){}}
\put(260,620){\vector( 0,-1){ 20}}
\put(240,585){\line(-1, 0){ 40}}
\put(200,585){\line( 0, 1){ 55}}
\put(200,640){\vector( 1, 0){ 20}}
\put(300,640){\vector( 1, 0){ 20}}
\put(257,635){$x$}
\put(190,645){$u$}
\put(325,645){$y$}
\put(257,580){$@D$}
\end{picture}
\end{center}

\Arr ``systems with disturbances'' (diff. inclusions)

\spp

since such systems must be introduced anyhow,

we present all our results (and definitions)

for systems with disturbances

\spp

core of proof is construction of $V$ for ``robustly detectable'' 

(more precisely, ``robust \oss'') 

for $@.x=g(x,d)$ obtained when $u=d@r(\abs{x})$ in original

with $d=d(!.)$ $=$ measurable \function w/values in unit ball

function $V$ must satisfy PDI:
\[
\nabla V (x) !. g(x,d) \;!<\; -@s_1(\abs{x})+@s
_2(\abs{y})
\]
for some \functions $@s_1$ and $@s_2$ $!e\ki$

\newpage

further reduction: ``relatively asymptotically stable'' systems:

if $@.V(x(t))!< -@s _1(\abs{x(t)})+@s_2(\abs{y(t)})$, then

\spp

$V$ {\em must decrease along trajectories as long as $y(t)\ll x(t)$}

\spp

\Arr ``global asymptotic stability modulo outputs''

\spp

construction of $V$ relies upon solution of an appropriate 

optimal control problem, for which $V$ is the value function\\

and first ``fuzzifying'' the dynamics near the set where $y\ll x$,

so as to obtain a problem whose value function is continuous

\spp

followed by use of techniques from nonsmooth analysis 

to obtain Lipschitz, and $\therefore$ by regularization argument $!C^!8$ $V$

starting from the $!C^0$ $V$ obtained from the OC problem

\newpage

``energy'' estimates for detectability also of interest; e.g.:

\spp

{\sc J. W. Helton and M. R. James}, {\em Extending $H_\infty$ Control to Nonlinear
  Systems}, SIAM, Philadelphia, 1999.

\spp

defined ``$L^2$-detectability'' by a requirement that 

state trajectory in $L^2$ if observations are

\sppp

we'll show: integral variant of \oss\ $\iff$ \oss\ \& UO

where UO = ``unboundedness observability'' (UO)

\spp

($\sim$ ``forward completeness modulo outputs'', due to:

{\sc F. Mazenc, L. Praly, and W. Dayawansa}, {\em Global stabilization by
  output feedback: Examples and counterexamples}, Systems Control
  Lett., 23 (1994), pp. 119--125)\\

\newpage

\section{Definitions \& Statements of Main Results}

\subsection{systems}

study systems w/2 types of inputs: 
{\em controls} and {\em disturbances}
\beo{udsystem}
 @.{x}(t)=f(x(t), \uc(t),\ddfun(t)), \;\;\; y(t)=h(x(t)).
\eeo
%

\spp

states evolve in $\states = \R^n$

controls: measurable,
ess bounded $\uc: \I =\R_{!>0} ->\controls := \R^{m_u}$

disturbances: measurable \functions $\ddfun: \I \to \Omegad$
%11/26

where $\Omegad$ $=$ compact, convex subset of $\R^{m_w}$

denote set of such functions by $\M_\Omegad$

\spp

if different interval 
$\I \subset \R_{!>0}$ of definition for control $\uc$,

apply definitions to extension of $\uc$ to $\R_{!>0}$ \ ($\uc ==0$ on $\R_{!>0} \setminus \I$)

\sppp

$f: \states !p\controls !p\Omegad \to \states$ locally Lipschitz in $(x, u)$ uniformly on $\dd$

jointly continuous in $x$, $u$, and $\dd$

and  s.t. $f(0, 0, \dd) = 0$ for any $\dd !e\Omegad$

$h: \states \to \outputs := \R^p$ is smooth ($C^1$) 
and vanishes at 0

\sppp

given state $@x!e\states$ and pair  $(\uc, \ddfun)$

denote $x(t, @x, \uc, \ddfun)=$ unique maximal solution

defined on maximal interval $[0, \tmax(@x, \uc, \ddfun))$

denote $y(t,@x, \uc, \ddfun):=h(x(t, @x, \uc, \ddfun))$

\spp

if clear, just ``$\tmax$'' \ ``$x(t)$'' \ ``$y(t)$''

\newpage

\subsection{comparison functions}

recall {\em global asymptotic stability (GAS)} of origin means:
\[
\abs{\xtnou} \;!<\; @b\left(\abs{\xo},t\right)
\]
($\forall\,\xo\,,\;\forall\,t!>0$)

for some  $@b(\nearrow,\searrow)!e\kl$ \ \ ($@b(0,!.)=0$, $!C^0$)

\beqn
\abs{\xtnou} &!<& @b\left(\abs{\xo},0\right)
\ \text{\Arr stab (small overshoot)}\\
\abs{\xtnou} &!<&  @b\left(\abs{\xo},t\right) \tzz
\ \text{\Arr \ attractivity}
\eeqn

\sppp

$@a:\R_{!>0}->\R_{!>0}$ {\em class $\kk$} if $!C^0$, pos def, strictly increasing

and {\em class $\ki$} if also unbounded

\begin{center}
\includegraphics[scale=0.01]{FIGURES/ki-fig}
\end{center}
%
\function $@b: \R_{!>0} !p\R_{!>0} \to \R_{!>0}$ {\em class $\kl$} if:

 $\forall$  $t !>0$, $@b(!., t)$ $!e\kk$,

 and $\forall$ $s !>0$, $@b(s,t)\searrow0$ as $t \to \infty$

\sppp


%remark: \ 
%$\forall\,@b!e\kl$, $\exists$ $@a_1,@a_2!e\ki$ s.t.\\
%$@b(r,t) \,!< \, @a_2\left(@a_1(r)e^{-t}\right)$ $\forall$
%$s!>0, t!>0$
%
%so could also write
%$\abs{\xtnou} !< @a_2\left(@a_1(\absn{\xo})e^{-t}\right)$

\newpage

\subsection{definition of uniform IOSS}

\sppp

$\forall$ $z(!.)$ measurable,
$\norma{z}{[t_1, t_2]}$ $:=$
$L_\infty$ norm of restriction
to $[t_1, t_2]$

\bd{def:UIOSS}
{\em uniformly input-output-to-state stable (UIOSS)}

system if $\exists$ $@b!e\kl$ and $\ulevel, \ylevel !e\kk$ s.t. estimate holds:
\beo{eq:UIOSS}
 \abs{x(t, @x, \uc, \ddfun)} !<\max \left\{ @b(\abs{@x}, t), 
\ulevel\left(\norma{\uc}{[0, t]}\right), 
\ylevel\left(\norma{y}{[0, t]}\right)\right\}
\eeo
$\forall$ initial state 
$@x!e\states$, control $\uc$, disturbance $\ddfun$\\
$\forall$ $t!e[0,\tmax(@x,\uc,\ddfun))$  
\eds

\spp

\bd{def:UIOSS-lyap}
$C^{\infty}$ $V: \states \to \R_{!>0}$ {\it UIOSS-Lyapunov function} if:
%add ''for above system'' if in paper
\bi
\item
$\exists$ $\ki$-functions $@a_1$, $@a_2$ s.t.:
\beo{V_proper}
@a_1(\abs{@x}) !<V(@x) !<@a_2 (\abs{@x})
\;\;\mbox{\small$\forall\,@x!e\states$}
\eeo
\item
$\exists$ $@a!e\ki$ and $\ugdissip,\ygdissip!e\kk$ s.t.:
\beo{pqUIOSS-dissip} 
\!\!\!\!\!\!\!\!
\nabla V(@x) !.f(@x, u, \dd)
 !< -@a(\abs @x)
        + \ugdissip(\abs{u}) + \ygdissip(\abs{h(@x)}) 
\;\mbox{\small$\forall\,@x!e\states,\,u !e \controls,\,\dd !e\Omegad$}
\eeo
\ei
\eds

\spp

note: first condition is just positive definiteness and properness

($\exists$ $@a_2$ is just $V(0) = 0$ (use $V$ $!C^0$), but convenient notation)

second is a {\em dissipation inequality} in Willems' sense

\spp

also note: IOSS, ISS, ISS-Lyap are particular cases

\newpage

\br{non-dissip}
if $V: \states \to \R_{!>0}$ 
smooth s.t. $@a_1(\abs{@x}) !<V(@x) !<@a_2 (\abs{@x})$,

$V$ is UIOSS-Lf iff $\exists$ $\xdissip$ $!e\ki$, and 
 $\ydissip$ and $\xurel$ $!e\kk$ s.t.
\be{eq:non-dissip}
\abs @x!>\xurel(\abs u) \;=>\;
\nabla V(@x) !.f(@x, u, \dd) !< -\xdissip(\abs @x) + 
 \ydissip(\abs{h(@x)})
%\;\mbox{\small$\forall\,@x!e\states,\,u !e \controls,\,\dd !e\Omegad$}
\ee
%for any  $@x!e\states$, $\dd !e\Omegad$,  and 
%$u !e\controls$ s.t.
%$\abs @x!>\xurel(\abs u)$.

\noindent{\bf Proof:}
if
\be{pqUIOSS-dissip} 
\nabla V(@x) !.f(@x, u, \dd)
 !< -@a(\abs @x)
        + \ugdissip(\abs{u}) + \ygdissip(\abs{h(@x)}) 
\ee
then~\rref{eq:non-dissip} holds with
 $\xdissip(!.):=@a(!.)/2$,
 $\ydissip ==\ygdissip$, and 
$\xurel:= @a^{-1} \circ(2\ugdissip)$

\spp

converse:
% To prove the other implication, assume now that~\rref{eq:non-dissip} holds
% with some $\xdissip !e\ki$ and $\ydissip, \,\xurel !e\kk$. 
define $\ugdissip(!.) = \max\cbrace{0, \hat{\ugdissip}(!.)}$, where
$\hat{\ugdissip}(r)$
is
\[
\max\cbrace{\nabla V(@x)!.f(@x, u, \dd) +  
             \xdissip(\xurel(\abs u))}
\]
over all
\[
             \abs u !<r\,, \;\; \abs@x!<\xurel (r)\,,\;\; \dd !e\Omegad
\]
$\ugdissip!e!C^0$, $\ugdissip(0) = 0$, and w.l.o.g. $\ugdissip!e\ki$ (majorize if not)
\spp

claim:~\rref{pqUIOSS-dissip} holds with $@a==\xdissip$ and $\ygdissip ==\ydissip$:

if $\abs @x!>\xurel (\abs {u})$, then~\rref{eq:non-dissip} holds, from which 
\rref{pqUIOSS-dissip} trivially follows

\spp
if $\abs @x< \xurel (\abs { u})$, then, by definition of $\ugdissip$,
%
\[
\ugdissip(\abs { u}) !> \nabla V(@x) !.f(@x,  u, \dd) 
 + \xdissip(\xurel(\abs { u}))
\]
%
for every $ \dd$, which, in turn, implies~\rref{pqUIOSS-dissip}
\er  

\newpage

%A few particular cases of the UIOSS property have been studied in the 
%literature.  
%If the system~\rref{udsystem} in consideration has no
%outputs and no disturbances, UIOSS reduces to the well-known ISS
%property, whose Lyapunov characterization was obtained in
%~\cite{ywscl}. In case~\rref{udsystem} is autonomous, UIOSS becomes
%OSS. This property was introduced in~\cite{ref8}
% where Lyapunov-type necessary and
%sufficient conditions were obtained. Finally, for systems with no 
%disturbances, UIOSS is just IOSS. This property was introduced 
%in~\cite{oss-scl}, where it was conjectured that any IOSS control 
%system admits a  smooth IOSS-Lyapunov 
%function. This conjecture will be proven here 
%in a more general setting,  for systems forced by both controls
%and disturbances. A few interesting applications of this Lyapunov 
%characterization were also discussed in~\cite{oss-scl}, one of them 
%to be defined next.

\subsection{norm-estimators}

\bd{norm-observer:def}
{\em state-norm-estimator} (or {\em state-norm-observer}):
\beo{norm-observer:eq}
\left(\nobserver\right)\quad
@.{p}=g(p,u, y)\,,
%\eqno(\nobserver)
\quad\quad k : \R^\ell !p\outputs \to \R_{!>0}
\eeo
%
evolving in $\R^\ell$, driven by the controls \& outputs of $@S$, s.t.:
\bi
\item
$\exists$ $\uobs,\yobs!e\kk$, $\hat @b!e\kl$ s.t. $\forall$
initial state $@z!e\R^\ell$, inputs $\uc$ and ${\bf y}$, $t$ in interval of
definition of solution $p(!., @z, \uc, {\bf y})$:
\beo{nobs-ios}
k\Bigg(p(t, @z, \uc, \yc), \yc(t)\Bigg) !<\hat @b(\abs @z, t) +
\uobs\rb{\norma{\uc}{[0, t]}} + \yobs \rb{\norma{\yc}{[0, t]}}
\eeo

\item $\exists$ $@r!e\kk$, $@b!e\kl$ s.t.

$\forall$ initial states $@x$ and $@z$ of system \& observer,
control
$\uc : \R_{!>0} \to \controls$ and 
disturbance $\ddfun !e\M_\Omegad$:
%
\beo{norm-estimate}
\abs{x(t, @x, \uc, \ddfun)} !<@b(\abs@x+ \abs@z, t) +
  @r\left({k\bigg(p(t, @z, \uc, \yc_{@x,\uc, \ddfun}),
                                  \yc_{@x,\uc, \ddfun}(t) \bigg)}\right)
\eeo
%
$\forall$ $t !e[0, \tmax(@x, \uc, \ddfun))$

($\yc_{@x,\uc, \ddfun}$ is output trajectory of $@S$, i.e.
$\yc_{@x, \uc, \ddfun}(t)
= y(t, @x, \uc, \ddfun)$)
\end{itemize}
\eds

\spp

observer output is IOS [stable] wrt the inputs $u$ and $y$

and true state asymptotically bounded in norm 

by function of norm of estimator

with transient (overshoot) depending on initial states of both

\newpage


\subsection{statement of the main theorem}
      
\bt{main_thm:UIOSS}
%Let $@S$ be a system of type~\rref{udsystem}. 
\ The following are equivalent:
\begin{enumerate}
\item[$1$.] \label{item1} 
 $@S$ is UIOSS. 
\item[$2$.] \label{item2}
 $@S$ admits a UIOSS-Lyapunov function. 
\item[$3$.] \label{item3}
 $\exists$ state-norm-estimator for $@S$.
\end{enumerate}
\ets

The main contribution is in showing that 1 implies 2; 

the remaining implications are much easier.

\newpage

\subsection{simple example: linear systems}
\renewcommand{\A}{{\bf A}}
\newcommand{\C}{{\bf C}}
\renewcommand{\B}{{\bf B}}
\renewcommand{\L}{{\bf L}}
\renewcommand{\P}{{\bf P}}

%A particular class of systems~\rref{udsystem} is as follows.
%A {\em linear}, time-invariant system $@S_{lin}$ with outputs is one for
%which $f$ and $h$ are linear, that is,
\beqn
@.x &=& \A x + \B u,\label{lsystem}\\ 
     y &=& \C x \nonumber
\eeqn
where $\A !e\R^{n !pn}$, $\B !e\R^{n !pm}$, 
and $\C !e\R^{p !pn}$ \ \ (assume $m_w=0$)

\spp

Recall:
{\em detectable}, or {\em asymptotically observable}, if
implication 
\beo{e-detect}
\C x(t) ==0  \;\; \Longrightarrow \;\;x(t) \to 0
\eeo   
holds for any  trajectory $x(t)$, corresponding to zero control $\uc ==0$
\ed
%
%The following result is a totally routine linear systems theory fact, but we include
%the proof as a motivation for the nonlinear material to follow.

\spp

\bp{detectimpliesIOSS}
for linear systems, detectable $=>$ IOSS
\eps

\spp

\bpr
 pick $\L!e\R^{n !pp}$ s.t. $\A + \L\C$ is Hurwitz, so observer:
\beo{observer}
@.z(t) = \A z(t) + \B \uc(t) + \L (\C z(t) -y(t)),
\eeo
is s.t. $\abs{x(t) - z(t)} \to 0$ (in particular
$x(0) = z(0)$ $=>$ $x(t) == z(t)$)

fix an initial state $@x$ and a control $\uc$

soln $x(t, @x, \uc)$ of system is also solution of observer with $z(0)=@x$,

so that 
\[
x(t, @x, \uc)= e^{t(\A+\L\C)}@x+ 
\int_0^t e^{s(\A+\L\C)}[\B \uc(t-s) - \L y(t-s)]ds
\]

choose two positive numbers $@d'$ and $@d$ 
so that $\Re @l!< - @d' < -@d$ 

for every 
eigenvalue $@l$ of $\A + \L \C$

then $\exists$ polynomial 
$P(!.)$ and $\therefore$ a constant $K$ s.t. $\abs{x(t, @x, \uc)} !<$
\[
P(t) e^{- @d' t} \abs @x+
  \int_0^t P(s)e^{-@d' s}
\sqb{\norm \B \abs{\uc(t-s)} + \norm \L \abs{y(t-s)}}ds
\]
\[
!< \; K e^{-@dt} \abs @x+ 
K \frac{\norm \B}{@d} \norma{\uc}{[0, t]} 
+ K \frac{\norm \L}{@d} \norma{y}{[0, t]}\label{obs-x}.
\]
%
thus IOSS estimate holds with the linear gains 

$@b(r,t) := K e^{-@dt} r$ and $\ulevel(r)=\ylevel(r):= K\frac{\norm\B}{@d} r$
\epr

\newpage

{\bf Remark:} observer is in particular a norm estimator:

solution $z(!., @z, \uc, \yc)$ of observer satisfies estimate 
\[
\abs{z(t,@z, \uc,\yc)}\le  K e^{-@dt} \abs @x+ 
K \frac{\norm \B}{@d} \norma{\uc}{[0, t]} 
+ K \frac{\norm \L}{@d} \norma{y}{[0, t]}
\]
where $K, @d$ are same as above; and with $e:= x-z$, we have:
\[
@.e(t) = (A+LC)e(t).
\]
so along any trajectory $(x(t), z(t))$:
\[
\abs{x(t)} \le K \abs{@x- @z}e^{-@dt} + \abs{z(t)}
\]
$\forall$ $t\ge 0$, where $K, @d$ are again as before

\spp

{\bf Remark:} IOSS-Lyapunov function is $V(x):= x'\P x$

where $\P !e\R^{n !pn}$ symmetric matrix s.t.:
\[
\P(\A+\L\C) + (\A+\L\C)' \P = -I
\]
since $\P(\A + \L\C) = ((\A + \L\C)'\P)'$,

$x'\P(\A +\L\C)x = x'(\A + \L\C)'\P x$; $\therefore$ $\nabla V(x) !.f(x, u) =$
\begin{eqnarray*}
%\nabla V(x) !.f(x, u) 
&=& 2 x' \P\, (\A x +\B u) \\
                          &=& 2 x' \P\, ((\A + \L\C)x + \B u - \L y)\\
                          &=& 2 x' \P (\A + \L\C)x + 2 x' \P \B u - 2x'\P\L y\\
                          &=& x' (\P (\A + \L\C) +(\A + \L\C)'\P) x  +
                                        2 x' \P \B u - 2x'\P\L y \\
                          &\leq& -\abs x^2 + 2 \norm \P \norm \B \abs u \abs x 
                                        + 2 \norm \P \norm \L \abs y \abs x\\
                          &\leq& -\abs x^2 + {\abs x}^2 /4 + 
                                        4\norm \P^2 \norm \B^2 {\abs u}^2 +
                           {\abs x}^2 /4 + 4\norm \P^2 \norm \L^2 {\abs y}^2\\
                          &\leq& -\abs x^2/2 
                                + 4\norm \P^2 \norm \B^2 {\abs u}^2
                                + 4\norm \P^2 \norm \L^2 {\abs y}^2.
\end{eqnarray*}
$\therefore$ UIOSS dissipation inequality holds for $V$ with gains defined by

$@a(r) = r /2$, $\ugdissip(r) = 4\norm \P^2 \norm \B^2 r^2$, $\ygdissip(r) = 4\norm \P^2 \norm \L^2 r^2$.

\newpage


%%%%%%%%%%%%%%%%%%%%%%%%%%%%%%%%%%%%%%%%%%%%%%%%%%%%%%%%%%%%%%%%%%%%%%%%%

\renewcommand{\A}{!A}
\renewcommand{\C}{!C}
\renewcommand{\B}{!B}
\renewcommand{\L}{!L}
\renewcommand{\P}{!P}

%%%%%%%%%%%%%%%%%%%%%%%%%%%%%%%%%%%%%%%%%%%%%%%%%%%%%%%%%%%%%%%%%%%%%%%%%

\subsection{systems without controls}

$@W$  compact metric space ($=\om$ unless specified otherwise)

consider systems of the type
%
\be{dsystem}
 @.{x}=f(x(t), \dfun(t)), \;\;\; y(t)=h(x(t))
\ee

where $f: \states !p@W$ is locally Lipschitz in $x$ uniformly on $d$ 

jointly continuous in $(x,d)$, $f(0, d) =0$ for any $d !e@W$

inputs: measurable \functions $\dfun: \I \to @W$, 

``{\em disturbances}'': $@W$-valued inputs

$\M_@W=$ collection of all such functions

is particular case of $@.{x}(t)=f(x(t),\uc(t),\ddfun(t))$

 (ignoring controls and driven only by disturbances)

for convenience, specialize definitions:
%

\bd{UOSS}
system UOSS (uniformly output-to-state stable) if $\exists$ $@b!e\kl$ and 
$\ylevel !e\kk$ s.t.
\beo{eq:UOSS}
 \abs{x(t, @x, \dfun)} !<\max \left\{ @b(\abs{@x}, t), 
\ylevel
\left(\norma{y}{[0, t]}\right)\right\}
\eeo
for any disturbance $\dfun$, initial state $@x\in\states$, and
$t\in\maxint$. 
\eds

\bd{def:UOSS-Lyap-smooth}
{\it UOSS-Lyapunov function} for~\rref{dsystem} is smooth 

\function $V: \states\to \R_{!>0}$ satisfying
$@a_1(\abs{@x}) !<V(@x) !<@a_2 (\abs{@x})$ and
\beo{UOSS-dissip-smooth}
\nabla V (x) !.f(x, d) !< -\xdissip(\abs x) 
+ \ydissip(\abs{h(x)}) \quad \forall\,x !e\states, \; \forall\,d 
!e@W,
\eeo 
with \functions $@a_i!e\ki$, \function $\ydissip!e\kk$

for systems with no disturbances, just ``$V$ is OSS-Lyap
function''
\ed

\newpage
\subsection{``Modulo outputs'' relative stability} \label{sec:GASMO}

classical uniform global asymptotic stability (GAS) for $@.{x}=f(x(t), \dfun(t))$

asks all sols $->$ equilibrium, short excursions

\spp

suppose does not matter how system behaves when information provided by output
is adequate (norm of output dominates norm of current state)

but want nice decay when the output does not help 
in determining how large the state is

 \bd{newprop:def}
 \newprop\  (global asymptotic stab modulo output):

if
 $\exists$ \function $@r$ $!e\ki$ and  {\afunction} 
$\klfun$ $!e\kl$ s.t., $\forall$ 
  $@x!e\states$, $\dfun !e\M_@W$, and any $T < \tmax(@x, \dfun)$,
if 
\[\abs{x(t, @x, \dfun)} !>@r(\abs{h(x(t, @x, \dfun))}) 
\quad \forall \; 0 !<t !<T,
\]
then the estimate
\beo{GASMO-est}
\abs{x(t, @x, \dfun)} !<\klfun (\abs{@x}, t) \quad \forall\; 0 !<t
!<T  
\eeo
holds
\ed

\br{} 
if a system has no outputs, then \newprop\ $==$ GAS
\er

\newpage
\subsection{``$\ve$-$@d$'' characterization of the \newprop\ property}

\bp{prop:newprop}
a system satisfies \newprop\ 
iff $\exists$ $@r!e\ki$ so that the 
following  two properties hold.
\begin{enumerate}
\item[$1$.]
 $\forall$ $\ve > 0$ and any $r > 0$, $\exists$ $T_{r, \ve}$ s.t. $\forall$ $\abs{@x}\le r$, any $\dfun$, 
  and any $T !e[0, \tmax(@x, \dfun))$ s.t.  $T \ge T_{r, \ve}$, if
\[
 \abs{x(t, @x, \dfun )} \;!>\;@r(\abs{y(t, \; @x, \dfun)})\;
\; \forall \ 0!< t!<T \,,
\]
then
\[
\abs{x(t,\; @x, \dfun )} \;< \;\ve\;
\; \forall \; t\in[T_{r,\ve}, T]\,.
\]
\item[$2$.] 
$\exists$  $\hbound!e\kk$ s.t. $\forall$ $@x\in
\states$, any disturbance $\dfun$, and any $T < \tmax(@x, \dfun)$ s.t.
\[
\abs{x(t,@x, \dfun )} \;!>\;@r(\abs{y(t,@x, \dfun)})\;
\; \forall \;0 !<t !<T\,,
\]
the following ``bounded overshoot'' estimate holds:
\[
\abs{x(T,@x, \dfun)} \;\leq\;\hbound(\abs@x)\,.
\]
\end{enumerate}
\eps

necessity obvious; for sufficiency use (from Lin/Wang/EDS SIAM96)

\bl{from_LSW_2}
if $@F(r, t) \,: \,(\R_{!>0})^2 \to \R_{!>0}$ s.t.
\begin{enumerate}
\item[$1$.] $\forall$ $\ve > 0$ and $\forall$ $R > 0$ $\exists$ $T$ s.t. 
$@F(r, t) < \ve$ $\forall$ $0 !<r !<R$ and $\forall$ $t !>T$,
\item[$2$.]  $\forall$ $\ve > 0$ $\exists$ $@d> 0$ s.t. if 
$r !<@d$, then $@F(r, t) < \ve$ $\forall$ $t > 0$. 
\end{enumerate}
then $@F$ can be majorized by a $\kl$-function
\el

now consider the function $@F(r, t):= \sup \{\abs{x(t, @x, \dfun)}\}$
over all
\[
\abs @x!<r,  \;
   \dfun !e\M_@W, \;\abs{x(s, @x, \dfun)}\ge@r(\abs{y(s, @x, \dfun)})
\,\forall s!e[0, t].
\]
then lemma applies by assumptions 1 and 2 of the proposition, so that one can
majorize  $@F$ by a $\klfun!e\kl$

\newpage
\subsubsection{integral/integral UOSS}

the UOSS property gives uniform estimates on states as a function of
uniform bounds on outputs

$\exists$ ``finite energy output implies
finite energy state'' version as well.

\bd{iiUOSS}
system  {\em integral-to-integral
uniformly output-to-state stable (iiUOSS)} if $\exists$ \functions $\iy$,
$\iini$ $!e\kk$, and $\ix !e \ki$ s.t.
%
\beo{eq:iiUOSS}
\int_0^t \ix(\abs{x(s, @x, \dfun)})\,ds \psp !<\psp \iini(\abs@x) + 
          \int_0^t \iy (\abs{h(x(s, @x, \dfun))})\,ds
\eeo
$\forall$ initial state $@x$, disturbance $\dfun !e\M_@W$, and
time $t !e[0, \tmax(@x, \dfun))$
\ed

w.l.o.g., $\iy$ and $\iini$ can be assumed to be of class
$\ki$.

\newpage
\subsection{unboundedness observability}

\bd{fc} 
system {\em forward complete}
 if for $\forall$ initial 
condition $@x$, input $\uc$,  disturbance
$\dfun$ defined on $[0, +\infty)$, trajectory 
$x(t, @x, \uc, \dfun)$  
defined $\forall$ $t !>0$, i.e., $\tmax(@x, \uc, \dfun) = + \infty$
\ed

\spp

strictly weaker property, introduced by Mazenc et.al.:

\bd{uo} 
{\em unboundedness observability} property (UO) if
\beo{uo:eq}
\limsup_{t\nearrow \tmax(@x, \dfun)}\abs{y(t, @x, \dfun)} = + \infty
\eeo
holds for $\forall$ initial state
$@x$ and disturbance $\dfun$ with $\tmax(@x, \dfun) < \infty$
\ed

\spp
notice that UOSS $=>$ UO

(UOSS: $\abs{x(t, @x, \dfun)} !<\max \left\{ @b(\abs{@x}, t), 
\ylevel
\left(\norma{y}{[0, t]}\right)\right\}$)

\spp
%
following characterization of UO is from:

{\sc D. Angeli and E. D. Sontag}, {\em Forward completeness, unboundedness
  observability, and their Lyapunov characterizations}, Systems
  Control Lett., 38 (1999), pp. 209--217.

\spp
%
\bp{prop:UO}
UO iff $\exists$
$@r_1, @c_1, @c_2 !e\kk$ and $c!>0$ s.t.:
% the following implication holds: 
\beqn
\abso{h(x(t, @x, \dfun))} &\leq& @r_1(\abs{x(t, @x, \dfun)}) \quad 
\forall\,t !e[0, T]\nonumber\\[3mm]
& =>&  \abs{x(t, @x, \dfun)} !<@c_1(t) + @c_2(\abs @x) + c
 \quad \forall t !e[0, T]\label{eq:UO}
\eeqn
$\forall$ $@x!e\states$, $\dfun !e\M_@W$, $T !e[0, \tmax(@x, \dfun))$
\ep

gives uniform bound on states that can
be reached in given time, from a given bounded set, via 
a trajectory {\em not dominated by the output}

\spp
{\em note: in particular, characterization of forward completeness}

\newpage
\subsection{MAIN RESULTS SUMMARY}

\bt{main_thm:UOSS}
equivalent for $@.{x}=f(x(t), \dfun(t)), \; y(t)=h(x(t))$:
%\vspace{-0.2 in}
\begin{enumerate}
\item[$1$.]  $@S$ is UOSS. \label{item1:UOSS}
\item[$2$.]  $@S$ is \newprop. \label{item2:UOSS}
\item[$3$.]  $@S$ is iiUOSS and UO. \label{item3:UOSS}
\item[$4$.]  $@S$ admits a UOSS-Lyapunov function. \label{item4:UOSS}
%\item  $@S$ admits an UOSS-Lyapunov-like function.
%\item  $\exists$ state-norm estimator for $@S$.
\end{enumerate}
\ets

recall
\noindent{\bf Theorem 1}, for
$@.{x}(t)=f(x(t), \uc(t),\ddfun(t)), \; y(t)=h(x(t))$,
states equivalence of:
\begin{enumerate}
\item[$1$.] \label{item1} 
 $@S$ is UIOSS. 
\item[$2$.] \label{item2}
 $@S$ admits a UIOSS-Lyapunov function. 
\item[$3$.] \label{item3}
 $\exists$ state-norm-estimator for $@S$.
\end{enumerate}

\spp
\noindent{\bf organization of proof:}

Theorem~\ref{main_thm:UIOSS}:  1 $=>$ 2 will be reduced via a
``small gain argument'' (sections~\ref{ROSS:section},~\ref{nec})
to Theorem~\ref{main_thm:UOSS}: 2 $=>$ 4

%(which is proved in section~\ref{no_controls})

\spp

Theorem~\ref{main_thm:UIOSS}: 2 $=>$ 3 $=>$ 1 
in section~\ref{nobservers}

\spp

implications Theorem~\ref{main_thm:UOSS}: 3 $=>$ 2, 1 $=>$ 2, and 4 $=>$ 3, in
 section~\ref{impliesgasmo}

\spp

finally, Theorem~\ref{main_thm:UOSS}: 4 $=>$ 1 follows from
Theorem~\ref{main_thm:UIOSS}: 2 $=>$ 1 

\newpage
\section{reduction to the case of no controls}

we show reduction to particular case of systems with no controls

$\controls_1 =$ closed unit ball 
$\cbrace {u !e\controls : \abs u !<1}$  in $\controls$

\subsection {robust output-to-state stability} \label{ROSS:section}

\bd{ROSS}
system
$@.{x}(t)=f(x(t), \uc(t),\ddfun(t)), \; y(t)=h(x(t))$
{\em robustly output-to-state stable 
(ROSS) } if $\exists$ locally Lipschitz $\sm!e\ki$
(``{\em stability margin}'') s.t. system 
\be{ROSSsystem}
@.{x}(t)= g(x(t),\dfun(t)) := f(x(t), \dufun(t)\sm(\abs{x(t)}),\ddfun(t))
\ee
w/disturbances 
$d:= [\du, \dd] !e\controls_1 !p\Omegad$,
  outputs $y=h(x)$,
 is UOSS
\eds

\spp
notice $\controls_1 !p\Omegad$ is compact, convex subset of $\R^{m_u + m_w}$

denote it by $@W$ in this section

dynamics $g$ of system~\rref{ROSSsystem} 
locally Lipschitz in $x$ uniformly in $d$

and also $g(0, d) = 0$ $\forall$ $d !e@W$

\spp

\bl{IOSSimpliesROSS}
for $@.{x}(t)=f(x(t), \uc(t),\ddfun(t)), \, y(t)=h(x(t))$

UIOSS $=>$ ROSS
\els

%The proof will follow from a few preliminary lemmas

\bpr
let $@b!e\kl$ and $\ulevel$, $\ylevel !e\ki$ be s.t.

$\abs{x(t, @x, \uc, \ddfun)} !<\max \left\{ @b(\abs{@x}, t), 
\ulevel\left(\norma{\uc}{[0, t]}\right),
\ylevel\left(\norma{y}{[0, t]}\right)\right\}$

\spp
let $\vartheta(r) = @b(r, 0)$

w.l.o.g.  $\vartheta$ is $\ki$ and
$\vartheta(r) !>r$ (s.t. $\vartheta^{-1}(r)!<r$)

\spp
main point is showing that $\sm$ is stability margin, where

$\sm$ is any locally Lipschitz $\ki$-function $!<
\ulevel^{-1}(\frac{1}{4}\vartheta^{-1}(!.))$ 

%and can be extended as Lipschitz function to neighborhood of $[0, \infty)$

\newpage
\subsection{details of proof IOSS implies ROSS}

\bp{p:x(t)<x/2}
Fix a $@x!e\states$, a control $\uc$, and a
disturbance  $\ddfun$,
and let $x(!.):= x(!., @x, \uc, \ddfun)$ be the corresponding solution 
of the system
$@.{x}(t)=f(x(t), \uc(t),\ddfun(t)), \; y(t)=h(x(t))$
 
Let $T !e[0, \tmax(@x, \uc, \ddfun))$. 
Then
if $\abs{\uc(t)} !<\sm(\abs{x(t)})$ for almost all 
$t!e[0, T]$,
the estimate
\be{e:x(t)<x/2}
\abs{x(t)} !<\max\left\{@b(\abs{@x}, t),\;
\ylevel(\norma{y}{[0, t]}), \; \frac{\abs{@x}}{4}\right\}
\ee
holds $\forall$ $t !e[0, T]$.
\eps

\bpr

{\em Claim $1$.} Suppose $T < \tmax(@x, \uc, \ddfun)$. If 
\be{u<sm}
\abs{\uc(t)} !<\sm(\abs{x(t)}) \mbox{ for almost all }
t!e[0, T),
\ee
 then, $\forall$  $ t!e[0, T)$,
\be{est1}
\abs{x(t)} !<\max \left\{\vartheta(\abs{@x}), 
2\ylevel(\norma{y}{[0,t]})\right\}.
\ee

{\em Proof of the claim.} 
Suppose first that $@x= 0$. In this case 
$x(t) ==0$. Indeed,  define $\dufun(!.)$ on $[0, \tmax(@x, \uc, \ddfun))$ 
by
\[
\dufun(t) = \twoif{0}{ \mbox{ if } x(t) = 0,}
{\uc(t)/\sm(\abs{x(t)})}{\mbox{ if } x(t) \neq 0.}
\]
Then~\rref{u<sm} implies that $\dufun !e\M_{\controls_1}$ and that 
$x(!.)$ is the solution of~\rref{ROSSsystem} with $@x= 0$, $\ddfun$
as we picked, and 
$\dufun(!.)$ as we defined. Noticing that the constant function equal to $0$ 
is also a solution of this system, with the same initial state and the 
same disturbance, we conclude by  uniqueness of solutions that 
$x(t) = 0$ $\forall$ nonnegative $t$, so that~\rref{est1} trivially holds.  
  
 Suppose now that $@x\neq 0$.   
Fix $\ve$, s.t. $1<\ve<2$. We will first show
that if $\abs{\uc(t)} !<\sm(\abs{x(t)})$
 for almost all $t < T$,
then the estimate
%
\be{est2}
\abs{x(t)} !<\max\left\{\ve\vartheta(\abs{@x}),
 2\ylevel(\norma{y}{[0, t]})\right\}
\ee
%
holds $\forall$ $t !e[0, T)$.
Indeed, notice that~\rref{est2} is true as a strict inequality at
$t=0$ because 
$\abs{@x} !<\vartheta(\abs{@x}) < \ve \vartheta(\abs{@x})$. 
 If~\rref{est2} 
fails at some $t !e[0, T)$, then $\exists$   
\[
t_0=\min \left\{t < T \;:\; x(t)=\max\left\{\ve\vartheta(\abs{@x}),
2\ylevel(\norma{y}{[0, t]})\right\}\right\}.
\]
Note that $t_0  >0$ because at $t = 0$ we have a strict inequality 
in~\rref{est2}.
So,~\rref{est2} holds $\forall$ $t !e[0, t_0)$ and  
$ x(t_0)=\max\left\{\ve\vartheta(\abs{@x}), 2\ylevel(\norma{y}{[0, t_0]})\right\}$.
$\therefore$ for almost all $t !e[0, t_0)$ we have
\begin{eqnarray*}
\lefteqn{\ulevel(\norma{\uc}{[0,t]})\leq
\ulevel(\norma{\sm(\abs{x(!.)})}{[0, t]})}\\ 
 & & !<\max\left\{\frac{1}{4}\vartheta^{-1}(\ve\vartheta(\abs{@x})), 
\frac{1}{4}\vartheta^{-1}(2\ylevel(\norma{y}{[0, t]}))\right\}\\
& & !<\max\left\{\frac{1}{4}\ve\vartheta(\abs{@x}),
\frac{1}{4}2\ylevel(\norma{y}{[0, t]})\right\}  \\
 & & !<\max\left\{\vartheta(\abs{@x}),\ylevel(\norma{y}{[0, t]})\right\}.
\end{eqnarray*}
Then, since our system is UIOSS and $x(!.)$ is continuous, $\forall$
 $t !e[0, t_0]$ we have 
%
\begin{eqnarray*}
\abs{x(t)}&\leq& \max\left\{\vartheta(\abs{@x}), 
\ulevel(\norma{\uc}{[0, t]}),
\ylevel(\norma{y}{[0, t]}) \right\} \\
& = &\max\left\{\vartheta(\abs{@x}), \ylevel(\norma{y}{[0, t]}) \right\}.
\end{eqnarray*}
%
On the other hand, $\abs{x(t_0)}= \max\left\{\ve\vartheta(\abs{@x}),
2\ylevel(\norma{y}{[0, t_0]})\right\}$ by definition of $t_0$.
%
The contradiction proves the estimate~\rref{est2}. Letting $\ve$ tend to $1$,
we conclude that estimate~\rref{est1} holds 
$\forall$ $t !e[0, T)$, completing the proof of Claim 1.

Hence, under the assumption of Claim 1, we have
\begin{eqnarray*}
\ulevel\rb{\norma{\uc}{[0, t]}} &\leq& 
\ulevel\rb{\norma{\sm(\abs{x(!.)})}{[0, t]}}\\
&\leq&\max \cbrace{\frac{\abs{@x}}{4}, 
\frac{1}{4}\vartheta^{-1}\rb{2\ylevel\rb{\norma{y}{[0,
t]}}}}
\end{eqnarray*}
$\forall$ $t$ in $\maxint$. So,
\[
\abs{x(t)} !<\max\left\{@b(\abs{@x}, t),\;
\ylevel\rb{\norma{y}{[0, t]}}, \; \frac{\abs{@x}}{4}, \;
\frac{1}{4}\vartheta^{-1}\left(2\ylevel\left(\norma{y}{[0, t]}\right)
\right)\right\}
\] 
$\forall$ $t\in\maxint$.
Noticing that 
\[
\frac{1}{4}\vartheta^{-1}\rb{2\ylevel\rb{\norma{y}{[0, t]}}} \leq
\ylevel\rb{\norma{y}{[0, t]}}
\]
(because $\vartheta^{-1}(r) !<r$), we arrive at~\rref{e:x(t)<x/2}. 
\epr

\bll{lem10} Given any $\hat{@b}!e\kl$, $\exists$ 
$@b!e\kl$, $@n!e\ki$  s.t. for
 any $@t> 0$, any continuous 
\function $@m\,:\,[0, @t] \to \R_{!>0}$, and any nonnegative 
constant $C$, the following implication holds:
%
\be{first}
\forall t_1, t_2,\;\; 0 !<t_1 < t_2 !<@t, \;\;\;
@m(t_2) !<\max \left\{ \hat{@b}(@m(t_1), t_2 - t_1), \;
 \frac{@m(t_1)}{2}, \;
C \right\}
\ee
implies
\be{second}
@m(@t) !<\max \left\{ @b(@m(0), @t), \; \;
@n(C) \right\}.
\ee
\ells

{\it Proof}.
By Proposition 7 in:

{\sc E. D. Sontag}, {\em Comments on integral
  variants of input-to-state stability}, Systems Control Lett., 34
  (1998), pp. 93--100

$\exists$ $@m_1$ and $@m_2 !e\ki$ 
s.t. 
\[
\hat{@b}(r, t) !<@m_1(@m_2(r)e^{-t}),
\] 
so, by majorizing $\hat @b$ as above if necessary, we can assume
 w.l.o.g. that  $\hat{@b}$ is
continuous in its second variable, and
 $\hat{@b}(r, 0) !>r$ $\forall$ $r$. 
% 
$\forall$ $r>0$ define $T_r$ to be the first time when
 $\hat{@b}(r, T_r) = r/2$.
%
By replacing $\hat @b$ with the 
$\tilde @b!e\kl$, defined by 
%
\[
\tilde @b(r, t) := 
\max \cbrace{\hat @b(r, t), \, \hat @b(r, 0) e^{-t}},
\]
we can assume w.l.o.g. that  the series
$\sum_{i=0}^{\infty}T_{\frac{r}{2^i}}$ diverges for every $r > 0$.

Define {\afunction} $@f: \R_{!>0} !p\R_{!>0} \to \R_{!>0}$ as follows:
\[
@f(r, t) = \twoif {\hat{@b}(r, t)} {\mbox{for $t !e[0, T_r),$}}
{\hat{@b}\left(\frac{r}{2^k}, t - \sum_{i=0}^{k-1}T_{\frac{r}{2^i}}\right)}
{\mbox{for } \, t !e\left[\sum_{i=0}^{k-1}T_{\frac{r}{2^i}}, 
\sum_{i=0}^{k}T_{\frac{r}{2^i}}\right), \ k !> 1}
\]

Notice that the following two conditions hold for $@f$.

(1) For every $R, \, \ve > 0$, $\exists$ $\tilde t >0$ s.t. 
$@f(r, t) !<\ve$ $\forall$  $r < R$ and $t > \tilde t$. 

Indeed, fix positive $R$ and $\ve$ and find $k !e\Z$ s.t.
$\hat @b(R/2^k, 0) < \ve$. 
Next, by continuity of $\hat @b$ and by compactness of 
$[0, R]$ we can find a $\tilde t$, s.t. 
$\sum_{i=0}^{k-1}T_{\frac{r}{2^i}} < \tilde t$ $\forall$ positive $r < R$. 
Then, if $r< R$ and $t > \tilde t$, then 
\[
@f(r, t) =  
{\hat{@b}\rb{\frac{r}{2^k}, t - \sum_{i=0}^{k-1}T_{\frac{r}{2^i}}}} 
<  \hat @b\rb{\frac{r}{2^k}, 0} < \ve.
\]

(2) $\forall$ $\ve > 0$ $\exists$ $@d> 0$ s.t. 
if $r !<@d$, then $@f(r, t) < \ve$ $\forall$ $t !>0$.

Indeed, $\forall$ $r$ and $t$, $@f(r, t) 
!<\hat @b_0 (r) := \hat @b(r, 0)$. $\forall$ positive $\ve$, take 
$@d= @d(\ve) := \hat @b_0^{-1} (\ve)$. Then, $\forall$ 
$t !>0$ we have 
\[
@f(r, t) !<\hat @b\rb{\hat @b_0^{-1} (\ve) , 0} !<\ve.
\]   

$\therefore$, by Lemma cited from Lin et. al., $@f$ can be majorized by a 
$@b!e\kl$

Let $@n(r) = \hat@b(2r, 0)$. 

Now pick any $@m$, $C$, and $@t$ satisfying~\rref{first}. 
Define $T=\min\left\{t:\, @m(t) !<2C \right\}$ and $T=@t$ if
 $@m(t) > 2C$ $\forall$ $t !>0$. 

$\forall$ $t_1$ and $t_2$ in $[0, @t]$ s.t. 
$0 !<t_1 !<t_2 !<T$, we have $@m(t_1) >2C$, so that   
$@m(t_1)/2 > C$, hence 
%
\be{est:t<T}
@m(t_2) !<\max \left\{\hat@b(@m(t_1), t_2-t_1), @m(t_1)/2 \right\}.
\ee

Suppose now that $@t= T$. If $0!<@t< T_{@m(0)}$, then~\rref{est:t<T}
 with  $t_1 = 0$, $t_2 = @t$ yields
%
\[
@m(@t) !<\max\left\{\hat@b(@m(0), @t), \frac{@m(0)}{2} \right\}
= \hat@b(@m(0), @t),
\]
where the equality follows from the definition of $T_{@m(0)}$. Likewise,
if   
\[
 @t!e\left[\sum_{i=0}^{k-1}T_{\frac{@m(0)}{2^i}}, 
\sum_{i=0}^{k}T_{\frac{@m(0)}{2^i}}\right),
\]
 then
%
\begin{eqnarray*}
@m(@t) &\leq& \max\left\{\hat@b\left(2^{-k}@m(0), 
@t- \sum_{i=0}^{k-1}T_{\frac{@m(0)}{2^i}}\right), 2^{-(k+1)}@m(0)\right\}
 \\
&=& \hat@b\left(2^{-k}@m(0),
@t- \sum_{i=0}^{k-1}T_{\frac{@m(0)}{2^i}}\right),
\end{eqnarray*}
where the inequality follows from~\rref{est:t<T} and the equality 
is implied by the definition of $T_{\frac{@m(0)}{2^k}}.$
$\therefore$ we have 
\be{BD1}
@m(@t) !<@f(@m(0), @t).
\ee

In case $@t> T$, inequality~\rref{first} implies
\begin{eqnarray}
@m(@t)& !<& \max\left\{\hat{@b}(2C, @t-T), @m(T)/2, C\right\} 
\nonumber\\
&=& \max\left\{\hat{@b}(2C, @t-T), C, C\right\} !<
\hat{@b}(2C,0) = @n(C). \label{BD2}
\end{eqnarray}
Combining~\rref{BD1} and~\rref{BD2} we obtain
%
\[
@m(@t) !<\max\left\{@f(@m(0), @t), @n(C)\right\}
 \leq\max\left\{@b(@m(0), @t),@n(C)\right\}.\mybox
\] 


{\it Proof of Lemma IOSS implies ROSS}:
We need to show that the system~\rref{ROSSsystem}, corresponding to
our system

$@.{x}(t)=f(x(t), \uc(t),\ddfun(t)), \; y(t)=h(x(t))$

with the stability margin $\sm$ we have
 defined, is UOSS

apply Lemma~\ref{lem10} to $\hat @b:= @b!e\kl$
to find appropriate \functions $@b_1 !e\kl$ and $@n!e\kk$.
  Assume given any 
initial state $@x$ and disturbance $\dfun = [\dufun, \ddfun]$, and let
$x(t):= x(t, @x, \dufun, \ddfun)$ be the corresponding solution.
  Fix any positive $t < \tmax(@x, \dfun)$, and define $\uc$ by 
\[
\uc(s) := \twoif{\sm(\abs{x(s)}) \dufun(s),}{s !<t}{0,}{s > t.}
\]
Then, $\forall$ $s !<t$ we have $x(s) = x(s, @x, \uc, \ddfun)$, where the latter
is the solution of the original system
$@.{x}(t)=f(x(t), \uc(t),\ddfun(t)), \; y(t)=h(x(t))$
with control 
$\uc$ and disturbance $\ddfun$.    
Let $C=\ylevel(\norma{y}{[0, t]})$. Then, $\forall$ $t_1$ and $t_2$ in $[0, t]$
 we have $C !>\ylevel(\norma{y}{[t_1, t_2]})$. So, since
$\abs{\uc(s)} !<\sm(\abs{x(s)})$ $\forall$ $s !e[0, t]$, 

the above Proposition will imply that   
\[
  \abs{x(t_2)} !<\max \left\{ @b(\abs{x(t_1)}, t_2 - t_1), \; 
\frac{\abs{x(t_1)}}{2}, \;
C \right\}.
\]
By the choice of $@b_1$ and $@n$ we then have
\[
\abs{x(t)} !<\max\left\{@b_1(\abs{@x}, t), 
@n(\ylevel(\norma{y}{[0, t]})) \right\},
\]
proving the UOSS property for system~\rref{ROSSsystem} corresponding 
to the original UIOSS system. Thus,   
$\sm$
 is indeed a stability margin for the original system, and the proof of 
Lemma IOSS implies ROSS is now complete.
\mybox

{\bf end of details of proof IOSS implies ROSS}


\newpage
\subsection{a UIOSS system admits a UIOSS-Lyapunov function}\label{nec}

%not sure now if this is exactly what is being proved here:
%Theorem~\ref{main_thm:UIOSS}:  1 $=>$ 2 reduced to
%Theorem~\ref{main_thm:UOSS}: 2 $=>$ 4

show now how main implication of Theo~\ref{main_thm:UIOSS}
follows from Theo~\ref{main_thm:UOSS}

\sppp

(1) showed above: UIOSS $=>$ ROSS

\spp

(2) assume Theo~\ref{main_thm:UOSS}: 2$=>$4 already proved: 

$\exists$ UOSS-Lf $V$ for UOSS auxiliary system

$@.{x}(t)= g(x(t),\dfun(t)) = f(x(t), \dufun(t)\sm(\abs{x(t)}),\ddfun(t))$

\spp

(3) only left to prove that $V$ is 
UIOSS-Lf for original system

\spp
%
%\bl{}
%%[see Lemma $2.13$ in $\cite{ywscl}$] \label{lem8}
%
%if system $@S$ $@.{x}(t)=f(x(t), \uc(t),\ddfun(t)), \; y(t)=h(x(t))$
%is ROSS
%
%$V$ UOSS-Lf for system
%
%$@.{x}(t)= g(x(t),\dfun(t)) = f(x(t), \dufun(t)\sm(\abs{x(t)}),\ddfun(t))$
%
%associated to $@S$ $=>$  $V$ UIOSS-Lf for $@S$
%\els
%

\spp

%\bpr
% let $\sm$ be a stability margin for $@S$
%
%as $V$ UOSS-Lf for
%
%$@.{x}(t)= g(x(t),\dfun(t)) = f(x(t), \dufun(t)\sm(\abs{x(t)}),\ddfun(t))$
%
%inequalities 
%
%$@a_1(\abs{@x}) !<V(@x) !<@a_2 (\abs{@x})$,
%$\nabla V (x) !.f(x, d) !< -\xdissip(\abs x) + \ydissip(\abs{h(x)})$
%
%hold with some $@a_1$, $@a_2$, $\xdissip$, and $\ydissip$

pick state $@x!e\states$, disturbance value $\dd !e\Omegad$
and control value $u !e\controls$
 with $\abs{u} !<\sm(\abs @x)$ 

then $\exists$ $\du !e\controls_1$
 s.t. $u = \du \sm(\abs @x)$,  so that by the 

dissipation inequality for $V$ 
(applied with $d:=[\du, \dd]$) we have
\[
\nabla V(@x) !.f(@x, u, \dd) = \nabla V(@x) !.g (@x, d)
!< -\xdissip(\abs @x) + \ydissip(\abs {h(@x)}),
\]   
i.e.
\[
\abs @x!>\xurel(\abs u) \;=>\;
\nabla V(@x) !.f(@x, u, \dd) !< -\xdissip(\abs @x) + 
 \ydissip(\abs{h(@x)})
%\;\mbox{\small$\forall\,@x!e\states,\,u !e \controls,\,\dd !e\Omegad$}
\]
with $\xurel = \sm^{-1}$, so 
$V$ is UIOSS-Lf by previous Remark
%\epr

%By  Theorem~\ref{main_thm:UOSS}, the system~\rref{ROSSsystem} admits a
%UOSS-Lyapunov function $V$. Hence the following corollary follows.
%\bc{lem8:cor}
%if system
%$@.{x}(t)=f(x(t), \uc(t),\ddfun(t)), \; y(t)=h(x(t))$ is ROSS, then it admits a UIOSS-Lyapunov function
%\ecs
%
%\sppp
%
%By Lemma~\ref{IOSSimpliesROSS}, every UIOSS system is also ROSS, hence
%the implication 1 $=>$ 2 
%of Theorem~\ref{main_thm:UIOSS} follows.     
%} %end \forpaper

\forthesis{
   \bl{lem8} (See Lemma 2.13 in \cite{ywscl})
suppose system $@S$ $@.{x}(t)=f(x(t), \uc(t),\ddfun(t)), \; y(t)=h(x(t))$
is ROSS

Let $V$ be a UOSS-Lyapunov-like function for the system~\rref{ROSSsystem}
 associated with $@S$. Then  $V$
is an UIOSS-Lyapunov-like function for $@S$. 
\els
%
\bpr
 Let $\sm$ be a stability margin for $@S$.
Since $V$ is a UOSS-Lyapunov-like function for~\rref{ROSSsystem},
inequalities $@a_1(\abs{@x}) !<V(@x) !<@a_2 (\abs{@x})$
 and~\rref{UOSS-dissip} 
hold with some $@a_1$,
$@a_2$, $\xdissip$ and $\ydissip$.  Pick an initial state $@x$,
 control $\uc$, and disturbance $\ddfun$, and write
$x(t):= x(t, @x, \uc, \ddfun)$. Take $@t$, 
$@s!e[0, \tmax(@x, \uc, \ddfun))$.
  If $\forall$ $t !e[@s, @t]$ we have
 $\abs{x(t)}!>\sm^{-1}(\abs{\uc(t)})$, then 
$\uc(t) = \sm(\abs{x(t)})\dufun(t)$ for some $\dufun !e\M_{\controls_1}$,
 that is, $x(t)$
 is a trajectory
of the UOSS system~\rref{ROSSsystem} corresponding to $@S$,
with the disturbance $\dfun=[\dufun, \ddfun]$.
 Along this trajectory, the
 estimate~\rref{IOSS-dissip}
for the function $V$ is the same as~\rref{UOSS-dissip}. So, $V$
is a UIOSS-Lyapunov-like function for $@S$, with 
$\xurel = \sm^{-1}$, and $@a_i$ and $\ydissip$ as before.
\epr
%
\bc{lem8:cor}
if system $@.{x}(t)=f(x(t), \uc(t),\ddfun(t)), \; y(t)=h(x(t))$
is ROSS, then it admits a UIOSS-Lyapunov function
\ecs

\bpr
By the implication \ref{item1:UOSS} $=>$ \ref{item4:UOSS} 
of Theorem~\ref{main_thm:UOSS}, the system~\rref{ROSSsystem} admits a
UOSS-Lyapunov function $V$. So, $V$ is a smooth UOSS-Lyapunov-like function 
for~\rref{ROSSsystem}, hence, by the previous lemma, it is a smooth 
UIOSS-Lyapunov-like function for $@S$. Finally, by 
Proposition~\ref{if_smooth},
$V$ is an UIOSS-Lyapunov function for $@S$.
\epr

By Lemma~\ref{IOSSimpliesROSS}, every UIOSS system is also ROSS, hence
the implication \ref{item1} $=>$ \ref{item2}
of Theorem~\ref{main_thm:UIOSS} follows.
}% end \forthesis

\newpage
\subsection{UOSS or iiUOSS $=>$ \newprop}\label{impliesgasmo}

\bl{lem1} for $@.{x}=f(x(t), \dfun(t)), \, y(t)=h(x(t))$,
UOSS $=>$ \newprop\
\els

\spp

\bpr
have
$\abs{x(t, @x, \dfun)} !<\max \left\{ @b(\abs{@x}, t),\ylevel\left(\norma{y}{[0, t]}\right)\right\}$

w.l.o.g. 
$\ylevel!e\ki$, \ $\hbound(s) = @b(s,  0)> s$ $\forall$ $s > 0$

let $@r=$
any $\ki$-function s.t.
$@r(s)>\hbound(4\ylevel (s))$ $\forall$ $s > 0$

\spp

{\em Claim:} $\forall$
  $@x\in\states$,  $\dfun !e\M_@W$,  $@t\in[0,
  \tmax(@x, \dfun))$,
\[
\mbox{ if   }
\abs{x(t, @x, \dfun )} \geq@r(\abs{y(t, @x, \dfun )}) \quad
 \mbox{ $\forall$ } 0!<t!<@t,
\]
\[
\mbox{ then   }
\ylevel(\abs{y(t, @x, \dfun )})\leq
  \abs{@x}/2 \quad \mbox{ $\forall$ } 0!<t!<@t, 
\]
\[
\mbox{ and hence   } \ \ 
\abs{x(t, @x, \dfun )}\le@b(\abs{@x}, 0) 
 = \hbound(\abs{@x}) \quad \mbox{ $\forall$ } 0!<t!<@t. 
\]  
 in particular: if $\abs{x(t, \;
  @x, \dfun )} \geq@r(\abs{y(t, @x, \dfun )})$
 $\forall$ $t!e[0, \tmax(@x, \dfun))$,
  then $\tmax(@x, \dfun) = \infty$

\spp

{\em Proof of the claim:}.
clear for $@x= 0$, so pick $@x\not=0$, $\dfun !e\M_@W$
%assume that 
%$\abs{x(t, @x, \dfun )} \geq@r(\abs{y(t, @x, \dfun )})$
%  $\forall$ $0!<t\leq
%@t$ for some $@t!e(0, \; \tmax(@x, \dfun))$.  Then, a
\[
\mbox{at $t=0$: }
\ylevel (\abs{y(0, @x, \dfun )})\, \le\,
\ylevel (@r^{-1}(\abs{@x}))\leq\ylevel (@r^{-1}(\hbound(\abs{@x})))
< \abs{@x}/4
\]
$\therefore$ $\ylevel (\abs{y(t, @x, \dfun )}) < \abs{@x}/4$
$\forall$ $t\in[0,  @d)$ for some $@d> 0$
\[
\mbox{let\ }
t_1 := \inf\left\{t > 0: \ylevel (\abs{y(t, @x, \dfun )}) \ge
\abs{@x}/2\right\}
\;>\;0
\]
Assume by contradiction $t_1 \le @t$; $\therefore$
$\ylevel (\abs{y(t_1, @x, \dfun )}) = \abs{@x}/{2}$,
\[
\ylevel (\abs{y(t, @x, \dfun )}) < \abs{@x}/{2}
\quad\forall\;t!e[0, \, t_1)
\]
\[
=>\quad \abs{x(t, @x,d )} \le \hbound(\abs{@x})
\quad\forall\;t!e[0, \, t_1)
\]
\[
\therefore
\ylevel (\abs{y(t, @x, \dfun )}) !<\ylevel (@r^{-1} (\abs{x(t, @x, \dfun )}))
!<\ylevel (@r^{-1}(\hbound(\abs{@x}))) < \abs{@x}/4
\]
so by $!C^0$
$\ylevel (\abs{y(t_1, @x, \dfun )}) \le \abs{@x}/{4}$, contradicting 
def of $t_1$

so impossible to have $t_1\leq@t$, and claim proved.

\newpage


recall ``$\ve$-$@d$'' characterization of \newprop:

$\exists$ $@r!e\ki$ s.t.

(1) $\forall$ $\ve > 0$, $r > 0$, $\exists$ $T_{r, \ve}$ s.t. $\forall$
$\abs{@x}\le r$, $\dfun$, $T !e[T_{r,\ve}, \tmax(@x, \dfun))$,
\[
 \abs{x(t, @x, \dfun )} \;!>\;@r(\abs{y(t, \; @x, \dfun)})\;
\; \forall \ 0!< t!<T \,,
\]
\[
=>\;\;\abs{x(t,\; @x, \dfun )} \;< \;\ve\;
\; \forall \; t\in[T_{r,\ve}, T]\,.
\]
(2) $\exists$  $\hbound!e\kk$ s.t. $\forall$ $@x\in
\states$, $\dfun$, $T < \tmax(@x, \dfun)$ s.t.
\[
\abs{x(t,@x, \dfun )} \;!>\;@r(\abs{y(t,@x, \dfun)})\;
\; \forall \;0 !<t !<T\,,
\]
\[
=>\;\;\abs{x(T,@x, \dfun)} \;\leq\;\hbound(\abs@x)
\]

\spp

part (2) done; part (1): 

$\forall$ $r > 0$  let
$T_r!>0$ be s.t.
$@b(r, t) < r/2$ $\forall$ $t\ge T_r$

\spp

now, given any $r>0$ and any $\ve >0$, 

let $k(@e):=$ any positive integer s.t. $2^{-k(\ve)}r < \ve$

def $T_{r, \ve}:=T_{r_1} + T_{r_2} + !...  + T_{r_{k(\ve )}}$
\ (where $r_i:=2^{1-i}r$)

\spp

pick any trajectory $x(t, @x, \dfun )$ as in the statement of the

defined
on an interval of the form $[0,T]$,
with $T!>T_{r, \ve}$, with initial condition $\abs{@x} !<r$,
and disturbance $\dfun !e\M_@W$

satisfying $\abs{x(t, @x, \dfun )} \ge @r(\abs{y(t, @x, \dfun )})$
$\forall$ $t !e [0,T]$

$\Therefore$ the above claim implies that $\ylevel (\abs{y(t, @x, \dfun )})<
\abs{@x}/2$ $\forall$ such $t$

$\therefore$ $\forall$ $t>T_{r_1}=T_r$,
\beqn
\abs{x(t, @x, \dfun )} & \le & \max\left\{@b(\abs{@x}, t), \; \abs{@x}/{2}\right\} \\
&\le &\max \left\{@b(r,  t), \; r/{2}\right\} \, !< r/2 \,
\eeqn
consider now restriction of trajectory to interval $[T_{r_1},T]$

this $=$ trajectory that starts from
state $x(T_{r_1},@x, \dfun )$, 

which has norm less than $r_1$, so by the same
argument and the definition of $T_{r_2}$ we have that

$\abs{x(t, @x,d )}!<r/4$ $\forall$ $t!>T_{r_2}$

repeating on each interval $[T_{r_i},T_{r_{i+1}}]$, we conclude

$\abs{x(t,@x, d )} < \ve$ $\forall$ $T_{r,\ve}\le t\le T$
\eprr

\newpage

\bl{iiUOSSimpliesnewprop}
for $@.{x}=f(x(t), \dfun(t)), \, y(t)=h(x(t))$,

\ \ \ \ \ \ \ \ \ \ \ iiUOSS and UO $=>$ \newprop
\els

\spp

recall:

{\bf UO} $=>$ $\exists$ $@r_1, @c_1, @c_2 !e\kk$ and $c!>0$ s.t.:
% the following implication holds: 
\beqn
\abso{h(x(t, @x, \dfun))} &\leq& @r_1(\abs{x(t, @x, \dfun)}) \quad 
\forall\,t !e[0, T]\nonumber\\[3mm]
& =>&  \abs{x(t, @x, \dfun)} !<@c_1(t) + @c_2(\abs @x) + c
 \quad \forall t !e[0, T]
\eeqn
%$\forall$ $@x!e\states$, $\dfun !e\M_@W$, $T !e[0, \tmax(@x, \dfun))$

\spp

{\bf iiUOSS:} $\exists$ \functions $\iy$, $\iini$ $!e\kk$, $\ix !e \ki$ s.t.
%
\[
\int_0^t \ix(\abs{x(s, @x, \dfun)})\,ds \psp !<\psp \iini(\abs@x) + 
          \int_0^t \iy (\abs{h(x(s, @x, \dfun))})\,ds
\]
%$\forall$ initial state $@x$, disturbance $\dfun !e\M_@W$, and
%time $t !e[0, \tmax(@x, \dfun))$

one shows \newprop\ with
$@r(!.) := \max\cbrace{\ix^{-1}(2\iy(!.)), @r_1^{-1}(!.)}$


\newpage

\subsection{details of proof of iiUOSS and UO $=>$ \newprop}

need variant of ``Barb\u{a}lat's lemma'':

\bp{barbalat}
Let 
$!X:=\cbrace{x_@a, \; @a!e\A
}$
 be a family of absolutely continuous
curves in $\states$, each of which is defined 
on an interval $\I_@a$, either half-open 
($\I_@a= [0, @l_@a)$) or closed 
($\I_@a= [0, @l_@a]$). Suppose the following. 
\begin{itemize}
\item $!X$ is closed with respect to shifts, that is,
$\forall$ $@a!e\A$ and $T !e\I_@a$, $\exists$
an $@a' !e\A$ s.t.  
$x_{@a'} =={x_@a}_T$, where ${x_@a}_T$ is 
 defined by ${x_@a}_T (t) := x_@a(t +T)$, and 
$@l_{@a'} = @l_{@a} - T$. 
%  
\item $\exists$ nonnegative, increasing \function $@n_3$ s.t. 
\[
\abs{@.{x}_@a(t)}
 !<@n_3(\abs{x_@a(t)}) \;\; 
\forall \,@a!e\A \;\; 
\text{{\it for almost all}}
 \,t !e\I_@a.
\]
%
\item $\exists$ \functions  $@k$ and $@c$ $!e\ki$ s.t.
\[
@k(\abs{x_@a(0)}) !>\int_0^t @c(\abs{x_@a(s)})\,ds \quad 
\forall \, @a!e\A, \; t !e\I_@a.
\]
\end{itemize}
$\Therefore$ $\forall$ two positive numbers $r$ and $\ve$, $\exists$ $T_{r, \ve}$, 
s.t. $\forall$ $@a!e\A$ and $t !e\I_@a$ the following holds:
\[
t !>T_{r, \ve} \mbox{ and } \abs{x_@a(0)} !<r =>
  \abs{x_@a(t)} < \ve.
\]      
\eps

\newpage
\noindent

\bpr

{\em Claim $1$}. Given any $\ve > 0$, $\exists$ $@d= @d(\ve)$ s.t.

$\abs{x_@a(0)} !<@d \;=>\;\abs{x_@a(t)} < \ve$
$\forall\,t !eI_@a$
 
{\em Proof of Claim $1$}.  Fix $\ve>0$, and set:
\[
@d(\ve) = \min \left\{ \frac{\ve}{2}, 
 @k^{-1}\left(\frac{\ve@c(\ve/2)}{2@n_3(\ve)}\right)\right\}.
\]
pick any $@a!e\A$  s.t. $\abs{x_@a(0)} !<@d$
 
suppose $\abs{x_@a(\tilde t_2)} !>\ve$ for some 
$\tilde t_2 !e\I_@a$

 $\Therefore$
$\exists$ $t_1$ and $t_2$ with  $t_1 < t_2 !<\tilde t_2$
 s.t. 

$\abs{x_@a(t_1)} = \ve/2$
and $\ve/2 < \abs{x_@a(t)} < \ve$ $\forall$ $t !e(t_1, t_2)$
\[
\Therefore\;
\frac{\ve}{2} = \abs{x_@a(t_2)} - \abs{x_@a(t_1)} \leq
     \abs{x_@a(t_2)-x_@a(t_1)} 
\]
\[
!< \sup_{t_1 !<t !<t_2}
\abs{@.x_@a(t)}(t_2 - t_1) \leq     @n_3(\ve)(t_2- t_1).
\]
So,     
\[
@k(@d) !>@k(\abs @x) \;!>\; 
\int_{t_1}^{t_2} @c(\abs{x_@a(s)})\,ds 
> \frac{1}{2} @c\rb{\frac{\ve}{2}} (t_2 -t_1) 
\]
\[
!>\frac{\ve@c(\ve/2)}{2@n_3(\ve)} !>@k(@d).
\]
The obtained contradiction proves the claim.   
 
{\em Claim $2$}.  $\forall$ $r,@d>0$, $\exists$
time $@t(r, @d)$ s.t.    

$\abs{x_@a(0)} !<r$ and $@t(r, @d) !e\I_@a$ $=>$
$\exists t_0 < @t(r, @d)$ s.t.
 $\abs{x(t_0, @x, \dfun)} !<@d$
 
\spp

{\em Proof of Claim $2$}.  
Take $@t(r, @d) = \frac{2@k(r)}{@c(@d)}$. $\Therefore$, if 
$@t(r, @d) !eI_@a$ and $\abs{x_@a(t)} > @d$ $\forall$ 
$t !e[0, @t(r, @d))$, $\therefore$ we have 
%
\[
@k(r) !>@k(\abs{x_@a(0)}) !>
\int_0^{@t(r, @d)} @c(\abs{x_@a(s)})\,ds > 
@c(@d) \frac{@k(r)}{@c(@d)} = @k(r).
\]
The obtained contradiction proves the claim.

\newpage

Fix arbitrary positive $r$ and $\ve$. 

By Claim 1, 

find $@d(\ve)$ s.t.
$\abs{x_@a(0)} < @d(\ve)$ $=>$  $\abs{x_@a(t)} !<\ve$
$\forall$ $t !e[0, @l_@a)$

define
$T_{r, \ve} := @t(r, @d(\ve))$, where $@t(r, @d(\ve))$ is
as in Claim 2

If $\abs{x_@a(0)} < r$, 
by Claim 2, $\exists$
$t_0 < @t(r, @d(\ve))$ with $\abs{x_@a(t_0)} < @d(\ve)$

consider {\afunction} $x_{@a'}(!.):= x_@a(t_0 + !.)$ 
(it belongs to $\cal X$ by assumption)

since 
$\abs{x_{@a'}} !<@d(\ve)$, Claim 1 ensures that 
%
\[
\abs{x_{@a}(t)} = \abs{x_{@a'}(t -t_0)} !<\ve \quad \forall \;
t !>t_0 !>T_{r, \ve}.
\]  
$\therefore$ $T_{r, \ve}$  satisfies the
conclusion of the proposition


\epr 

\newpage

We now return to the proof of Lemma iiUOSS $=>$ \newprop

\spp

\bpr
\[
\io{@x, \dfun}:= \inf\{t !e[0, \tmax(@x, \dfun)): 
\abs{x(t,@x, \dfun )} !<@r(\abs{y(t,@x, \dfun)})\}
\]
and $\io{@x,\dfun} = \tmax$ if $\abs{x(t,@x, d )} > @r(\abs{y(t,@x, \dfun)})$ $\forall$
$t !e[0, \tmax)$

note: given $@x$ and $\dfun$, $\forall$ $t < \io{@x, \dfun}$
we have 

$\ix(\abs{x(t, @x, \dfun)}) > 2 \iy(\abs{h(x(t, @x, \dfun))})$, so
that:
%
\[
\qquad \iini(\abs @x) !>
\int_0^t\left(\ix(\abs{x(s, @x, \dfun)}) - \iy(\abs{h(x(s, @x,
\dfun))})\right)\,ds 
\]
\[
> 
\frac{1}{2}\int_0^t \ix(\abs{x(s, @x, \dfun)})\,ds.
\]
Let $@n_3!e\kk$ s.t.
$\max_{d\in@W} \abs{f(x, d)} !<@n_3(\abs x)$

Write $x_{@x, \dfun}(!.) := x(!., @x, \dfun)$

family 
$\cbrace{x_{@x, \dfun}(!.),\; @x!e\states,\, \dfun !e\M_@W}$
with $\I_{@x, \dfun} := [0, \io{@x, \dfun})$ satisfies all the assumptions
 of previous Proposition
(with ``$@k$'' $= 2\iini$)

given any positive $r, \ve$, above Proposition furnishes 
$T_{r, \ve}$

$T_{r, \ve}$ satisfies first condition in the
$\ve$-$@d$ characterization of \newprop\

to find {\afunction} $\hbound$ for second part of that
 characterization, recall that, 

UO $=>$ $\exists$ $\ki$-functions $@r_1$,
$@m_1$, $@m_2$, and a constant $c > 0$ s.t.

$\forall$ $@x!e\states$, $\dfun !e\M_@W$, $T !e[0, \tmax(@x, \dfun))$:
% 
\beqn
\abso{h(x(t,@x, \dfun))} &\leq& @r_1(\abs{x(t,@x, \dfun)}) \;\; 
\forall t !e[0, T] \\ 
&=>&
\abs{x(t,@x, \dfun)} !<@m_1(t) + @m_2(\abs @x) + c \;\; 
\forall t !e[0, T].
\eeqn
$\therefore$, if $\abs @x!< r$ and  $T_{r, r/2}$ is as defined above, $\therefore$
$\forall$ $t !e[0, \io{@x, \dfun})$ we have
%
\[
\abs{x(t, @x, \dfun)} !<@m_1(T_{r, r/2}) + @m_2(r) +c
\;\;\mbox{if } t < T_{r, r/2} 
\]
and 
\[
\abs{x(t, @x, \dfun)} !<r/2 
\;\;\mbox{if } t !>T_{r, r/2}. 
\]
Thus, the following estimate holds $\forall$ such $t$:
\[
\abs{x(t, @x, \dfun)} !<\tilde \hbound(\abs @x) :=
\max\cbrace{\abs @x/2,\;  
@m_1(T_{\abs @x, \abs @x/2}) + @m_2(\abs @x) +c}.
\]

Next, take a sequence $\cbrace{\ve_k}, \; k=0, \,1,\, 2,!...,$ strictly 
decreasing to $0$, with $\ve_0 = 1$. $\forall$ $\ve_k$, find 
$@d_k = @d(\ve_k)$ as in the proof of Claim 1. Since 
$@d_k !<\ve_k/2$,
the sequence $\cbrace{@d_k}$ converges to $0$ as well. Find {\afunction}
$\hbound$ $!e\kk$, s.t.

(1) $\hbound(@d_{k+1}) >  \ve_k \;\; \forall k >0,$ \newline
(this will ensure that 
$\abs{x(t, @x, \dfun)} !<\hbound(\abs@x) \mbox{ $\forall$ } @x\mbox{ with} 
\abs @x< @d_0, \mbox{ $\forall$ } t !e[0, \io{@x, \dfun}]$) 

(2) $\hbound(s) !>\tilde \hbound(s) \;\;  \forall s > @d_0$.

$\Therefore$ $\hbound$ satisfies the second condition in the 
$\ve$-$@d$ characterization of \newprop\

\eprr

\noindent{\bf end of proof iiUOSS and UO $=>$ \newprop}

\newpage

\subsection{UO needed for ii statement}

example of iiOSS system which
 fails to be OSS 

(equivalently, fails to be \newprop)

%$1_A(!.):=$ indicator function of a set $A$, and let
 $@f_\ve:=$ $C^\infty$-bump function with support in
$(-\ve, \ve)$:
%
\be{bump}
@f_\ve(@x):= \twoif{e^{-\frac{\abs @x^2}{\ve^2-\abs @x^2}},}
{\abs @x< \ve,}{0,}{\abs @x!>\ve.}
\ee
%  
fix an arbitrary positive $\ve < 0.25$ and consider one
dimensional autonomous system:
\[
@S: \quad\quad @.x = f(x) ,\quad y = h(x),
\] 
where
%
\begin{eqnarray*}
 f(x)& = &x^3\left[1_{(-\infty, -1]}(x)(1 - @f_\ve(x+1)) + 
1_{[1, +\infty)}(x)(1 - @f_\ve(x-1)) \right] \\
&&-x\left[1_{(-1, 1)}(x) (1 - @f_\ve(x+1))(1-@f_\ve(x-1))\right],
\end{eqnarray*}
%
and $h!e!C^!8$ s.t. $h(x) = x$ $\forall$ $x$ in
 $[-2, 2]$, and $h(x) = 0$ if $\abs{x} !>3$.
 
\includepicx{0.6}{picture1}

\newpage

{\em Claim $1$}. System $@S$ is not OSS. 

{\it Proof}.
pick initial state $@x>0$ large enough s.t. $h(@x) =0$

$\Therefore$ $h(x(t, @x))=0$ $\forall$ $t < \tmax(@x) =
@x^{-2}/2$

If $@S$ were OSS, there would exist
$@b!e\kl$ s.t. 
$\abs{x(t, @x)} !<@b(\abs@x, t)$, but $\abs{x(t, @x)}$ tends to
$\infty$ as $t \to \tmax(@x)$, 

whereas 
$@b(\abs@x, t) !<@b(\abs@x, 0)$. 

This contradiction proves
the claim.

\newpage


{\em Claim $2$}. The system $@S$ is iiOSS.

{\em Proof}. $@S$ has stable equilibrium at $x = 0$ and two
unstable ones at $1$ and $-1$

if $\abs{x} < 1$, $\sign(x) = - \sign(f(x))$, so, if $\abs @x!<1$, then
$\abs{x(t, @x)} \leq 1$ $\forall$ nonnegative $t$

Thus, if $@x!e[-1, 1]$, then
 $\forall$ $t !>0$:
%
\be{x=h}
 \int_0^t \abs{x(s, @x)}\,ds = \int_0^t \abs{h(x(s, @x))}\,ds,
\ee
%
so,  iiUOSS estimate trivially follows $\forall$ $@x!e[-1, 1]$,
$t !e\tmax(@x)$ with $\iy = Id$, $\iini !e\kk$

\spp

on $\abs@x!>1+\ve$, $f(x) = x^3$, so
\[
x(t, @x) = \frac{\sign(@x)}{\sqrt{@x^{-2} -2t}}.
\]
Thus, in this
case the solution $x(t, @x)$ is defined $\forall$ nonnegative
$t < \tmax(@x)= @x^{-2}/2$ and 
\[
\int_0^t \abs{x(s, @x)}\,ds  !<\int_0^{\tmax(@x)} \abs{x(s, @x)} \,ds
 = \frac{1}{@x} !<\frac{1}{1+\ve}.
\] 
and OK again (just $@b$ term used)

\spp
now suppose $1< \abs{@x}< 1+\ve$. 

let $\iini!e\kk$ s.t.
$\iini(1) !>(1+\ve)^{-1}$. 

Let $\hat t$ be the time when
$\abs{x(\hat t, @x)} = 1+\ve$. $\Therefore$ 
$\tmax(@x) = \hat t + (1+\ve)^{-2}/2$. Also, $x(s, @x) = h(x(s, @x))$ $\forall$ $s !e[0, \hat t]$, so, in particular, equality~\rref{x=h} holds $\forall$ 
$t < \hat t$, which, again, trivially implies iiUOSS    
with $\iy = Id$ and any $\iini !e\kk$.

If $t > \hat t$:
% 
\begin{eqnarray*}
\int_0^t \abs{x(s,@x)}ds &=& \int_0^{\hat t}
\abs{x(s,@x)}ds + \int_{\hat t}^{t} \abs{x(s,@x)}ds  \\
&\leq& \int_0^{\hat t}
\abs{x(s,@x)}ds + \int_{\hat t}^{\tmax(@x)} \abs{x(s,@x)}ds  \\ 
&=& \int_0^{\hat t} \abs{h(x(s,@x))}ds + \int_0^{\tmax(1 +\ve)}
\abs{x(s,1 + \ve)}ds \\
&\leq& \int_0^{\hat t} \abs{h(x(s,@x))}ds + \iini(@x)\\
& !<& \int_0^t \abs{h(x(s,@x))}ds + \iini(@x).
\end{eqnarray*}
This shows that $@S$ is iiOSS, as iiOSS estimate holds for 
$@S$ with the $\iini$ that we constructed and $\iy = Id$.\mybox

\spp
\er

\noindent{\bf end of counterexample}

\newpage

\subsection{implication 4 $=>$ 3
of Theorem~\protect{\ref{main_thm:UOSS}}}

i.e. 
$@S$ admits a UOSS-Lyapunov function $=>$ iiUOSS and UO

\bpr
already 
remarked $@S$ must be UO, so need iiUOSS

by 
assumption $\exists$  smooth \function
$V$ satisfying

$\nabla V (x) !.f(x, d) !< -\xdissip(\abs x) + \ydissip(\abs{h(x)})$

and $@a_1(\abs{@x}) !<V(@x) !<@a_2 (\abs{@x})$
 with some
$@a_1$, $@a_2$, $\xdissip$, and $\ydissip$

 pick any $@x$, $\dfun$, and
$t !e[0, \tmax(@x, \dfun))$

integrating 
inequality
$\nabla V (x) !.f(x, d) !< -\xdissip(\abs x) + \ydissip(\abs{h(x)})$

along the trajectory 
$x(!., @x, \dfun)$ over $[0, t]$ we get
%
\begin{eqnarray*}
&&\int_0^t \xdissip(\abs{x(t, @x, \dfun)})dt \\
&\leq&
        V(x(0, @x, \dfun)) -V(x(t, @x, \dfun))+ 
          \int_0^t \ydissip(\abs{h(x(t, @x, \dfun))})\,dt \\
&\leq& @a_2(\abs@x) + \int_0^t \ydissip(\abs{h(x(t, @x, \dfun))})\,dt,
\end{eqnarray*}
proving the iiUOSS inequality
with $\ix = \xdissip$ and $\iini = @a_2$
\eprr

\newpage
\subsection{a remark on the \newprop\ $=>$ UOSS implication}

\br{alt_proof}
once Lyapunov characterization for UOSS is proved,
implication \newprop $=>$ UOSS of
Theorem~\ref{main_thm:UOSS} will automatically follow by appealing to the
converse Lyapunov theorem 
(2~$=>$~4~$=>$~1)

however, \newprop\ $=>$ UOSS 
can easily be proved directly:
 
fix $@x\in\states$ and
$\dfun !e\M_@W$

take any $t !e[0, \tmax(@x, \dfun))$. 

$\io{@x, \dfun}= \inf\{t !e[0, \tmax(@x, \dfun)): 
\abs{x(t,@x, \dfun )} !<@r(\abs{y(t,@x, \dfun)})\}$

If $t !<\io{@x, \dfun}$, 
\newprop\ property provides estimate
\be{case_less_io}
\abs{x(t, @x, \dfun)} !<\klfun(\abs @x, t).
\ee
suppose now  $t > \io{@x, \dfun}$. Obviously,
 at least one of the two following conditions
must be satisfied:

(1) $\abs{x(t, @x, \dfun)} > 2@r(\abs{h(x(t, @x, \dfun))})$,

(2) $\abs{x(t, @x, \dfun)} !<3@r(\abs{h(x(t, @x, \dfun))})$.

\noindent if condition (2) applies, we have a bound
\be{if_case_2}
\abs{x(t, @x, \dfun)} !<3@r(\abs{h(x(t, @x, \dfun))})
!<3@r(\norma{y}{[0, t]}).
\ee

\noindent In case condition (1) holds, let
\[
\tilde t := \max \cbrace{s, \, \io{@x, \dfun} < s < t\, : \,
 \abs{x(s, @x, \dfun)} = 2@r(\abs{h(x(s, @x, \dfun))})}.
\]       
$\Therefore$ again by the \newprop\ property applied with initial state
$x(\tilde t, @x, \dfun)$, we have
\beq
\abs{x(t, @x, \dfun)} &\leq&
\klfun(\abs{x(\tilde t, @x, \dfun)}, t - \tilde t)\nonumber\\ &\leq&
\klfun(2@r(\abs {y(\tilde t, @x, \dfun)}), 0) \leq
\klfun(2@r(\norma{y}{[0, t]}), 0).\label{if_case_1}
\eeq

Combining estimates~\rref{case_less_io}, \rref{if_case_2},
and~\rref{if_case_1}, we conclude:

$\abs{x(t, @x, \dfun)} !<\max \left\{ @b(\abs{@x}, t),
\ylevel\left(\norma{y}{[0, t]}\right)\right\}$

holds for $@S$ with
$@b(!.) := \klfun (!.)$ and

$\ylevel(!.) :=
\max \cbrace{\klfun(2@r(!.), 0), \; 3@r(!.)}$
\er   

\newpage

\subsection{construction of $V$}

very long and detailed argument, so only brief
outline provided:

suppose $(@S)$ $@.{x}=f(x(t), \dfun(t)), \, y(t)=h(x(t))$
\newprop:
\beqn
&&
\hspace{-100pt}
\abs{x(t, @x, \dfun)} !>@r(\abs{h(x(t, @x, \dfun))}) 
\,\;\;(\forall \; 0 !<t !<T)
\\
&=>&
\abs{x(t, @x, \dfun)} !<\klfun (\abs{@x}, t) \;\;\; (\forall\; 0 !<t!<T)
\eeqn
($@r!e\ki, \;\klfun!e\kl$;
$\forall$   $@x!e\states$, $\dfun !e\M_@W$, $T < \tmax(@x, \dfun)$)

introduce
%
\[
\D \,:=\;\left\{@x!e\states \colon \;\; \abs{@x} \leq
@r(\abs{h(@x)})\right\}
\]

$\forall$ $\dfun !e\M_@W$ and $@x\not!e\D$, define
\beo{e-d-io}
\io{@x, \dfun} = \inf \left\{t !e\maxint \;\colon\;
 x(t, @x, \dfun) !e\D\right\},
\eeo
with convention $\io{@x, \dfun} = \tmax(@x, \dfun)$ if 
trajectory never enters $\D$

The \newprop\ property then implies  
\beo{e-gasmob}
\abs{x(t, @x, \dfun)} \le \klfun(\abs{@x}, t) \qquad \forall\,@x\in\E, \
\forall\, \dfun \in\M_@W, \ \forall\, t!e[0, \io{@x, \dfun})
\eeo
for some $\klfun\in\kl$

and this provides the necessary integrability estimates along trajectories,
needed in what follows, using Proposition 7 in:

{\sc E. D. Sontag}, {\em Comments on integral
  variants of input-to-state stability}, Systems Control Lett., 34
  (1998), pp. 93--100.

which says that $\exists$ $\ki$-functions $@m_1$ and $@m_2$ s.t. 
\beo{mu1mu2} 
\klfun(r, t) !<@m_1( @m_2(r) e^{-t}) \qquad \forall \, r,\ t !>0.
\eeo
(or, redefining $@l$, $\klfun(r, t) !< @m_2(r) e^{-t}$, and hence
$\klfun(r, t)e^{t}!e!L^1$)

key: find a $V$ which decreases while trajectories are outside $\D$

(but note that $\D$ is not stable, or even invariant)

modify dynamics near $\D$ by making system 

completely controllable (adding controls $u$) 

next define a function $V_0$ by solving an optimal control problem

$\sup_d\inf_u \int L(x(t,d,u))\,dt$ 

along all trajectories that stay outside 

(allow relaxed disturbances to compactify space of disturbances; 

affine in $u$ so also (weak) compactness in $u$)

since system modified, $V_0$ for the original system will only decrease while
outside a ``thick'' version of $\D$, but this is OK (redefine $@r$)

finally, use nonsmooth analysis techniques (inf-convolution) to 

obtain a locally Lipschitz $V$, and from there 

a smoothing technique via standard convolutions and partition of unity to
obtain smooth $V$

\newpage
\section{norm-observers}\label{nobservers}

let ($@S$)\ $@.{x}(t)=f(x(t), \uc(t),\ddfun(t)), \; y(t)=h(x(t))$
 
$C^1$ $V: \states \to \R_{!>0}$ is
{\em exponential decay UIOSS-Lf} 
for $@S$

if  $@a_1(\abs{@x}) !<V(@x) !<@a_2 (\abs{@x})$
 ($@a_1,@a_2!e\ki$) and
% the following version of inequality~\rref{pqUIOSS-dissip}:
\beo{exp-pqUIOSS-dissip}
\nabla V(x) !.f(x, u, \dd) \;!<\; - V(x) + \ugdissip(\abs{u}) + 
\ygdissip(\abs{h(x)}) 
\eeo

($\forall x !e\states,  u !e\controls,  \dd !e\Omegad$,
with $\ugdissip,\ygdissip !e\kk$)
\eds

\spp
\bl{IOSS:exp-decay}
$V$ UIOSS-Lf $=>$ $\exists$ $@r!e\ki$ s.t.
$W:=@r!oV$ is expo decay UIOSS-Lf
\els 

then
\beo{ne} 
\nobserver: \quad  @.p = -p + \ugdissip(\abs u) + \ygdissip(\abs y),
\quad
k(s, r)=\abs{s}
\eeo
is a norm-estimator

\eps
\bpr
assume w.l.o.g. $@a_2$ in def of $V$
satisfies $r !<@a_2(r)$

system~\rref{ne} 
 is ISS w.r.t. $u,y$, since asymptotically
stable linear system driven by input 
$(\ugdissip(\abs u),\ygdissip(\abs y))$

\spp

pick any
initial states $@x$, $@z$,
control $\uc$, disturbance $\ddfun$;

then for trajectory
$(x(t), p(t))$ of the composite
system:
 \beo{V-z_dissip}
\frac{d}{dt}(V(x(t))-p(t)) !<
 - (V(x(t))-p(t))
\eeo
for almost all $t !e[0, \tmax(@x, \uc, \ddfun))$.
Thus
\[
V(x(t)) \;!<\;
p(t) + e^{-t}(V(@x)-@z) \;!<\;
\abs{p(t)} + 2e^{-t}@a_2(\abs{@x}+\abs{@z})
\]
(using $r!<@a_2(r)$)

This can be written in def of estimator
with $@r:=@a_1^{-1}(2(!.))$
and $@b(s,t):=@a_1^{-1}(4e^{-t}@a_2(s))$.

\newpage

\subsection{$\exists$ estimator $=>$ UIOSS}

assume that $(\nobserver, k)$ is norm-estimator for $@S$

pick any initial state $@x$ for $@S$, any input $\uc$, 
disturbance $\ddfun$, and the special initial state $@z= 0$ 
for $\nobserver$

so
\be{normeq1}
 {k(p(t,0,\uc, \yc_{@x,\uc, \ddfun}),  y(t, @x, \uc, \ddfun))}
\,!<\,
 \uobs\left(\norma{\uc}{[0,t]}\right) +
 \yobs\left(\norma{\yc_{@x,\uc, \ddfun}}{[0, t]}\right)
\ee
$\forall$ $t!e\maxint$, with some  \functions $\uobs,\yobs!e\kk$
(the $\kl$-term vanishes because $@z= 0$)

$\Therefore$ combining~\rref{normeq1}
 with the norm-estimate:
\[
\abs{x(t, @x, \uc, \ddfun)} \leq @b(\abs @x, t) + 
@r({k(p(t,0,\uc, \yc_{@x,\uc, \ddfun}), y(t, @x, \uc, \ddfun) )})
\]
\[ !<
\max\cbrace{ 2@b(\abs@x, t),\, 
4@r(\uobs(\norma{\uc}{[0, t]})), \,
4@r(\uobs(\norma{y_{@x, \uc, \ddfun}}{[0, t]}))}
\]
showing UIOSS
\epr

\newpage

\subsection{details of construction of $V$}

By majorizing $@r$ with
another $\ki$-function if necessary, we will assume that
$@r$ is smooth
%
when restricted to $s>0$
and also
  $@r(s) > s$ $\forall$ positive $s$. 

let

\[
\E \,:=\;\X\backslash \D \,,
\]
and
\[
\E_1 \,:=\, \left\{@x!e\X \; \colon \; 
\abs{@x} > 2@r(\abs{h(@x)})\right\} \,.
\]
%

%
If $\D = \states$, $\therefore$ any proper, smooth, and positive definite \function 
$V :\states \to \R$ is a UOSS-Lyapunov function for~\rref{dsystem}. 
Indeed, because it is proper and finite,  $V$ obviously 
satisfies $@a_1(\abs{@x}) !<V(@x) !<@a_2 (\abs{@x})$
 for some $@a_1$ and $@a_2$.
Since $V$ is smooth, $\abs{\nabla V(@x)}$ is bounded above by a nondecreasing
continuous \function $@n(\abs{@x})$ and 
\[
\frac{d}{dt}V(x(t)) \nsp \psp =\nsp \psp \nabla V(x(t)) !.f(x(t), \dfun(t))
 \nsp \psp !<\nsp \psp
@n(\abs{x(t)})@n_3(\abs{x(t)}),
\]
where $@n_3(\abs{!.})!e\kk$ majorizes $f(!., d)$ 
$\forall$ $d !e@W$. 
$\Therefore$ since $\abs{x} !<@r(\abs{h(x)})$ $\forall$ $x !e\states$,
we have
%
\[
\frac{d}{dt}V(x(t))
\;\leq\;
 - @n(\abs{x(t)})@n_3(\abs{x(t)}) \nsp \psp + \nsp \psp
2@n(@r(\abs{h(x(t))}))@n_3(@r(\abs{h(x(t))})) \,.
\]
So, $V$ satisfies inequality
\[
\nabla V (x) !.f(x, d) !< -\xdissip(\abs x) 
+ \ydissip(\abs{h(x)}) \quad \forall\,x !e\states, \; \forall\,d 
!e@W
\]
(with 
$\xdissip(!.) = @n(!.) @n_3(!.)$ and
$\ydissip(!.) = [2 @n!o@r(!.)][ @n_3 !o@r(!.)]$),
which is the same as~\rref{pqUIOSS-dissip} for systems of type~\rref{dsystem}. 

Suppose now that $\D \neq \states$. 
Recall that we have defined, $\forall$  $@x\not!e\D$ and 
$d !e\M_@W$,
$\io{@x, \dfun} = \inf \left\{t !e\maxint \;\colon\;
 x(t, @x, \dfun) !e\D\right\},$
with the convention $\io{@x, \dfun} = \tmax(@x, \dfun)$ if the trajectory never 
enters $\D$.

The \newprop\ property then implies  
\be{e-gasmob}
\abs{x(t, @x, \dfun)} \le \klfun(\abs{@x}, t) \qquad \forall\,@x\in\E, \
\forall\, \dfun \in\M_@W, \ \forall\, t!e[0, \io{@x, \dfun})
\ee
for some $\klfun\in\kl$.

Note that, because of property~\rref{e-gasmob}, the system
cannot have any equilibrium in $\E$, that is, 
\[
f(@x,  d) !=0
\]
for every $@x!e\E$ and every $ d\in@W$.
Moreover, replacing $@r(s)$ by $c@r(s)$
for some $c>1$ if necessary, one may also assume that $f(@x, d)\not=0$
$\forall$ $@x\in\partial\D\setminus\{0\}$, $d\in@W$
%

We introduce an auxiliary system $\hat @S$ which 
slows down the motions of the original one:
%
\be{slowsystem}
@.{z}\;=\;\fhat(z, d)\; =\; \frac{1}{1+\abs{f(z, d)}^2 + @k(z)} f(z, d),
\ee
%
where $@k$ is any smooth \function $\states \to [0, \infty)$ with the 
property
that 
\be{kappa_def}
@k(@x) \;!>\;2 \max_{d \in@W} 
\abs{\nabla(@r\circ
\abs{h})(@x) !.f(@x, d)}
\ee
%
whenever $\abs{h(@x)} !>1$. 
(Recall that $@r$ was assumed, w.l.o.g., to be smooth for
positive arguments.)
$\forall$ disturbance $\hat \dfun$ 
(defined on $\R_{!>0}$) denote by
\[
z(s, @x, \hat{\dfun})
\]
the value at time $s$ of the solution
of the equation $@.{z}=\fhat(z, \hat{\dfun})$ with initial state
$@x$. Observe that, as $\fhat$ is bounded, this solution exists $\forall$ 
nonnegative $s$.


{\em Claim $1$}. $\forall$ $@x$ and each $\dfun$,
\be{e-xz}
x(t, @x, \dfun) = z(\tslow{@x}{\dfun}(t), @x, \dfun \circ\tslow{@x}{\dfun}^{-1})
\qquad \forall\,t\in[0, \tmax(@x, \dfun)),
\ee
where
$\tslow{@x}{\dfun}: [0, \tmax(@x, \dfun)) \to \R_{!>0}$ is defined by
%
\[
\tslow{@x}{\dfun}(t) = \int_0^t 
\sqb{1+\abs{f(x(s, @x, \dfun), \dfun(s))}^2 + @k(x(s, @x, \dfun))}ds.
\]
Moreover,
$\tslow{@x}{\dfun}(t)->\infty$ as $t->\tmax(@x,\dfun)$,
so, we can define $\tslow{@x}{\dfun}(\tmax(@x, \dfun)) := +\infty$ 
for convenience.

{\it Proof of Claim $1$}.
Indeed, writing $s = \tslow{@x}{\dfun}(t)$ and computing the derivative
of $x(\tslow{@x}{\dfun}^{-1}(s), @x, \dfun)$ with respect to $s$, one
has
\begin{eqnarray*}
f(x(t, @x, \dfun), \dfun(t)) &=& \frac{d}{dt}x(t, @x, \dfun)
 = \frac{d}{dt}x(\tslow{@x}{\dfun}^{-1} !o\tslow{@x}{\dfun}(t), @x, \dfun) \\
& = & \frac{d}{ds}x\rb{\tslow{@x}{\dfun}^{-1}(s), @x, \dfun}!.
\frac{d}{dt} \tslow{@x}{\dfun}(t)
\end{eqnarray*} 
\[ 
=  \frac{d}{ds}x(\tslow{@x}{\dfun}^{-1}(s), @x, \dfun)
       \sqb{1+\abs{f(x(t, @x, \dfun), \dfun(t))}^2 + @k(x(t, @x, \dfun))}.
\]
$=>$
\begin{eqnarray*}
& &\frac{d}{ds}x(\tslow{@x}{\dfun}^{-1}(s), @x, \dfun)  = 
\frac{f(x(t, @x, \dfun), \dfun(t))}
{1+\abs{f(x(t, @x, \dfun), \dfun(t))}^2 + @k(x(t, @x, \dfun))}\\[2mm]
& &=
\frac{f \rb{x(\tslow{@x}{\dfun}^{-1}(s), @x, \dfun), \dfun !o
\tslow{@x}{\dfun}^{-1}(s)}}
{1+\abs{f\rb{x(\tslow{@x}{\dfun}^{-1}(s), @x, \dfun), 
\dfun !o\tslow{@x}{\dfun}^{-1}(s)}}^2 
+ @k\rb{x\rb{\tslow{@x}{\dfun}^{-1}(s), @x, \dfun}}} \\[2mm]
& &\qquad =\widehat f\rb{x\rb{\tslow{@x}{\dfun}^{-1}(s), @x, \dfun}, \dfun 
!o\tslow{@x}{\dfun}^{-1}(s)}
\forall\,0\le s <
\tslow{@x}{\dfun}(\tmax(@x, \dfun)). 
\end{eqnarray*}
Thus, the functions $z(s, @x, \dfun !o\tslow{@x}{\dfun}^{-1})$ and
$x(\tslow{@x}{\dfun}^{-1}(s), @x, \dfun)$ satisfy the same
differential equation
\rref{slowsystem} with initial state $@x$ on $[0,
\;\tslow{@x}{\dfun}(\tmax(@x, \dfun)))$, $\therefore$ they coincide on
$[0, \;\tslow{@x}{\dfun}(\tmax(@x, \dfun)))$. 

To show the limit property of $\tslow{@x}{\dfun}$, suppose
\[
\lim_{t->\tmax(@x, \dfun)}\tslow{@x}{\dfun}(t) = b < \infty
\]
(note that the limit exists because $\tslow{@x}{\dfun}(!.)$ is
increasing).  Let
\[
K=\{z(s, @x, \dfun\circ\tslow{@x}{\dfun}^{-1}):  0\le s < b\}.
\]
$\Therefore$ $K$ is bounded, so $\bar K$ is compact,
 and by~\rref{e-xz}, $x(t, @x, \dfun)!eK !c\bar K$ $\forall$
$t!e[0, \tmax(@x, \dfun))$, contradicting the maximality of 
$\tmax(@x, \dfun)$.\mybox

$\forall$ initial state $@x$ and each disturbance \function $\dfun$, define
\be{e-ioslow}
\ioslow_{@x, \dfun} = \inf \left\{t \ge 0
\;\colon\; z(t, @x, \dfun) !e\D\right\}, 
\ee
where $\ioslow_{@x, \dfun} = \infty$ if $z(t, @x, \dfun)!ne\D$ $\forall$
$t\ge 0$. Note that $\ioslow_{@x, \dfun} > 0$ $\forall$ $\dfun !e\M_@W$,
$@x!e\E$ because $\E$ is open. 

Observe also:  if $\dfun_1 = \dfun!o\tslow{@x}{\dfun}$,
\[
\ioslow_{@x, \dfun} = \tslow{@x}{\dfun_1}(\io{@x, \dfun_1}).
\]
%

%\noindent
{\em Claim $2$}. System $\hat @S$ satisfies the \newprop\ property.

{\em Proof}. 
According to~\rref{e-gasmob} and~\rref{e-xz}, we have, for every
$@x\in\E$ and each $\dfun$,
\beqn
\abs{z(t, @x, \dfun)}& =& \abs{x(\tslow{@x}{\dfun_1}^{-1}(t), @x,
  \dfun_1)}\\
&\le& \klfun(\abs{@x}, \tslow{@x}{\dfun_1}^{-1}(t)) \le \hbound(\abs{@x}),
\eeqn
$\forall$ $t!e[0, \ioslow_{@x, \dfun})$, where $\hbound(s) = \klfun(s, 0)$,
and   
$\dfun_1$ is s.t. $\dfun = \dfun_1\circ\!\tslow{@x}{\dfun_1}^{-1}$.
Let 
\[
M_r = 1 + \max_{d!e@W, |@x| !<\hbound(r)} 
\abs{f(@x, d)}^2 + \max_{|@x| !<\hbound(r)} @k(@x).
\]
$\Therefore$ $\forall$ $@x\in\E$ with $\abs{@x}\le r$, it holds that
\[
\tslow{@x}{\dfun_1}(t) = \int_0^t\left[
1+ \abs{f(x(s, @x, \dfun_1), \dfun_1(s))}^2 + @k(x(s, @x, \dfun_1))
\right]\,ds \psp \le \psp M_rt
\]
$\forall$ $t!e[0, \io{@x, \dfun_1})$,
and hence, $\tslow{@x}{\dfun_1}^{-1}(t) \ge \frac{t}{M_r}$ $\forall$
$\abs{@x}\le r$, $t!e[0, \ioslow_{@x, \dfun})$.
Consequently, we have
\be{e-slowgasmo}
\abs{z(t, @x, \dfun)}\le \newklfun(\abs{@x}, t)\qquad \forall\,
t!e[0, \ioslow_{@x, \dfun}),
\ee
where $\newklfun(s, t) = \klfun(s,
\frac{t}{M_s})$ is clearly $!e\kl$. 

this shows system $\hat@S$ is GASMO



From now on, let  $\klfun$ $!e\kl$ be as in
Definition of \newprop\ for the system $\hat @S$, that is, the
following estimate holds for system~\rref{slowsystem}:
\be{e-zgasmo}
\abs{z(t, @x, \dfun)} \le \klfun(\abs{@x}, t) \qquad\forall\,t\in[0,
@q_{@x, \dfun}), 
\ee
$\forall$ $@x\in\E$, $\dfun \in\M_@W$.

According to Proposition 7 in:

{\sc E. D. Sontag}, {\em Comments on integral
  variants of input-to-state stability}, Systems Control Lett., 34
  (1998), pp. 93--100.

$\exists$ $\ki$-functions $@m_1$ and $@m_2$ s.t. 
\be{mu1mu2} 
\klfun(r, t) !<@m_1( @m_2(r) e^{-t}) \qquad \forall \, r,\ t !>0.
\ee
Define 
\be{precomp_def} 
\precomp(s) := @m_1^{-1} (s). 
\ee 

The proof will now develop as follows. We first construct a 
continuous Lyapunov-like
\function $V_0$, defined on the set $\E_1$. Next $V_0$ is approximated by a
 Lipschitz continuous function (by the methods of nonsmooth analysis).
 The resulting function is then approximated by a smooth \function $V_1$.
 Finally  we extend $V_1$ to the rest of the state space, obtaining a 
 Lyapunov-like function, smooth away from the origin, which is then 
 approximated by a smooth Lyapunov function.

\newpage
\subsection{definitions and basic facts on relaxed controls}

Recall that our disturbances $\dfun$ are measurable \functions 
$\R_{!>0} \to  @W$, a compact, convex subset of $\R^m$. 
 
Let  $\rpm\left(@W\right)$ be the set of all Radon probability
measures on $@W$. Bishop's theorem furnishes a weak norm on 
$\rpm\left(@W\right)$, whose corresponding metric topology coincides
 with the weak star topology on 
$\rpm\left(@W\right)$ (see e.g. Warga). 

$\forall$ $T > 0$, we define
$\relaxed_T$ to be the set of
all measurable functions from $[0, T]$ to $\rpm\left(@W\right)$, and
$\relaxed$ to be the set of all measurable functions from 
$\R_{!>0}$ to  $\rpm\left(@W\right)$. 
We topologize $\relaxed_T$  by  weak convergence:
$\cbrace{@n_k(!.)} \to @n(!.)$ in $\relaxed_T$ if and only if  
\[
\int_0^T \int_@Wg(t, @w)d[@n_k(t)](@w)\,dt \to 
 \int_0^T \int_@Wg(t, @w)d[@n(t)](@w)\,dt 
\]
$\forall$ \functions $g: [0, T] !p@W\to \R$ which are continuous in
$@w$, measurable in $t$, and s.t. 
\[
\max\cbrace{\abs {g(t, @w)}, \; @w!e@W}
\]
is integrable 
on $[0, T]$.  We say that $\{@n_k\}->@n$ weakly in
$\relaxed$ if, for every $T >0$, the sequence 
$\cbrace{@n_k \vert_{[0, T]}}$ of restrictions of $@n_k$ to $[0, T]$ 
converges to $@n\vert_{[0, T]}$ in $\relaxed_T$.  

Notice that, since every element of $@W$ can be identified with the
 $@d$-measure, concentrated in it, $@W$ can be embedded into 
$\rpm\left(@W\right)$, $\M_@W$ into $\relaxed$, and $\M_@W^T$ into
$\relaxed_T$ in the obvious way, where $\M_@W^T$ is the set of
functions in $\M_@W$ restricted to $[0, T]$.

$\forall$ $@n!e\rpm\left(@W\right)$,
we denote
\be{rel_f_def}
f(x, @n) \;=\; \int_{@W}f(x, r) \, d@n(r) \,.
\ee
Notice that $\forall$ relaxed control $@n(!.)$, the function 
$f(!., @n(!.)): \, (x, t) \to f(x, @n(t))$ 
is Lipschitz in $x$ and measurable in $t$. Moreover, as $\forall$ 
$x !e\states$, $@n!e\rpm(@W)$, we have a bound
\[
\abs {f(x, @n)} !<\max_{d !e@W} \abs{f(x, d)},
\]
the solution of the ``system'' $@.x(t) = f(x(t), @n(t))$ exists
$\forall$ initial condition $@x$ and relaxed control $@n$ on some
maximal interval $[0, \tmax(@x, @n))$. Write $x(!., @x, @n)$ to
denote this solution. Just as in the case with ordinary controls, 
we will define 
\be{lambda_xi_nu}
\io{@x, @n} := 
\inf \cbrace{t !<\tmax \; : \; x(s, @x, @n) !e\D}.   
\ee

The basic three facts we will be using in the next section are as follows.

{\em Fact $1$}. $\forall$ $T > 0$, the space $\relaxed_T$ is sequentially
compact (see e.g. Warga).  Consequently,
$\relaxed$ is sequentially compact by a diagonalization argument. 

{\em Fact $2$}. $\forall$ $T > 0$,  the set $\M_@W^T$ of ordinary controls
on $[0, T]$  is dense in $\relaxed_T$
%(see \cite[p.\ 691]{Artstein}, also \cite{Warga}).  

Consequently,
$\M_@W$ is dense in $\relaxed$. The topology of $\relaxed$ induces a 
topology on the subspace $\M_@W$. This topology is stronger than the 
topology of $L^p$ $\forall$ positive $p$: in fact, 
$\cbrace{\dfun_k(t)} \to \dfun$
 would imply
\[
\int_0^T g(\dfun_k(t)) dt \to \int_0^T g(\dfun(t))\,dt
\]
$\forall$ continuous \function $g: \controls \to \R$ and any positive $T$.

{\em Fact $3$}. 
$\forall$ $T > 0$,
the mapping $(t, @x, @n(!.)) |->z(t, @x, @n)$ is continuous
on $[0, T]!p\states!p\relaxed_T$ 
%(see \cite[Lemma 3.12]{ALS}).
%
 
Some relevant immediate consequences of Facts 1, 2, and 3
are as follows.
 

{\em Claim $1$}. The function $(@x, @n) \to \io{@x, @n}$, mapping 
initial states and disturbances to the first hitting times as defined 
in~\rref{lambda_xi_nu}, is
 lower semicontinuous on both $@n$ and $@x$.
% 
(This easily follows from continuous dependence on initial
conditions and control, and from the fact that the set $\D$ is
closed.)

 
{\em Claim $2$}.  If system~\rref{dsystem} is \newprop, $\therefore$ $\forall$ initial 
value $@x!e\E$, a  relaxed disturbance $@n\in\relaxed$, and time 
$T !<\io{@x, @n}$  
  we have an estimate
%
\[
 x(T, @x, @n) !<\klfun(|@x|, T).
\]
\bpr
Pick $@x!e\states$ and $@n!e\relaxed$.
Let $\cbrace{\dfun_k}$ be a sequence of ordinary disturbances, 
converging to $@n$ in $\relaxed$. $\Therefore$ for large enough $k$ we have
$\tmax(@x, \dfun_k) > T$, and $x(!., @x, \dfun_k)$ converge  to 
$x(!., @x, @n)$ uniformly on $[0, T]$. Also, by Claim 1, 
$\liminf_{k \to \infty} \io{@x, \dfun_k} !>\io{@x, @n}$. 
Since we assume the system to be 
\mbox{\newprop,}
the  estimate

$\abs{x(t, @x, \dfun)} !<\klfun (\abs{@x}, t) \quad \forall\; 0 !<t
!<T$

holds for $\dfun:= \dfun_k$ $\forall$ large enough $k$ and $\forall$ $t !<T$, so 
the claim follows.
\epr  

With the previous claim in mind, we can say that the system~\rref{dsystem} 
with $d !e\relaxed$ is \newprop\   
(this is a slight abuse of terminology because, strictly speaking,
\rref{dsystem} with relaxed disturbances is not really a ``system,''
as that would mean, by definition, that disturbances take values in a
finite dimensional space).
Consequently, the auxiliary 
system~\rref{slowsystem} is also \newprop\ for $\dfun\in\relaxed$.  
Thus, we can assume that~\rref{e-zgasmo} holds $\forall$ $@x\in\E$ and
all $\dfun\in\relaxed$. 

\newpage
\subsection{constructing a continuous Lyapunov-like function on {\protect\boldmath{$\E_1$}}} \label{dV0}
\renewcommand{\ioslow}[1]{{@q_{#1}}} 

we introduced the system $\hat @S: $
$@.z = \hat{f}(z, d):=\frac{f(z, d)}{1+\abs{f(z, d)}^2 +
@k(z)}$, which slows down the motions of the original system $@S$.
%
Recall also that $\D:=\left\{@x: \abs@x\leq
@r(\abs{h(@x)})\right\}$, and define the set 
\[
\B\,:=\;\{@x: @r(\abs{h(@x)}) !<\abs @x!<1.5@r(\abs{h(@x)})\}
\,.
\]

Let $f_0:\states->\R$ be defined by 
\[
f_0(@x) = \max_{\rd\in@W}\abs{\newf(@x, \rd)}.
\]
Note that $f_0$ is locally Lipschitz, and recall that we have assumed 
with no loss of generality that $@S$ has no equilibria on the set 
$\cbrace{x !e\states \;:\;\abs x !>@r(\abs{h(x)})}$ 
(otherwise replace $@r(!.)$ by $c@r(!.)$, where $c > 1$). 
In particular,  $f_0(@x)\not=0$ $\forall$
$@x\in\partial\D\setminus\{0\}$. 
Let $\mix: \states\setminus\{0\} \to [0, 1]$ be smooth and
\[
\mix(x) \;=\; \twoif{1,}{x !e\D,}{0,}{x !e\states \setminus (\D \cup \B).}
\]
Now introduce another system, on the state space 
$\states \setminus \cbrace{0}$,  
\[
\tilde @S: \;\;\; @.z = \tf(z, d, v),
\] where
disturbances $d$ are as before, and auxiliary controls $\vc$ are
measurable functions of time,  taking values in
$[-1, 1] ^n$
(note that the dimension of the control set for $v$'s is the same as
that of $\states$), where, $\forall$ $i =1, 2, !..., n$,
\[
\tf_i(z, d, v) := \newf_i(z, d) + 2\mix(z)f_0(z)v_i.
\]

$\forall$ $T > 0$, let $\pcow_T$ denote the set of the auxiliary controls
defined on $[0, T]$
(i.e., measurable functions $\vc: [0, T] \to [-1, 1]^n$) 
 equipped with the weak convergence topology; that is, ``$\{\vc_k\}$
converges weakly to $\vc$ in $\pcow_T$'' means
\[
\int_0^T\vf(s)\vc_k(s)\,ds\;->\int_0^T\vf(s)\vc(s)\,ds
\]
$\forall$ \functions $\vf$ that are integrable over $[0, T]$.  With the
weak topology, $\pcow_T$ is sequentially compact 
%(cf.~\cite[Proposition 10.1.5]{mct}).  

Consequently, given any sequence $\{\vc_k\}$ of 
controls defined on $[0, \infty)$, $\exists$ control $\vc$ and a
subsequence $\{\vc_{k_j}\}$ s.t. $\vc_{k_j}->\vc$ weakly on every
interval $[0, T]$.  

We denote the set of the auxiliary controls defined on $[0, \infty)$ by
$\pcow$.  Let $\{\vc_k\}\subset\pcow$ and $\vc !e\pcow$.  We say that
$\{\vc_k\}$ weakly converges to $\vc$ if, for every $T > 0$, the sequence 
of restrictions $\{\vc_k \vert_{[0, T]}\}$
weakly converges to $\vc \vert_{[0, T]}$  in $\pcow_T$.

Recall that $z(t, @x, \dfun)$ denotes the solution of $\hat @S$.
Write $z(t, @x, \dfun, \vc)$ for the value at time $t$ of the solution of 
$\tilde@S$ with initial state $@x\not=0$, 
disturbance $\dfun !e\M_@W$, and
auxiliary control $\vc !e\pcow$.

Observe the following.
\begin{itemize}
\item $\tilde @S$ is affine in $v$.

\item since $\forall$ $@x!e\states$ and $d !e[-1, 1]^m$,
$|{\hat f(@x, d)}| !<1$, we have $|{\tf(@x, d, v)}| !<3$
$\forall$ $@x$, $d$ and $v$. In particular, this implies that
$\tilde@S$ is forward complete.

\item Suppose $@x\not!e\D\cup\B$ and pick $\dfun !e\M_@W$ 
and $\vc !e\pcow$.
 $\Therefore$ there is some $t_0 > 0$ s.t.   
$z(t, @x, \dfun) \not!e\D\cup\B$ $\forall$ $t !e [0, t_0]$, and
$z(t, @x, \dfun, \vc) ==z(t,@x, \dfun)$ on $[0, t_0]$.
\end{itemize}

To extend the definition of $z(t, @x, \dfun, \vc)$ to the case when 
$\dfun !e\relaxed$, we let, for a fixed $\dfun !e\relaxed$ and 
$\vc !e\pcow$, 
\[
g_{\dfun, \vc}(z, t):= \hat f(z, \dfun(t)) + 2\mix(z)f_0(z) \vc(t)
\]
(where $f(z, \dfun(t))$ is as defined in~\rref{rel_f_def} for 
$@n:= \dfun(t) !e\rpm(@W)$). $\Therefore$  
$g_{\dfun, \vc} : 
\states \setminus \cbrace{0} !p\R_{!>0} \to \states $ 
is locally Lipschitz in its first variable, and
$\abs{g_{\dfun, \vc} (z, t)} !<3$ $\forall$ 
$(z, t) !e\states \setminus \cbrace{0} !p\R_{!>0}$.
Hence, the solution of 
\begin{eqnarray*}
@.z(t) &=& g_{\dfun, \vc}(z, t), \\
z(0) &=& @x
\end{eqnarray*}
exists $\forall$ $@x!e\states \setminus \cbrace{0}$ and $t!>0$. We 
will denote it by $z(t, @x, \dfun, \vc)$.

Observe that $\tilde f: \states \setminus\cbrace{0}!p@W!p[-1,1]^n
->\R^n$ is continuous, and $\tilde f(z, d, v)$ is locally Lipschitz
in  $z$ on $\states\setminus\{0\}$
uniformly on $(d, v)!e@W!p[-1, 1]^{n}$.  
The system $\tilde @S$ evolves in the state space 
$\states \setminus \cbrace{0}$. As $|{\tilde f}| !< 3$ everywhere, 
trajectories are defined and unique $\forall$ initial value 
$@x!e\states \setminus \cbrace{0}$
and each pair of inputs $\dfun, \vc$.
 Moreover, if $z(!.)$ is a maximal such 
trajectory, $\therefore$ either $z(t)$ is defined $\forall$ $t !>0$, or there 
is some $T > 0$ s.t. $\lim_{t \to T} z(t) = 0$. We prove next that 
this last case cannot happen.

\bll{z>w:lemma}
For every ball $U$ around $0$, $\exists$ constant
 $c$ s.t. $\forall$ 
$@x!eU$, $@x\neq 0$, $\dfun !e\relaxed$, 
$\vc !e\pcow$, and $t !>0$, we have
 a lower bound
\be{z>w}
\abs{z(t, @x, \dfun, \vc)} > \frac{1}{2}\abs{@x} \,e^{-ct}.
\ee
\ells

\bpr 
Since $\hat f$ is locally Lipschitz in $x$ uniformly in $d$, 
and 
$\hat f(0, d) = 0$ 
$\forall$ $d$,
we can find a positive constant $c$,
s.t. $|{\hat f(z, d)}| !<c|z|/3$ $\forall$ $z !eU$,
$d !e@W$

$\Therefore$ also 
$|f_0(z)| !<c|z| /3$; $\therefore$
\[
\abs {\tilde f(z, d, v)} !<\abs{\hat f(z, d)} + 2\abs{f_0(z)}
!<c|z|
\] 
$\forall$ $v !e[-1, 1]^{n}$, all $d !e@W$, and hence, all 
$d !e\rpm(@W)$.
Now fix $@x!eU \setminus \cbrace 0$, $\dfun !e\relaxed$,
 and $\vc !e\pcow$,
and write $z(t):=z(t, @x, \dfun, \vc).$
Note that the inequality~\rref{z>w} holds for $t=0$; $\therefore$ it holds 
$\forall$ small enough $t>0$. Suppose that~\rref{z>w} fails at some $t_2 >0$,
s.t. 
\be{fails_at_t2}
\abs{z(t_2)} 
        !<\frac{1}{2}\abs{@x} \,e^{-ct_2} < \abs {@x}.
\ee
$\Therefore$ $\exists$  $t_1 < t_2$ s.t. 
$\abs{z(t_1)}= \abs @x$ and  
$\abs{z(t)} !<\abs{@x}$ $\forall$ $ t !e[t_1, t_2]$.
Let 
\[
w(t):= \abs{z(t)}^2/2.
\]
$\Therefore$ for almost all $t !e[t_1, t_2]$ 
\[
\abs{@.w(t)} = \abs {z(t) !.@.z(t)} !< c \abs{z(t)}^2 = 
2c w(t).
\]
In particular, this implies that $@.w(t) + 2cw(t) !>0$. So,
$\forall$ $t !e[t_1, t_2]$ we have
\[
0 !<e^{2ct}(@.w(t) +2cw(t)) = \frac{d(e^{2ct}w(t))}{dt}, 
\]
implying that $e^{2ct}w(t) !>e^{2ct_1}w(t_1)$ $\forall$ $t !e[t_1, t_2]$.
 Thus, 
\[
\frac{1}{2}\abs{z(t_2)}^2=w(t_2) !> e^{2ct_1}w(t_1) e^{-2ct_2}
 = \frac{1}{2}\abs{z(t_1)}^2 e^{-2c(t_2-t_1)}
!> \frac{1}{2}\abs{@x}^2 e^{-2ct_2},
\]
s.t. $\abs{z(t_2)} !> e^{-ct_2} \abs{@x}$, 
contradicting~\rref{fails_at_t2}.
\epr

\bc{l-reach0}
For every $r > 0$ and 
$T > 0$ $\exists$ $@s= @s(r, T)  > 0$ 
s.t. $\forall$ $\dfun !e\relaxed$, $\vc !e\pcow$, 
$\abs {@x} !>r$, and $t !<T$ we have
%
\[
 \abs {z(t, @x, \dfun, \vc)} !>@s.
\]
\ec

\bll{uni_con}
$\forall$ $\dfun\in\relaxed_T$ and each $\vc\in\pcow_T$, if
$@x_k->@x$ in $\states \setminus \cbrace{0}$, 
$\dfun_k ->\dfun$ in $\relaxed_T$, and $\vc_k ->\vc$ in $\pcow_T$, 
$\{z(t, @x_k, \dfun_k, \vc_k)\}$ converges to $z(t, @x, \dfun, \vc)$ uniformly 
on $[0, T]$.
\ells

\bpr
Assume w.l.o.g. that 
$@x_k !eU := B_{\frac{\abs @x}{2}}(@x)$. Since 
$|{\tilde f} !<3|$ and by Corollary~\ref{l-reach0}, 
$!R_T\rb{U} !cB_{1.5 \abs{@x} + 3T}(0) \setminus B_@s(0)$,
where $@s= @s(\abs{@x}/2, T)$ is as in Corollary~\ref{l-reach0}.
Let $M_1$ and $M_2$ be Lipschitz constants for 
$\hat f(!., d)$ (uniformly for $d !e@W$)  and $2 f_0 \, \mix$,
 respectively, on $B_{1.5 \abs{@x} + 3T}(0) \setminus B_@s(0)$. 
Write $z(t) := z(t, @x, \dfun, \vc)$, and $z_k(t):= z(t, @x_k, \dfun_k, \vc_k)$.
Now, $\forall$ $t !e[0, T]$
we have
%
\begin{eqnarray}
\abs{z(t) -z_k(t)} & = & 
\abs{@x_k -@x+ 
\int_0^t\rb{\tilde f(z_k(t), \dfun_k(t), \vc_k(t)) - 
            \tilde f(z(t), \dfun(t), \vc(t))}dt} \nonumber\\
& !<& \abs{@x_k -@x} +
\abs{\int_0^t\rb{\tilde f(z(t), \dfun_k(t), \vc_k(t)) - 
            \tilde f(z(t), \dfun(t), \vc(t))}dt} \nonumber \\
&& + \int_0^t\abs{\tilde f(z_k(t), \dfun_k(t), \vc_k(t)) - 
            \tilde f(z(t), \dfun_k(t), \vc_k(t))}dt \nonumber \\
& !<& \abs{@x_k -@x} +
\abs{\int_0^t \rb{\hat f(z(t), \dfun_k(t)) - 
            \hat f(z(t), \dfun(t))}dt} \label{uni_con:1} \\
&& + 
2\abs{\int_0^t \rb{ f_0(z(t))\mix(z(t))\vc_k(t) 
  - f_0(z(t)) \mix(z(t))\vc(t)}dt} \label{uni_con:2} \\
&&  + \int_0^t(M_1 +M_2)\abs{z_k(t) -z(t)}dt. \nonumber
\end{eqnarray} 
%
The integrals in~\rref{uni_con:1} and~\rref{uni_con:2} tend to $0$ because
of the convergence of $\cbrace {\dfun_k(!.)}$ to $\dfun$ in $\relaxed$ 
and the weak convergence of $\vc_k$ to $\vc$. So, $\forall$ $\ve >0$ we can find 
a $K$ s.t., $\forall$ $k !>K$, 
\begin{eqnarray*}
&&\abs{@x_k -@x} +
\abs{\int_0^T \rb{\hat f(z(t), \dfun_k(t)) - \hat f(z(t), \dfun(t))}dt} \\
&+& 2\abs{\int_0^t \rb{ f_0(z(t))\mix(z(t))\vc_k(t) 
  - f_0(z(t)) \mix(z(t))\vc(t)}dt} !<\ve e^{-(M_1 + M_2)T}.   
\end{eqnarray*}
$\Therefore$ $\forall$ $k !>K$ and $t !e[0, T]$ we have, by the Gronwall inequality,
\[
\abs{z(t) -z_k(t)} !<\ve e^{-(M_1 + M_2)T} \,e^{(M_1 + M_2)t} !<\ve.
\]
As $\ve$ was arbitrary, this proves uniform convergence.
\epr

Also, since $\forall$ $@x\in!dD\setminus\{0\}$, $f_0(@x)\not=0$,
and $f_0(@x) \ge |{\hat f(@x,d)}|$ $\forall$ $d\in@W$, we have the
following controllability property on $\partial\D\setminus\{0\}$.
\bll{l-contr}
Let $@x\not=0$ be on $\partial\D$.  $\Therefore$ $\forall$ $@t>0$, $\exists$
a neighborhood $U$ of $@x$, s.t. $\forall$ $@h!eU$ and any
$\dfun\in\M_@W$, there is some control $\vc$, and some $0\le t_1\le @t$, 
s.t. $z(t_1, @h, \dfun, \vc) = @x$  and $z(t, @h, \dfun, \vc)\in
U$ $\forall$ $0\le t\le t_1$.
\ells

\bpr
Since $\mix(@x) \, f_0(@x) \neq 0$ and the function 
$\mix(!.) \, f_0(!.)$ 
is continuous, we can find a ball
$U_1$ centered at $@x$ and a constant $c_1$ s.t. $\mix(z) f_0(z) > c_1$
for every $z !eU_1$. Since $\mix (@x) = 1$ and $\mix$ is continuous, 
one could also find a ball $U_2 !cU_1$ centered at $@x$, s.t.
%
\[
\abs{\newf(z, d)} !<1.5 \mix(z) f_0(z) \;\;\forall z !eU_2, \, 
d !e@W.
\]
%
Fix $@t> 0$ and let $B(@x)$ be the ball of radius $@tc_1/2$ centered
 at $@x$. Define $U:= B(@x) \cap U_2$. Pick a point $@h!eU$.
$\Therefore$ $\abs{@x-@h} < @tc_1/2$, so, 
\[
\bar v_2 := \frac {2 (@x-@h)}{@tc_1}
\]
has norm smaller than $1$.
Consider the ``feedback law'' 
\[
k(z, d) = \frac{1}{2}(1.5 k_1(z, d) + 0.5\bar v_2),
\] where 
\[
k_1(z ,d) := -\frac{\newf(z, d)}{1.5 \mix(z) f_0(z)}.
\]
Notice that $\forall$ $z !eU$ and $d !e@W$ we have
 $\abs {k_1(z, d)} !<1$ and 
\begin{eqnarray*}
\tilde f(z, d, k(z, d)) 
        & = & \newf(z, d) + 2 \mix(z) f_0(z) k(z, d) \\
        & = & \newf(z, d) + 1.5 \mix(z) f_0(z) k_1(z,d) 
                        + 0.5 \mix(z) f_0(z) \bar v_2  \\
        & = & 0.5 \mix(z) f_0(z) \bar v_2 = 
                \frac{(@x- @h)\mix(z) f_0(z)}{@tc_1}.
\end{eqnarray*}

Thus, with any initial condition $@h!eU$ and disturbance 
$\dfun !e\M_@W$, if the control $v(t):= k(z(t), d(t))$ is applied, 
$\therefore$ the trajectory of the system $\tilde @S$ will be the 
line segment, connecting $@x$ and $@h$, transversed with a velocity 
greater than $(@x- @h)/@t$. So, $\exists$  $t_0 !<@t$
s.t. $z(t_0, @h, \dfun, \vc) = @x$ and, since $U$ 
is convex, $z(t, @h, \dfun, \vc) !eU$ $\forall$ $t !>t_0$.  
\epr  
%

$\forall$ $\dfun !e\relaxed$, let
 \[
\ioslow{\dfun}(@x, \vc)=\inf\left\{t\ge 0: z(t, @x, \dfun, \vc) !e\D
\right\},
\] 
where as before, $\ioslow{\dfun}(@x, \vc) = \infty$ if
the trajectory never reaches $\D$.

\bll{l-ioslowsc}
The map $(@x, \vc, \dfun) |->\ioslow{\dfun}(!., !.)$ is lower
semicontinuous on $\E!p\pcow!p\relaxed$.
\ells

\bpr Let $\{@x_k\}\subset\E$, $\{\vc_k\}\subset \pcow$ and 
$\{\dfun_k\}\subset\relaxed$
be s.t. $@x_k->@x, \vc_k->\vc$, and $\dfun_k->
\dfun$ for some $@x\in\E$, $\vc\in\pcow$, and $\dfun\in\relaxed$.
We need to show that
\be{e-someeq}
\ioslow{\dfun}(@x, \vc) \le \liminf_{k->\infty}\ioslow{\dfun_k}(@x_k, \vc_k).
\ee
Let $\ioslow{k} = \ioslow{\dfun_k}(@x_k, \vc_k)$.  Without
loss of generality, we may assume that
\[
\liminf_{k->\infty}\ioslow{k} = \ioslow{0} < \infty.
\]
Passing to a subsequence if necessary, we assume that $\ioslow{k}
->\ioslow{0}$. Thus, $\exists$ $K$ s.t. 
$\ioslow{k}\le\ioslow{0} + 1$ $\forall$ $k\ge K$.  Since $\{z(t, @x_k, \dfun_k,
\vc_k)\}$ converges to $z(t, @x, \dfun, \vc)$ uniformly on $[0,
\ioslow{0}+1]$, it follows that 
\[
z(\ioslow{0}, @x, \dfun, \vc) = \lim_{k->\infty}
z(\ioslow{k}, @x_k, \dfun_k, \vc_k).
\]
Since $\D$ is closed and $z(\ioslow{k}, @x_k, \dfun_k, \vc_k)!e\D$ $\forall$
$k$, we know that $z(\ioslow{0}, @x, \dfun, \vc)!e\D$, and hence,
$\ioslow{\dfun}(@x, \vc)\le \ioslow{0}$.
\epr

Define, for $@x\in\E, \dfun\in\relaxed$, and $\vc\in\pcow$,
%
\[
V_\vc(@x, \dfun):= \int_0^{\ioslow{\dfun}(@x, \vc)}\precomp(\abs{z(t, @x, \dfun,
  \vc)})\,dt
\]
(where $\precomp$ is as in~\rref{precomp_def}) and, for $@x\in\E$ 
and $\dfun\in\relaxed$,
\[
\tilde V_0(@x, \dfun):= \inf_{\vc !e\pcow} V_\vc(@x, \dfun).
\]
Note that for some $\vc$ and $\dfun$, $V_\vc(@x,\dfun)$ may take $\infty$ as its
value, but $\tilde V_0(@x, \dfun)$ is always finite, since
\be{V_0finite1}
\tilde V_0(@x, \dfun)\le V_O(@x, \dfun),
\ee
 where $O(!.)$
is the control identically equal to $0$. 

Recall that by definition of $\precomp$ we have
$\therefore$ $\forall$ $@x!e\E_1$ and $\dfun !e\relaxed$,  
\begin{eqnarray}
&& V_O(@x, \dfun) = \int_0^{\ioslow{\dfun}(@x, O)}\precomp(\abs{z(t, @x, \dfun,
O)})\,dt \nonumber \\
 &=& \int_0^{\ioslow{\dfun}(@x, O)}\precomp(\abs{z(t, @x, \dfun)})\,dt \\
&!<& 
\int_0^{\infty} \precomp (@m_1(@m_2(\abs @x) e^{-t}))dt !<@m_2(\abs @x), \label{V_0finite2}
\end{eqnarray}
where $@m_1$ and $@m_2$ are as in \rref{mu1mu2}.

\bll{l-lsc} The function $V_{(!.)}(!., !.):\
\E!p\relaxed!p\pcow \to \R_{!>0}$ is lower semicontinuous.
\ells

\bpr
Let $(@x, \dfun, \vc)\in\E!p\relaxed!p\pcow$,
and let $\{@x_k\}->@x$, $\{\dfun_k\}\to
\dfun$ and $\{\vc_k\}\to \vc$, where $@x_k\in\E$ $\forall$ $k$.  

{\it Case $1$}. \ $V_\vc(@x, \dfun) < \infty$.  In this case, $\forall$
$\ve>0$, $\exists$ $0< T <\ioslow{\dfun}(@x, \vc)$ s.t.
\[
V_\vc(@x, \dfun) \psp= \psp \int_0^{\ioslow{\dfun}(@x,\vc)}\precomp(\abs{z(t, @x, \dfun,
  \vc)})\,dt
\psp \le \psp  \int_0^T\precomp(\abs{z(t, @x, \dfun,  \vc)})\,dt + \ve.
\]
w.l.o.g. we can assume that all $@x_k$ are within the unit 
distance from $@x$. Recall that the reachable set from the unit ball 
around $@x$ is bounded.
Since $z(t, @x_k, \dfun_k, \vc_k)$ converges to $z(t, @x, \dfun, \vc)$
uniformly on $[0, T]$, and 
$\precomp(!.)$ is uniformly continuous
on compacts, $\exists$ $K > 0$ s.t.
\[
\abs{\precomp(\abs{z(t, @x, \dfun,  \vc)}) - 
\precomp(\abs{z(t, @x_k, \dfun_k,  \vc_k)})} < \frac{\ve}{1+T} \;\;
\forall k > K, \; \forall t !e[0, T].
\]
This implies that
\[
\int_0^T\precomp(\abs{z(t, @x_k, \dfun_k,  \vc_k)})\,dt \psp \ge \psp
\int_0^T\precomp(\abs{z(t, @x, \dfun,  \vc)})\,dt - \ve
\]
$\forall$ $k\ge K$.  By Lemma~\ref{l-ioslowsc}, $\exists$ $K_1\ge
  K$ s.t. $\ioslow{\dfun_k}(@x_k, \vc_k) > T$ $\forall$ $k\ge K_1$.  Thus,
  $\forall$ $k\ge K_1$,
\beqn
V_{\vc_k}(@x_k, \dfun_k) &\ge&
\int_0^T\precomp(\abs{z(t, @x_k, \dfun_k,  \vc_k)})\,dt\\
&\ge&
\int_0^T\precomp(\abs{z(t, @x, \dfun,  \vc)})\,dt - \ve
 \psp \ge \psp
V_\vc(@x, \dfun) - 2\ve.
\eeqn
As $\ve$ was arbitrary, we conclude that 
\[
V_\vc(@x, \dfun) \le \liminf V_{\vc_k}(@x_k, \dfun_k).
\]


{\it Case $2$}. \ $V_\vc(@x, \dfun) = \infty$.

In this case, $\ioslow{\dfun}(@x, \vc) = \infty$.  Fix an integer
$k\ge 0$. $\exists$ $T_k$ s.t.
\[
\int_0^{T_k}\precomp(\abs{z(t, @x, \dfun,  \vc)})\,dt\ge k.
\]
Repeating the same argument used above, one sees that 
\[
\int_0^{T_k}\precomp(\abs{z(t, @x_l, \dfun_l,  \vc_l)})\,dt\ge k -1
\]
$\forall$ $l\ge L_0$ for some $L_0$.  By Lemma~\ref{l-ioslowsc}, $\exists$
$L_1\ge L_0$ s.t. $\ioslow{\dfun_l}(@x_l, \vc_l)\ge T_k$ $\forall$
$l\ge L_1$.  Consequently, $\forall$ $l\ge L_1$,
\[
V_{\vc_l}(@x_l, \dfun_l) \ge \int_0^{T_k}\precomp(\abs{z(t, @x_l, \dfun_l,
  \vc_l)})\,dt\ge k -1. 
\]
Since $k > 0$ can be picked arbitrarily, it follows that
\[
\liminf\,V_{\vc_l}(@x_l, \dfun_l) = \infty.
\]
In both cases, we have shown that $\liminf\,V_{\vc_l}(@x_l, \dfun_l) \ge
V_\vc(@x, \dfun)$.  The lower semicontinuity property follows readily.
\epr

\bll{Opt_v_exists}
For every $@x!e\E$, $\dfun !e\relaxed$, there 
exists a control $\bar \vc$ s.t.
\be{e-opt-v_exists}
V_{\bar \vc}(@x, \dfun) = \tilde V_0(@x, \dfun).
\ee
\ells

\bpr
Let $\vc_k$ be a sequence of controls s.t. 
$V_{\vc_k}(@x, \dfun)\searrow \tilde V_0(@x, \dfun)$

w.l.o.g. assuming all these controls are defined $\forall$ positive $t$ (by
letting them equal to $0$ where they are not defined)

Extract from $\{\vc_k\}$
a subsequence $\{\vc_{k_l}\}$ converging weakly to some limit $\bar v$ in
$\pcow$.  Without relabeling, we assume that $\vc_k\to \bar \vc$

By Lemma~\ref{l-lsc},
\[
V_{\bar \vc}(@x, \dfun)\le \lim_{k->\infty}V_{\vc_k}(@x, \dfun) = \tilde
V_0(@x, \dfun).
\]
Combining this with the fact that $\tilde V_0(@x, \dfun)\le V_\vc(@x, \dfun)$ $\forall$
$\vc\in\pcow$, one thus proves \rref{e-opt-v_exists}.
\epr

\bc{from_Opt_v_exists}
$\forall$
$@x\in\E, \dfun \in\relaxed, 
\vc\in\pcow$, and $0\le T < \ioslow{\dfun}(@x, \vc)$, it
holds that 
\be{dynprog}
\tilde V_0(@x, \dfun) \; \le \; \int_0^T\precomp(\abs{z(s, @x, \dfun, \vc)})\,ds +
\tilde V_0(z(T, @x, \dfun, \vc), \dfun_T),
\ee
where $\dfun_T(t) = \dfun(t+T)$ $\forall$ $t\ge 0$.
\ecs

\bpr
Suppose the assertion is not true. $\Therefore$ $\exists$ $@x!e\E$, $T >0$, 
 $\vc !e\pcow$, and $\dfun !e\relaxed$ s.t.~\rref{dynprog} fails. 
By Lemma~\ref{Opt_v_exists}, one can find a control $ \vc_1$ s.t.
$\tilde V_0(z(T, @x, \dfun, \vc), \dfun_T) = V_{\vc_1}(z(T, @x, \dfun, \vc), \dfun_T)$.
Define $\bar \vc$ to be the concatenation of $\vc$ and $\vc_1$. 
$\Therefore$ letting $@q:= \ioslow{\dfun_T}(z(T, @x, \dfun, \vc), \vc_1)$ and
noticing that $\ioslow{\dfun}(@x, \bar \vc) = @q+T$, we get, by our
assumption,
\begin{eqnarray*}
\tilde V_0(@x, \dfun)  &>&  \int_0^T\precomp(\abs{z(s, @x, \dfun, \vc)})\,ds +
\tilde V_0(z(T, @x, \dfun, \vc), \dfun_T) \\
&=&  \int_0^T \precomp(\abs{z(s, @x, \dfun, \vc)})\,ds +
 \int_0^@q\precomp(\abs{ z(t, z(T, @x, \dfun, \vc), \dfun_T, \vc_1)}) \,dt \\
&=& \int_0^{@q+T}\precomp(\abs{z(s, @x, \dfun, \bar \vc)})\,ds,
\end{eqnarray*}
which contradicts with the minimality of $\tilde V_0(@x, \dfun)$.
\epr 


\bll{l-VdContD}
$\forall$ $@x\in!d\D$, the following holds:
\be{e-Vcb}
\lim_{{@h->@x}}\tilde V_0(@h, \dfun) = 0
\ee
uniformly in $\dfun\in\relaxed$, that is, $\forall$ $\ve > 0$, $\exists$
neighborhood $U$ of $@x$ s.t. $\tilde V_0(@h, \dfun) < \ve$ $\forall$
$@h!eU\cap \E$, $\dfun\in\relaxed$
\ells

\bpr 
If $@x= 0$, the result follows from~\rref{V_0finite1} and~\rref{V_0finite2}.  Suppose
now that $@x\not=0$.

Let $\ve>0$ be given.  Let $@t=\frac{\ve}{\precomp(\abs{@x}+1)}$.
Find a neighborhood $U$ of $@x$ as in Lemma~\ref{l-contr}.  Shrinking
$U$ if necessary, we assume that $\abs{@x- @h}\le 1$ $\forall$
$@h!eU$.

Suppose $@h!eU\cap\E$ and  $\dfun\in\M_@W$.  By the controllability
property, there is some 
control $v$ s.t. $z(t_1, @h, \dfun, \vc)=@x\in\partial\D$ for some 
$t_1\in[0, @t]$,
and that $z(t, @h, \dfun, \vc)!eU$ $\forall$ $t!e[0, t_1]$.  Thus,
\[
V_{\vc}(@h, \dfun) \le \int_0^{t_1}\precomp(\abs{z(s, @h, \dfun, \vc)})\,ds
\le @t!.\precomp(\abs{@x}+1) \le \ve,
\]
{} from which it follows that $\tilde V_0(@h, \dfun) \le \ve$.  

The above shows that $\tilde V_0(@h, \dfun)\le \ve$ $\forall$ nonrelaxed
 $\dfun\in\M_@W$,
$@h!eU\cap\E$.
Pick a relaxed $\dfun\in\relaxed$. $\Therefore$ $\exists$
a sequence 
$\{\dfun_k\}\subset\M_@W$ s.t. $\dfun_k->\dfun$
in the topology of relaxed controls.  Let $@h\in
U\cap\E$.  Let $\vc_k$ be s.t. $\tilde V_0(@h, \dfun_k) = V_{\vc_k}(@h,
\dfun_k)$.  By weak sequential compactness of $\pcow$, one may assume, 
after taking a subsequence, that $\vc_k ->
\bar \vc$ for some $\bar \vc$.   By Lemma~\ref{l-lsc},
\[
V_{\bar \vc}(@h, \dfun)\le\liminf_{k\to\infty} V_{\vc_k}(@h, \dfun_k)\le\ve,
\]
and consequently, $\tilde V_0(@h, \dfun)\le \ve$.  This shows that
$\tilde V_0(@h, \dfun)\le \ve$ $\forall$ $@h!eU \cap \E$, $d\in\relaxed$
\epr

To prove the continuity of $\tilde V_0$, we also need the following
result.
\bll{l-sl1}
Suppose for some $@x\in\E$, $\dfun\in\relaxed$, and $\vc\in\pcow$,
$V_\vc(@x, \dfun) < \infty$. $\Therefore$ $\exists$ $@x_0!e!d\D$
s.t.
\be{e-sl1-liminf}
\lim_{t->\ioslow{\dfun}(@x, \vc)} z(t,@x,\dfun,\vc) = @x_0. 
\ee
\ells

\bpr
  Suppose $V_\vc(@x, @n) < \infty$.  This means that
\be{e-sl1-conv}
\int_0^{\ioslow{@n}(@x, \vc)}\precomp(\abs{z(s,@x,@n, \vc)})\,ds < \infty.
\ee
If $\ioslow{\dfun}(@x, \vc) < \infty$, $\therefore$~\rref{e-sl1-liminf} follows from
the continuity of $z(!., @x, @n, \vc)$ with $@x_0 = z(\ioslow{\dfun}(@x,
\vc))$.  

Suppose now that
 $\ioslow{\dfun}(@x, \vc) =  \infty$. 
Since the integral in~\rref{e-sl1-conv} converges, 
it follows that
\[
\int_t^!8\precomp(\abs{z(s,@x,\dfun, \vc)})\,ds
->0 \quad \text{as}\ t->\infty.
\]
%decreases to $0$ as $t$ tends to $\infty$. 
Consider the family of functions $\cbrace{x_t(!.), \; t >0}$, defined by
$x_t(s) := z(t+s, @x, \dfun, \vc)$, $\I_t = [0, \infty)$. 
%REV1
By Lemma~\ref{z>w:lemma}, the trajectory $z(s,@x, \dfun, \vc)$
does not reach the origin in finite time, hence, $\exists$  
positive, strictly 
decreasing \function $\vf$ s.t. $\vf(s) < \abs{z(s,@x, \dfun, \vc)}$ $\forall$
 $s > 0$. Find a $@k!e\ki$ s.t. 
\[
@k(\vf(t)) > \int_t^!8\precomp(\abs{z(s,@x,\dfun, \vc)})\,ds=
 \int_0^!8\precomp(\abs {x_t(s)})\, ds.
\]
$\Therefore$ the family $\cbrace{x_t(!.), \; t > 0}$ satisfies all
the conditions of 
Proposition~\ref{barbalat} (with $@c:= \precomp$). Take $r := \abs {@x}$.
$\Therefore$ $\forall$ $\ve > 0$, $\abs{z(t, @x, \dfun, \vc)} < \ve$ $\forall$ 
$t > T_{r, \ve}$. So, the conclusion of the lemma follows.  
\epr



\bp{VdCont}
The function $\tilde V_0: \ \E!p\relaxed\to\R$ is continuous.
\eps

\bpr
Fix $@x\in\E$, $\dfun \in\relaxed$.  
Suppose $@x_k \to @x, \dfun_k\to \dfun$,
where $@x_k\in\E$.  Let $\{k_j\}$ be a subsequence of
$\{k\}$ s.t.
\be{eq1:VdCont}
\lim_{j\to\infty}\tilde V_0(@x_{k_j}, \dfun_{k_j}) =
\liminf_{k\to\infty}\tilde V_0(@x_{k}, \dfun_{k}).
\ee
$\forall$ $k$, let $\vc_k$ be s.t.
$\tilde V_0(@x_k, \dfun_k) = V_{\vc_k}(@x_k, \dfun_k)$. 
Notice that 
$\lim_{k\to \infty} V_{\vc_k}(@x_k, \dfun_k)$ exists,
because  of~\rref{eq1:VdCont}. 

By sequential compactness of $\pcow$, $\exists$
a subsequence of $\{\vc_{k_j}\}$ converging to some $\bar \vc\in\pcow$.
Without relabeling, we assume that $\vc_{k_j}\to \bar \vc$.
It then follows from Lemma~\ref{l-lsc} that
\[
V_{\bar \vc}(@x, \dfun) \le \lim\,V_{\vc_{k_j}}(@x_{k_j}, \dfun_{k_j}) = \liminf\,
\tilde V_0(@x_k, \dfun_k).
\]
Consequently, $\tilde V_0(@x, \dfun) \le \liminf\, \tilde V_0(@x_k, \dfun_k)$.  To
complete the proof, we will show that
\be{e-upsemi}
\tilde V_0(@x, \dfun) !>\limsup_{k \to \infty}
\tilde V_0(@x_k, \dfun_k).
\ee

Let $\vc$ be a control s.t.
$\tilde V_0(@x, \dfun) = V_\vc (@x, \dfun)$.
Let $\ve > 0$ be given.  By Lemma~\ref{l-sl1}, there is some
$@x_0\in!d\D$ s.t.~\rref{e-sl1-liminf} holds.  
By Lemma~\ref{l-VdContD}, $\exists$ neighborhood $U$ of $@x_0$ s.t.
\be{eq2:VdCont}
\tilde V_0(@h, @n) < \ve/4 \; \mbox{ $\forall$}
 @h!eU \cap \E, \mbox{ $\forall$}
@n\in\relaxed. 
\ee
Let $0< T <\ioslow{\dfun}(@x, \vc)$ be s.t. 
$z(T, @x, \dfun, \vc)!eU$.  $\Therefore$ since $\{z(t, @x_k, \dfun_k, \vc)\}$
converges to $z(t, @x, \dfun, \vc)$ uniformly on $[0, T]$, it follows that
$z(T, @x_k, \dfun_k, \vc)!eU$ for $k\ge K_1$ for some $K_1$.  By
Lemma~\ref{l-ioslowsc}, one may assume that $T < \ioslow{\dfun_k}(@x_k, \vc)$
$\forall$ $k\ge K_1$.  Consequently, $@h= z(T, @x_k, \dfun_k, \vc)$
is also in $\E$, so, applying~\rref{eq2:VdCont} with $@n= (\dfun_k)_T$, 
we have 
\[
\tilde V_0(z(T, @x_k, \dfun_k, \vc), (\dfun_k)_T) < \ve /4 \qquad 
\forall\, k\ge K_1,
\]
where $(\dfun_k)_T(t) = \dfun_k(T+t)$.
By the uniform convergence property of $\{z(t, @x_k, \dfun_k, \vc)\}$,
 it follows that there is some compact set $!K$ s.t.
$z(t, @x_k, \dfun_k, \vc)!e!K$ $\forall$ $k$, $t!e[0, T]$

 Using also the
uniform continuity of $\precomp(!.)$ on compacts, one sees that there
is some $K_2\ge K_1$ s.t.
\[
\int_0^T\precomp(\abs{z(s, @x_k, \dfun_k, \vc)})\,ds \psp \le \psp 
\int_0^T\precomp(\abs{z(s, @x, \dfun, \vc)})\,ds + \ve/2
\]
$\forall$ $k\ge K_2$. Thus,  \rref{dynprog} implies 
\beqn
\tilde V_0(@x_k, \dfun_k) &\le &
\int_0^T\precomp(\abs{z(s, @x_k, \dfun_k, \vc)})\,ds + \tilde V_0(z(T, @x_k, \dfun_k,
\vc), (\dfun_k)_T) \\
&\le & V_\vc(@x, \dfun) + \ve/2 + \ve/2 \psp =  \psp \tilde V_0(@x, \dfun) + \ve
\eeqn
$\forall$ $k\ge K_2$.  From this it follows that
\[
\tilde V_0(@x, \dfun) \ge \limsup\,\tilde V_0(@x_k, \dfun_k) -\ve.
\]
Letting $\ve\to 0$, one proves~\rref{e-upsemi}. 
\epr


$\forall$ $@x\in\E$, define
\be{e-vod}
\displaystyle{V_0(@x) := \sup_{\dfun\in\M_@W} \tilde V_0(@x, \dfun)}.
\ee
Note that the supremum is finite $\forall$ $@x\in\E$; in fact, 
by~\rref{V_0finite1} and~\rref{V_0finite2}, we have the upper bound
\be{e-V0b}
V_0(@x) \le @m_2(\abs{@x}) \qquad\forall\,@x\in\E,
\ee
where $@m_2$ is as in inequality~\rref{mu1mu2}. 
Also observe that the same \function $V_0$ results if the supremum
in~\rref{e-vod} is taken over the set $\relaxed$.  This follows from
the fact that $\M_@W$ is dense in $\relaxed$ and the continuity
property of $\tilde V_0$.

By sequential compactness of $\relaxed$ and by continuity of
$\tilde V_0$, we get the following.

\bc{relaxed_d}
$\forall$ $@x$ in $\E$ $\exists$  (possibly relaxed)
disturbance $\dfun$ s.t. $V_0(@x) = \tilde V_0(@x, \dfun)$, that
is,  $V_0(@x) = \max_{\bar \dfun\in\relaxed}\tilde V_0(@x,\bar \dfun)$.
  Moreover,
$V_0: \ \E\to\R$ is continuous.
\ecs

\bpr
Fix $@x!e\E$. By definition of $V_0$, $\exists$  sequence of 
disturbances $\cbrace{\dfun_k(!.)}$ s.t. 
$\tilde V_0(@x, \dfun_k) \nearrow V_0(@x)$. By sequential compactness of
 $\relaxed$, we can extract a subsequence $\cbrace{\dfun_{k_i}}$, converging 
in $\relaxed$ to some $\dfun$. By continuity of $\tilde V_0$,
\[
 V_0(@x) = \lim_{i \to \infty} \tilde V_0(@x, \dfun_{k_i}) = 
\tilde V_0(@x, \dfun),
\]
proving the first statement of the corollary.

Now fix $@x!e\E$. Take any sequence $\cbrace{@x_k} !e\E$, 
converging to $@x$ and s.t. the sequence $\cbrace {V_0(@x_k)}$ 
converges. Let $\dfun_k(!.)$ and $\dfun !e\relaxed$ be 
maximizing disturbances for $@x_k$ and $@x$, respectively, that is, 
\[
V_0(@x_k) = \tilde V_0(@x_k, \dfun_k) \mbox{  and  }
V_0(@x) = \tilde V_0(@x, \dfun).
\]
Extracting a subsequence if necessary, let $\hat \dfun$ be a limit of 
$\{\dfun_k\}$ in $\relaxed$. $\Therefore$ by continuity of $\tilde V_0$ and by 
definition of $V_0$, we have
\[
V_0(@x) \psp !>\psp \tilde V_0(@x, \hat \dfun) \psp = \psp  
\lim_{k \to \infty}  \tilde V_0 (@x_k, \dfun_k) \psp = \psp
 \lim_{k \to \infty} V_0(@x_k).  
\]
Consequently, if $\cbrace{@x_i} \subset \E$ is any sequence, converging to 
$@x$, $\therefore$ 
\[
 V_0(@x) \psp !>\psp  \limsup_{i \to \infty} V_0(@x_i), 
\]
proving upper semicontinuity of $V_0$. 

On the other hand, again by continuity of $\tilde V_0$, we have
 \[
\liminf_{k \to \infty} V_0 (@x_k) 
        !>\lim_{k \to \infty} \tilde V_0(@x_k, \dfun) = V_0(@x),
\]     
showing the lower semicontinuity of $V_0$. Thus we conclude that $V_0$ is
continuous on $\E$.
\epr        

\bll{bounded_below}
$\exists$  $\floor!e\ki$ s.t.
\[
\floor (\abs @x) !<V_0(@x)
\] 
$\forall$ $@x!e\E_1$.
\elll

Notice that if $\abs{@x} !>1.6 @r(\abs{h(@x)})$ and $@x\neq 0$, 
$\therefore$ $@x!e\E$, because 
$\abs{@x} > \abs{@x}/1.6 !>@r(\abs{h(@x)}).$
To prove Lemma~\ref{bounded_below}, we first prove a technical lemma.

\bll{tl}
 Suppose $@S\,: \, @.z = g(z, d), \; y = h(z)$ is a system
of type~\rref{dsystem}, and $p(!.)$ is a smooth \function $!e\ki$,
 s.t. the following conditions hold:
\begin{itemize}
\item $\abs{g(@x, d)} !<1 $ $\forall$ $@x!e\states$ and 
$d !e@W$,
\item $\abso{\nabla (p !o\abs h ) (@x) !.g(@x, d)} !<1 $ 
$\forall$ $d !e@W$, $@x$ with $\abs{h(@x)} !>1$,
\item $p(s) !>s$ $\forall$ $s > 0$.
\end{itemize}
%
Pick any constant $a >0$ and  define $K_0 = p(1) + (2+a)/a + 1$. $\Therefore$
$\forall$ $@x!e\states$ s.t.   
\[
\abs @x!>(1+a) p(\abs{h(@x)}) \mbox{ and } \abs @x!>K_0,
\] 
it holds that 
\[
\abs {z(t, @x, \dfun)} > p(\abs{h(z(t, @x, \dfun))})
\]
 $\forall$
 $t !e[0, 1)$  and any $\dfun !e\M_@W$.
\ells

\bpr
Fix $a$, $@x$, and $\dfun$ as in the formulation of the lemma, and
define 
\[
@q:= \min \cbrace{t \,:\, \abs{z(t, @x, \dfun)} 
!<p(\abs{h(z(t, @x, \dfun))})}
\]
with the convention $@q= +\infty$ if the inequality never holds
for $t !>0$.
Assume the lemma is false, s.t. $@q< 1$.       

 Let $@h= z(@q,@x,  \dfun )$, and
let $\hat{\dfun}$ be the shift of $ \dfun$ by $@q$, that is, 
$\hat{\dfun}(t) =\dfun(t + @q)$.   Since
$\abs{g(z, d)}\le 1$ $\forall$ $z \in\states$, $d !e@W$, it
holds that 
\[
\abs{@h} \ge \abs{@x} - @q\ge  K_0 - 1.
\]
By the definitions of $@h$ and $@q$,
one has
\be{e-eta}
p (\abs{h(@h)}) = \abs{@h} \ge K_0 - 1 ,
\ee
so also $\abs{h(@h)} \ge p ^{-1}(K_0 - 1) > 1$.
Thus, $|{h(z(s,@h, \hat \dfun ))}|>1$ $\forall$ $s$ near zero.


 {\it Claim}. \ $|{h(z(s, @h, \hat \dfun ))}| > 1$ $\forall$ $s\in
[-1, \; 0]$.

Assume the claim is false.  $\Therefore$ there must exist some $-1\le s_0 < 0$ s.t.
\[
s_0 = \max\left\{s\le 0: \ |{h(z(s,  @h, \hat \dfun ))}| \le 1\right\}.
\]
We have that $\forall$ $s!e(s_0, \, 0]$,
$|{h(z(s,  @h, \hat \dfun ))}|  > 1$.

Recall that
$\abs{\nabla(p !o\abs{h})(z )f(z, d)}!<1$ 
$\forall$ $z$ with $\abs{h(z )} \ge 1$, $d !e@W$

Thus
\[
\abs{
\frac{d}{ds}p\left(\abs{h(z(s,  @h, \hat \dfun ))}\right)
}
\le 1 \quad
\forall\,s!e(s_0, \, 0].
\]
This, in turn, implies that 
\[
 p\left(\abs{h(z(s_0,  @h, \hat d ))}\right) \ge p (\abs{h(@h)}) +
 s_0 \ge 
K_0 +s_0-1 > p(1),
\]
and so, since $p$ is strictly increasing, 
$|{h(z(s_0,  @h, \hat \dfun ))}| > 1$, thus
contradicting the definition of $s_0$.  This proves the claim.

It follows from the claim that
   $\abs{h(z(s, @x,  \dfun))} > 1$ $\forall$ $s\in[0,@q]$.
Thus,
\beqn
p (\abs{h(@h)}) &=&\abs{@h} \;\ge\; \abs{@x} -@q\\
&\ge& (1+a)p (\abs{h(@x) }) - @q 
\ge (1+a) p(\abs{h(@h)}) - (1+a)@q- @q.
\eeqn
(The last inequality used the fact that 
$\abs{\frac d{ds}p(\abs{h(z(s, @x,  \dfun))})} !<1$ $\forall$ 
$s !e[0, @q]$.) 
  It follows that $p (\abs{h(@h)}) \le \frac{2+a}{a}@q$,
so from~\rref{e-eta} we know that 
\[
K_0!<1 + p(\abs{h(@h)}) !<1+ \frac{2+a}{a}@q,
\]
 contradicting the choice of $K_0$.
This shows that it is impossible to have $@q< 1$.
\epr

{\it Proof of Lemma~$\ref{bounded_below}$}. 
Recall that if $@x\not!e\D\cup\B$, $\therefore$ 
$\tilde f(@x, d, v) = \hat f(@x, d)$ $\forall$ $d$ and $v$, s.t. 
\[
\abs{\nabla (1.5 @r!o\abs{h})(@x) !.\tilde f(@x, d, v)}
 = \frac{1.5 \abs{\nabla(@r!o\abs h)(@x) !.f(@x, d)}}
        {1 + \abs{f(@x, d)}^2 + @k(@x)} !<1,
\]
where the last inequality follows from~\rref{kappa_def}. 
$\therefore$ the assumptions of Lemma~\ref{tl} are satisfied with $p := 1.5 @r$
and  $f(@x, d) := \tilde f(@x, d, v) =  \newf (@x, d)$.  
By Lemma~\ref{tl}, we can find a constant $K_0$ s.t. if
$\abs @x> K_0$ and  $\abs @x!>1.6 @r(\abs{h(@x)})$, $\therefore$
$z(t, @x, \dfun, v) \not!e\D\cup\B$ $\forall$ $t !e[0, 1)$. In particular,
for such a $@x$ we will have $\ioslow{\dfun}(@x, \vc) > 1$ and 
$\abs{z(t, @x, \dfun, \vc)} > \abs{@x} -1$ $\forall$ positive  $t < 1$ . 
Hence, the inequality
\[
\precomp(\abs {z(t, @x, \dfun, \vc)}) >  \precomp(\abs{@x} -1) 
\; \;\; \forall t !e[0, 1)
\]
holds $\forall$ $\dfun\in\M_@W$, $\vc\in\pcow$, and any $@x$
s.t. 
\be{xi_away_from_0} 
\abs{@x} > \max\cbrace{ K_0 +1, \;1.6 @r(\abs{h(@x)})}.
\ee
$\therefore$ $\forall$ $@x$ as in~\rref{xi_away_from_0} and any
$\dfun \in\M_@W$, the  following estimate holds:
\[
V_0(@x) \; !>\; \tilde V_0(@x, \dfun) \;\geq\;
\int_0^1 \precomp(\abs {z(t, @x, \dfun, \vc)}) dt \;!>\;
\precomp(\abs @x-1). 
\]
Next, notice that $V_0$ is strictly positive on $\E$. Indeed, 
$|{\tilde f(@x, d, v)}| !<3$ $\forall$ $@x!e\E$, $d !e@W$,
and $v !e[-1, 1]^n,$  s.t. 
\[
\abs{z(s, @x, \dfun, \vc) - @x} !<\frac 1 2 \dist(@x, \D) \quad
\forall \; s !<\frac{\frac 1 2 \dist(@x, \D)}{3}, \; 
\dfun !e\M_@W,\; \vc !e\pcow.
\]
$\therefore$ 
$\ioslow{\dfun}(@x, \vc) !>\frac{1}{6}\dist(@x, \D)$ and 
$\abs{z(s, @x, \dfun, \vc)} !>@x- \frac 1 2 \dist(@x, \D)$
$\forall$ $s !<\frac{1}{6}\dist(@x, \D)$, all 
$\dfun !e\M_@W$, and $\vc !e\pcow.$ So, we have
\[
V_\vc(@x, \dfun) \;!>\;
\int_0^{{\rm dist}(@x, \D)/6} \precomp\rb{\abs{z(s, @x, \dfun, \vc)}}\,ds
\]
\[
\;\geq\;
 \frac{\dist(@x, \D)}{6}\, 
  \precomp(\abs @x- \dist(@x, \D)/2).
\]
%
This shows that 
\be{infinf}
\inf_{\dfun !e\M_@W} \inf_{v !e\pcow} V_\vc(@x, \dfun) 
!>\frac{\dist(@x, \D)}{6}\, 
        \precomp \! \rb{\abs @x- \frac{\dist(@x, \D)}{2}}.
\ee
Thus also the $\sup_{\dfun !e\M_@W} \inf_{v !e\pcow} V_\vc(@x, \dfun)$
satisfies~\rref{infinf}, and hence  the same inequality 
holds for $V_0$.  

Since $V_0$ is lower semicontinuous, it attains its minimum on any 
compact set. $\forall$ positive $l$ define
\[
r_l := \frac 1 l \max\cbrace{ K_0 +1, \;1.6 @r(\abs{h(@x)})} \; \mbox{ and } \;
\]
\[
m_l = \inf \cbrace{V_0(z)\,:\,
 z !e\E_1, \, r_l !<\abso z !<r_1} \,.
\]
Since the sequence $\cbrace{m_l}$ is nonincreasing and positive, and 
$\precomp$ is $!e\ki$, we can find a $\floor!e\ki$
 s.t. 
\[
\floor(s) < m_l \;\;\forall \, s !e[r_l, r_{l-1}], \ \forall \, l>1
\]
and 
\[
\floor(s) < \precomp(s - 1) \forall \, s !>\max\cbrace{ K_0 +1, \;1.6 @r(\abs{h(@x)})}.
\]
By construction, $\floor$ will be a lower bound for $V_0$ on $\E_1$.  
\mybox 

Combining Lemma~\ref{bounded_below} with~\rref{e-V0b}, we get the
following, using $\pot:=@m_2$:
\be{e-V0bd-2sd}
\floor(\abs{@x}) \le V_0(@x) !<\pot(\abs{@x})
\qquad\forall\,@x\in\E_1.
\ee

The following lemma and corollary  summarize the dissipation properties for
$V_0$.
\bll{lem:V0dissip:2}
$\forall$ $@x!e\states \setminus (\D \cup \B)$ and $\dfun !e\M_@W$, and 
any $t_0$ s.t. $z(t, @x, \dfun) \not !e\D\cup\B$ $\forall$ 
$t !e[0, t_0]$,
the following dissipation inequality holds:
%
\[
 V_0(z(t, @x, \dfun)) -V_0(@x) !< -\int_0^{t_0} \precomp(\abs{z(t,
@x, \dfun)})\,dt.
\] 
\ells

\bpr
Fix $@x!e\states \setminus(\D \cup \B)$, $\dfun!e\M_@W$,
 and any positive $t_0$ as in the formulation of the lemma.
Let $\ve > 0$ be given. Find $\dfun_1$ s.t. 
%
\[
V_0(z(t_0, @x, \dfun)) - \tilde V_0 (z(t_0, @x, \dfun), \dfun_1)<\ve,
\]
and let $\vc_1 !e\pcow$ be a control s.t.  
$\tilde V_0 (z(t_0, @x, \dfun), \dfun_1) = V_{\vc_1}(z(t_0, @x, \dfun), \dfun_1)$.
%
Let $\tilde \dfun$ be defined by
\[
\tilde \dfun(t) = \twoif{\dfun(t)}{\mbox{if $0\le t\le t_0$,}}
{\dfun_1(t-t_0)}{\mbox{if $t > t_0$.}}
\]

$\Therefore$ $z(t, z(t_0, @x, \dfun),  \dfun_1, \vc_1) = z(t + t_0, @x, \tilde \dfun, \tilde \vc)$
$\forall$ $t !>0$.
 By assumption,  
\[
z(t, @x, \dfun) \not  !e\D \cup \B \quad\quad \forall \;t !e[0, t_0],
\]
and $\therefore$ we have
\be{v_dnm}
z(t, @x, \dfun) = z(t, @x, \dfun, \vc) \quad\quad \forall \;t !e[0, t_0], \quad \forall \ \vc !e\pcow.
\ee 
Notice also that~\rref{v_dnm} implies that $\ioslow{\tilde \dfun}(@x, \vc) > t_0$ $\forall$ 
$\vc !e\pcow$ and 
\begin{eqnarray*}
&& \tilde V_0(@x, \tilde \dfun) =  \min_{\vc !e\pcow} 
\int_{t_0}^{\ioslow{\tilde \dfun}(@x, \vc)}\precomp\rb{\abs{z(s, @x, \tilde \dfun, \vc)}}\,ds \\
&=&  \int_0^{t_0}\precomp\rb{\abs{z(s, @x, \tilde \dfun)}}\,ds + \min_{\vc !e\pcow} 
\int_{t_0}^{\ioslow{\tilde \dfun}(@x, \vc)}\precomp\rb{\abs{z(s, @x, \tilde \dfun, \vc)}}\,ds \\
&=&  \int_0^{t_0}\precomp(\abs{z(s, @x, \dfun)})ds + \min_{\vc !e\pcow}
\int_0^{\ioslow{\dfun_1}(z(t_0,@x, \dfun), \vc)} 
\precomp(\abs{z(s, z(t_0,@x,  \dfun), \dfun_1, \vc)})ds \\
&=& \int_0^{t_0}\precomp(\abs{z(s, @x, \dfun)})ds + \tilde V_0(z(t_0, @x, \dfun), \dfun_1).
\end{eqnarray*} 
Consequently, one has
\[
\tilde V_0(z(t_0,@x,\dfun),\dfun_1) = \tilde V_0(@x, \tilde \dfun) - 
\int_0^{t_0}\precomp\left(\abs{z(s,@x,\dfun)} \right)\,ds.
\]
Thus, 
\beqn
V_0(z(t_0,@x, \dfun)) &\le& \tilde V_0(z(t_0,@x,\dfun),\dfun_1) + \ve\\
&=& \tilde V_0(@x, \tilde \dfun) - 
\int_0^{t_0}\precomp(\abs{z(s,@x,\dfun)})\,ds + \ve\\
&\le& V_0(@x) -
\int_0^{t_0}\precomp(\abs{z(s,@x,\dfun)})\,ds + \ve.
\eeqn
Letting $\ve->0$, we get the desired inequality.
\epr

Thus, we have proven that the UOSS dissipation inequality 
holds for $V_0$ along the trajectories of the slower
system $\hat @S$ which are entirely contained in 
$\states \setminus (\D \cup \B)$. It follows
 immediately that 
the same estimate holds along the trajectories of the
original system $@S$.

\bc{cor:V0dissip:2}
$\forall$ $@x!e\states \setminus (\D \cup \B)$ and $\dfun !e\M_@W$, and 
any $t_0$ s.t. $x(t, @x, \dfun) \not !e\D\cup\B$ $\forall$ 
$t !e[0, t_0]$,
the following dissipation inequality holds:
\be{V0_dissip3}
V_0(x(t,@x,\dfun))-V_0(@x) !< -\int_0^{t_0}\precomp(\abs{x(t,@x,\dfun)})\,dt.
\ee
\ecs

{\it Proof}.
Pick an initial state $@x!e\states \setminus (\D \cup \B)$, 
a disturbance $\dfun !e\M_@W$, and an appropriate $t_0$. $\Therefore$ 
%
\begin{eqnarray*}
\lefteqn{V_0(x(t_0, @x, \dfun))-V_0(x(@x))}\\
&=& V_0(z(\tslow{@x}{\dfun}(t_0), @x, \dfun 
!o\tslow{@x}{\dfun}^{-1}  )) 
  -V_0(z(\tslow{@x}{\dfun}(t_1), @x, \dfun !o\tslow{@x}{\dfun}^{-1} )) \\
 & \leq&  -\int_{0}^{\tslow{@x}{\dfun}(t_0)}
 \precomp(\abs{z(s, @x, \dfun !o\tslow{@x}{\dfun}^{-1} )})\,ds \\
& = & -\int_{0}^{t_0}
 \precomp(\abs{z(\tslow{@x}{\dfun}(t), @x, \dfun !o\tslow{@x}{\dfun}^{-1}
 )})\,d\tslow{@x}{\dfun}(t)\\
 & =& -\int_{0}^{t_0}
 \precomp(\abs{x(t, @x, \dfun)})\,\frac{d}{dt}{\tslow{@x}{\dfun}}(t)\,dt\\
&=&  -\int_{0}^{t_0} 
 \precomp(\abs{x(t, @x, \dfun)})[1+\abs{f(x(t, @x, \dfun), \dfun(t))}^2 
  + @k(x(s, @x, \dfun))]\, dt \\
 & \leq& -\int_{0}^{t_0}
 \precomp(\abs{x(t, @x, \dfun)})\,dt.\mybox
\end{eqnarray*}


\newpage
\subsection{some definitions and facts from nonsmooth analysis}

\bd{def:prosub-super}
A vector $@z!e\R^n$ is a {\em proximal subgradient}
(respectively, {\em proximal supergradient}) of the function $V: \R^n
\to (-\infty, +\infty]$ at $x$ if $\exists$ positive $@s$
s.t., $\forall$ $x'$ in some neighborhood of $x$,
%
\be{eq:prosub}
V(x') !>V(x) + @z!.(x' -x) -@s\abso{x'-x}^2
\ee
\be{eq:prosuper}
(\mbox{correspondingly, }\; V(x') !<V(x) + @z!.(x'-x) 
+ @s\abso{x' -x}^2).
\ee

The (possibly empty) set of all proximal subgradients (respectively,
supergradients) of $V$ at $x$ is called the {\em proximal subdifferential}
and is denoted $\partial_P V(x)$ (respectively, {\em proximal
superdifferential}, denoted $\partial^P V(x)$). Note that the
definitions imply that if the function $V$ is differentiable at $x$,
then both the subdifferential and superdifferential sets must be subsets
 of the 
$\{\nabla V(x)\}$.
\ed
%
\bll{ipsdissip}
%
Let $@S$ be a system of type~\rref{dsystem},  $@x$ be a vector 
in $\states$, $\bar d !e@W$, and $V  :  \states \to \R$. $\Therefore$ if 
$\exists$ continuous 
$@a_@x\,: \, \R_{!>0} \to \R$ and  $\ve > 0$ s.t. 
the following inequality holds $\forall$ $@t< \ve$,
%
\be{alpha_xi_dissip}
V(x(@t, @x, \dfun)) - V(@x) \psp !<\psp  \int_0^@t@a_@x(t)dt
\ee
(where $\dfun$ is the constant disturbance equal to $\bar d$), $\therefore$ 
$\forall$ $@z!e\partial_P V (@x)$ the proximal form of 
inequality~\rref{alpha_xi_dissip} holds:
%
\be{pro_alpha_xi_dissip}
 @z!.f(@x, \bar d) !<@a_@x(0).
\ee
\ells

\bpr
It follows from~\rref{eq:prosub} and~\rref{alpha_xi_dissip} that, 
$\forall$ $@t$ close enough to $0$ we have
%  
\[
\int_0^@t@a_@x(t)dt \psp !> \psp V(x(@t, @x, \dfun)) - V(@x)
\psp  !> \psp
@z!.(x(@t, @x, \dfun) - @x) - @s\abso{x(@t, @x, \dfun) -@x}^2.
\]
%
Dividing by $@t$ and passing to the limit as $@t$ tends to $0$, we get
\rref{pro_alpha_xi_dissip}.
\epr


\newpage
\subsection{smoothing out a continuous Lyapunov function}

\renewcommand{\xdissip}{{\Theta}}
\newcommand{\upso}{{\mbox{$\Upsilon_1$}}}
\newcommand{\upst}{{\mbox{$\Upsilon_2$}}}
\newcommand{\xddissip}{\xdissip}
\newcommand{\xddissipx}{\xddissip(x,d)}
\renewcommand{\GO}{!O}

The next result shows how to approximate a continuous function $V$ by
a locally Lipschitz one in a weak $C^1$ sense. The function $V$ is
assumed to be bounded below, or up to a translation by a constant, 
nonnegative.

\bll{contimplieslipschitz}
Let $@S:\,@.x = f(x, d)$ be a system,
with $x\in\states=\R^n$ and $d\in@W$,
a compact metric space,
s.t. $f(x, d)$ is locally Lipschitz in $x$ uniformly on
$d$ and jointly continuous in $x$ and $d$.
Assume that we are given
\begin{itemize}
\item  an open subset $\GO$ of $\states$;
\item a continuous, nonnegative \function $V: \GO \to \R$ 
satisfying 
\be{dissip:loc}
 @z!.f(x, d) !<\xddissipx
 \;\;\forall x !e\GO, \;
@z!e\prosub V(x), \; d \in@W
\ee
with some continuous \function $\xddissip \,:\, \GO!p@W\to \R$;
%
\item two positive, continuous \functions $\upso$ and $\upst$ on 
$\GO$.
\end{itemize}

$\Therefore$ $\exists$ {\afunction} $\tilde V: \GO \to \R$, locally Lipschitz on 
$\GO$, s.t.
%
\be{tilde_V_close_to_V}
0 !<V(x) -\tilde V(x) !<\upso(x) \;\;\forall x !e\GO
\ee
and 
\be{tilde_V_dissip}
L_{f_d}\tilde V(x) !< \xddissipx + \upst(x)
\;\; \mbox{ $\forall$}d \in
@W\mbox{ and for almost all } x !e\GO,
\ee
%
where $f_d$ is the vector field defined by $f_d = f(!., d)$. 
\elll

Note that, by Rademacher's theorem, the directional derivative 
$L_{f_d} \tilde V (x)$, given by $\nabla V(x) !.f(x, d)$,
is defined for almost all $x$ because $\tilde V$ is locally Lipschitz.  
The proof will follow closely along the lines of the proof of the similar
result for CLFs, found in Clarke-Ledyaev-Sontag-Subbotin

We will first prove a ``local'' version of the result.
$\forall$ $K$ which is
 a compact subset of $\GO$ and $r>0$,
we 
introduce the following notations: 
\begin{itemize}
\item $\bar B_r(K) := \cbrace{x!e\states \,:\, \exists @x!eK \mbox{ with }
\abs{x-@x} !<r}$, the closed $r$-fattening of $K$, for $r > 0$,
\item $@b(K):=\sup_{x !eK}V(x),$

\item $m_{K} := 
\frac{1}{4}\min\cbrace{\upst(x): x !eK},$

\item $\ell_K :=$ a  Lipschitz constant for $f$ with respect to $x$
 in  $K$, that is,
\[
|f(x, d) - f(x', d)| !<\ell_K |x -x'| \;\;\forall d !e@W, \;
         \forall  x !eK,
\]

\item $@w_K(!.) :=$ the modulus of continuity of $V$ on 
$K$, that is,
\[
@w_K(@d):=\sup\cbrace{V(x)-V(x')\,:\, |x-x'| !<@d; x, x' \in
K},
\]
%
\item $@p_K(!.) :=$ the modulus of continuity of $\xddissip$
 on $K!p@W$, that is,
\beqn
& &@p_K(@d):=\sup\{\xddissip(x , d)-\xddissip(x',d')\,:
\, |x-x'| + {\rm dist\,}(d,d') !<@d; \\
& &\hspace{1.6in} x, x' !eK, d,d'\in@W\}.
\eeqn
\end{itemize}

To approximate the given continuous \function $V$ by a locally 
Lipschitz one, we
would like to use the
notion of ``Iosida--Moreau inf-convolution,'' well known in convex analysis.
 Fix a parameter
$@a!e(0, 1]$. Suppose for the moment that $V$ is defined on
the whole $\states$. Define
%
\[
V_@a(x) := \min_{y !e\states}
\left[V(y) + \frac{1}{2@a^2}|y-x|^2 \right].
\]
%
$\forall$ fixed $x$, the set of points $y$ where the minimum is attained
is nonempty because $V$ is bounded below. Denote one of them by $y_@a(x)$.

Fix a compact $K$ and  $@a!e(0, 1]$, let 
\[
K_{@a} := {\bar B}_{@a\sqrt{2@b(K)}}(K).
\]
The following four claims summarize some of the useful properties of
$V_@a$, proved in Clarke-Ledyaev-Sontag-Subbotin

{\em Claim $1$}.  $\forall$ $x !eK$,
%
\[
\abs{y_@a(x)-x}^2 \psp !< \psp 
\min \cbrace{2 @a^2 @b(K), \;
2@a^2 @w_{K_{@a}}
\left(@a\sqrt{2@b(K)}\right)}.
\]

{\em Proof of Claim $1$}. By definition of $V_@a$ and $@b(K)$,
we have
%
\be{est_min}
\frac{1}{2@a^2}\left|y_@a(x)-x \right|^2 !<V(x)-V(y_@a(x))
!<V(x) !<@b(K), 
\ee
s.t. $\left|y_@a(x)-x \right|^2 !<2@a^2@b(K)$.
On the other hand, the first inequality in~\rref{est_min} implies
also that 
\begin{eqnarray*}
\left|y_@a(x)-x \right|^2 &\leq& 2@a^2 (V(x)-V(y_@a(x))) \\
&\leq& 2@a^2@w_{K_{@a}}
\left(\left|y_@a(x)-x \right|\right)
\psp !<\psp  2@a^2 @w_{K_{@a}}
\left(@a\sqrt{2@b(K)}\right),
\end{eqnarray*}
proving the claim.

Let \[
@z_@a(x) = \frac{x-y_@a(x)}{@a^2}.
\]

{\em Claim $2$}.  
$\forall$ $x !e\states$, 
\be{zeta_in_prosub}
@z_@a(x) !e\prosub V(y_@a(x)) \;\;\mbox{and}
\ee
\be{zeta_in_prosuper}
@z_@a(x) !e\prosuper V_@a(x). 
\ee

{\em Claim $3$}.
$\forall$ $x !eK$, 
\[
 V_@a(x) !<V(x) !<V_@a(x) + 
@w_{K_{@a}}\left(@a\sqrt{2@b(K)}\right).
\]

{\em Claim $4$}.
$V_@a$ is locally Lipschitz.
%

Now recall that, in our setting, $V$ is defined on an open subset $\GO$
of $\states$, so,
we can't minimize the expression in the definition of $V_@a$ over
the whole state space. However, $\forall$ compact subset $K$ of
 $\GO$ we
can choose $@a$ small enough s.t.
\[
K_{@a} \subset \GO,
\]
 hence, $\forall$ $x$ in $K$
we could define $V_@a$ minimizing over $y !e\GO$,
 and the same \function $V_@a$ results on $K$. 

\bll{sub_est}
Assume that $V$ satisfies~\rref{dissip:loc}, a compact $K$ is fixed, and 
  $@a!e(0, 1]$ is chosen to
 satisfy 
\begin{itemize}
\item   
$K_{@a} = {\bar B}_{@a\sqrt{2@b(K)}}(K) \subset \GO$,
\item $@p_{K_{@a}}\left(@a\sqrt{2@b(K)}\right) !<m_K$, and
\item $2\ell_{K_{@a}}@w_{K_{@a}}\left(@a\sqrt{2@b(K)}\right)
 !<m_K$.
\end{itemize}
$\Therefore$ the function 
\be{V_alpha_def}  
V_@a(x) := \inf_{y !e\GO}
\left[V(y) + \frac{1}{2@a^2}|y-x|^2 \right]
\ee
will possess the following property:

$\forall \ x !eK, \forall \ d !e@W$, and $\forall \  
@z!e\prosub V_@a(x)$,
\[
@z!.f(x, d) !<\xddissipx + 2m_K.
\]
\ells

\vspace{-.04pc}
{\it Proof}.
Fix any $x !eK$. The choice of $@a$ ensures that the infimum 
in~\rref{V_alpha_def} is a minimum, and it is achieved at some 
$y_@a(x) !eK_{@a}$. The
definition of $@z_@a(x)$ and Claim 1 imply that
\be{zeta<omega1}
\abso{@z_@a(x)}\abso{y_@a(x) -x} = 
\frac{\abso{y_@a(x) -x}^2}{@a^2} !<
2@w_{K_{@a}}\left(@a\sqrt{2@b(K)}\right).
\ee
Using again the fact  that 
$y_@a(x) !eK_{@a}$, 
%
\[
\abso{f(x, d)-f(y_@a(x), d)} !<\ell_{K_{@a}} \abso{x-y_@a(x)}.
\]
Combining the last inequality with~\rref{zeta<omega1} we obtain
\be{zeta<omega2}
\abso{@z_@a(x)}\abso {f(x, d)-f(y_@a(x), d)} \leq
2\ell_{K_{@a}}@w_{K_{@a}}\left(@a\sqrt{2@b(K)}\right).
\ee
Now, by Claim 2, $@z_@a(x) !e\prosub V (y_@a(x))$. Hence,
\begin{eqnarray*}
@z_@a(x) !.f(y_@a(x), d) &\leq&
   \xddissip(y_@a(x), d) \\ % + \ydissip(\abs{h(y_@a(x))}))\\
&\leq& \xddissipx + @p_{K_{@a}}(\abso{x-y_@a(x)})\\ 
%
&\leq& \xddissipx + @p_{K_{@a}}\left(@a\sqrt{2@b(K)}\right)\\
%
&\leq&  \xddissipx + m_{K}.  %+\ydissip(\abs{h(x)}).
\end{eqnarray*}
So, by the last inequality, \rref{zeta<omega2}, and the choice of
$@a$, we have
%
\begin{eqnarray}
@z_@a(x) !.f(x, d) &= &
@z_@a(x) !.f(y_@a(x), d) +
@z_@a(x) !.\rb{f(x, d) -f(y_@a(x), d)} \nonumber \\
&\leq& 
 \xddissipx + m_{K}  + m_{K}. \label{V_alpha-dissip}
\end{eqnarray}

Next, we show that 
$\prosub V_@a(x) !c\left\{@z_@a\right\}$
$\forall$ $x !eK$.
Pick any $@z!e\prosub V_@a(x)$. By definition of the proximal
subgradient, the inequality
\be{prosub:eq}
@z!.(y-x) !<V_@a(y) -V_@a(x) + o(\abso{y-x})
\ee 
holds $\forall$ $y$ near $x$.
Since $@z_@a(x)$ is a proximal supergradient of $V_@a$ at
$x$, we also have
\be{prosuper:eq}
-@z_@a(x) !.(y-x) !< -V_@a(y) + V_@a(x) + o(\abso{y-x}).
\ee
Adding \rref{prosub:eq} and \rref{prosuper:eq}, we get
\be{sub+super}
(@z-@z_@a(x))!.(y-x) !<o(\abso{y-x})
\ee
$\forall$ $y$ sufficiently close to $x$.
Substituting $y=x+h(@z-@z_@a(x))$ in \rref{sub+super}
 and letting $h$ tend to $0$, 
we arrive at   $@z=@z_@a(x)$.
\mybox

Now we are ready to prove the main lemma of the section.

{\it Proof of Lemma~$\ref{contimplieslipschitz}$}.
For every $x !e\GO$ find an $r_x > 0$ small enough s.t.
$\bar B_{r_x}(x) \subset \GO$. $\Therefore$ the collection of open balls
$\cbrace{B_{r_x} (x), \, x !e\GO}$ forms an open covering of $\GO$.  
By paracompactness of $\GO$ we can find a locally finite refinement
$\cbrace{B_i, i !e\N}$ of $\cbrace{B_{r_x} (x), \, x !e\GO}$
%(cf.~\cite[Lemma 4.1]{boothby}).

Moreover, since
$\cup_{x !e\GO} B_{r_x} =\GO$, we also have $\cup_{i} B_i =\GO$.
Let $\cbrace{\vf_i, i !e\N}$ be a partition of unity,
 subordinate to $\{B_i\}$.
$\forall$ index $i$, let 
\[
\J_i = \cbrace{ j !e\N \,:\, B_i \cap B_j \neq \emptyset }.
\]
Notice that $\J_i$ is finite $\forall$ $i$, because of local finiteness of 
the covering $\{B_i\}$. For every $i !e\N$ define
%
\[
M_i : =  \sup \abso{\nabla \vf_i( x)}\abso{f(x, d)},
\]
where the supremum is taken over all 
$x !e\cup_{j !e\J_i} \bar B_j$, $d !e@W$; and
\[
N_i = \max_{j !e\J_i} {\rm card} (\J_j),
\]
that is, $N_i$ denotes the maximum
cardinality of $\J_j$ for $j$ s.t. $B_j$ intersects $B_i$.
Next, $\forall$ compact $K = \bar B_i, \,i !e\N$, choose an $@a_i$ 
as in the formulation of Lemma~\ref{sub_est}, satisfying the following
two additional conditions: 
\be{add_cond1}
N_i \,
@w_{K_{@a_i}}\left(@a_i \sqrt{2@b(\bar B_i)}\right)M_{i} 
< \frac{1}{2} \inf_{x !e\bar B_i}\upst(x)
\ee
and
\be{add_cond2}
@w_{K_{@a_i}}\left(@a_i\sqrt{2@b(\bar B_i)}\right) 
< \frac{1}{2}\inf_{x !e\bar B_i}\upso(x)
\ee
(where $K_{@a_i}$ denotes 
$B_{@a_i\sqrt{2@b(\bar B_i)}}(\bar B_i)$).

Next,  $\forall$ $i !e\N$, define $V_{@a_i}$ on $\bar B_i$ as 
in~\rref{V_alpha_def} (this can be done, because, by the choice of $@a_i$,
$K_{@a_i} !c\GO$). Since $V_{@a_i}$ is locally Lipschitz, it is
differentiable almost everywhere by Rademacher's theorem. 
Lemma~\ref{sub_est} implies that for almost all $x !e\bar B_i$, $d !e@W$
\[
L_{f_d}V_{@a_i}(x) !<\xddissipx + \upst(x)/2.
\] 

Define \[
\tilde V:= \sum_{i=1}^{+\infty} V_{@a_i} \vf_i.
\]
Strictly speaking, this does not make sense, because $V_{@a_i}$ is
not defined outside $\bar B_i$, but that does not matter because $\vf_i$
vanishes outside $B_i$ anyway.

Since each $V_{@a_i}(x) !<V(x)$ $\forall$ $x !eB_i$ and $\vf_i$'s 
add up to $1$,  the definition of $\tilde V$ shows that 
$\tilde V(x) !<V(x)$ $\forall$ $x !e\GO$. It is also clear from the 
definition of $\tilde V$ that $\tilde V$ is locally Lipschitz on $\GO$. 
 
We claim that for almost all $x !e\GO$, $d !e@W$ the
following hold:
\begin{enumerate}
\item $0 !<V(x) - \tilde V(x) !<\upso(x)$;

\item $\nabla \tilde V(x) !.f(x, d) !< \xddissipx + \upst(x).$ 
\end{enumerate}
Take any $x !e\GO$ and find $i !e\N$ s.t. $x !eB_i$. 
Define also $\J_x := \cbrace{j !e\N\,:\, x !eB_j}$. Note that
$\J_x !c\J_i$ and that 
 $\tilde V(x):= \sum_{j !eJ_x} V_{@a_j}(x) \vf_j(x).$
$\Therefore$
 Claim 3 and  the choice of 
$@a_j$ (condition~\rref{add_cond2}) imply 
\[
V(x) - V_{@a_j} (x) \leq
 @w_{K_{@a_j}}\left(@a_j\sqrt{2@b(\bar B_j)}\right)
  !<\upso(x)/2
\]
$\forall$ $j !e\J_x$.
Thus,
\begin{eqnarray*}
V (x) - \tilde V(x)
& =  &  V(x) -  \sum_{j=1}^{+\infty} V_{@a_j} \vf_j \\
& =  &  V(x) - \sum_{j !e\J_i} V_{@a_j} \vf_j\\
&\leq&  V(x) - \min_{j !e\J_x}\cbrace{ V_{@a_j}(x)}
         \rb{\sum_{j !e\J_x}\vf_j(x)}\\  
& = &   V(x) - \min_{j !e\J_x}\cbrace{ V_{@a_j}(x)} \\
&!<& \frac{\upso(x)}{2},
\end{eqnarray*}
proving the first statement.

To prove the second statement, write
\begin{eqnarray}
&&L_{f_d}\tilde V(x) = \sum_{j !e\J_x}L_{f_d}V_{@a_j}(x)\,\vf_j(x) + 
 \sum_{j !e\J_x} V_{@a_j}(x) \,L_{f_d}\vf_j(x) \nonumber \\
%
&=& \sum_{j !e\J_x}L_{f_d}V_{@a_j}(x)\,\vf_j(x) 
+  \sum_{j !e\J_x} V_{@a_j}(x) \,L_{f_d}\vf_j(x) - 
 V(x)\sum_{j !e\J_x}L_{f_d}\vf_j(x) \label{l2} \\
%
&\leq& \max_{j !e\J_x} L_{f_d} V_{@a_j}(x)\sum_{j !e\J_x}\vf_j(x)  
+ \sum_{j !e\J_x} \rb{V_{@a_j}(x) - V(x)} \,L_{f_d}\vf_j(x) \nonumber \\
%
&\leq& \xddissipx+  \upst(x)/2 + \sum_{j !e\J_i} 
\abs{V_{@a_j}(x) - V(x)} \,\abs{L_{f_d}\vf_j(x)} \nonumber \\
%
&\leq& \xddissipx+  \upst(x)/2 +
\sum_{j !e\J_i} M_j @w_{K_{@a_j}}
\rb{@a_{j}\sqrt{2@b(\bar B_j})}
\label{nttl} \\
& !<& \xddissipx + \upst(x)/2 + \sum_{j !e\J_i} \frac{\upst(x)}{2 N_j}
\label{l3}\\
&\leq& \xddissipx + \upst(x)/2 + \upst(x)/2 \label{l}\\
&\leq& \xddissipx + \upst(x),  \nonumber
\end{eqnarray}
%
where the equality~\rref{l2} follows from the fact that
$ \sum_{j !e\J_x}L_{f_d}\vf_j(x) = 0$,
inequality~\rref{nttl} follows from Lemma~\ref{sub_est} 
and Claim 3; 
 inequality~\rref{l3} follows from the choice of $@a_{i}$
(condition~\rref{add_cond1});  and~\rref{l} is implied by the fact that 
\[
{\rm card} \J_i !<N_j \;\;\;\forall j !e\J_i.
\]  
This completes the proof of the lemma.
\mybox

The lemma we have just proved will provide  the 
``continuous $=>$ locally Lipschitz away from zero'' step in the 
smoothing process. To obtain a smooth Lyapunov function, we will
 use the following simple smoothing result.
The proof is given, for the special case when $@a$ does not depend on $d$,
in Lin-Sontag-Wang, but the general case ($@a$ depends on $d$) is proved in
exactly the same manner, so we omit the proof here.

\newcommand{\Dold}{@W}
\bll{lipschitzimpliessmooth}
 Let $\cal O$ be an open subset of $\R^n$, and let $\Dold$ be a compact subset
of $\R^l$, and assume the following as given:
\bi
\item
a locally Lipschitz \function $@F: \;\cal O \,->\,\R$;
\item
a continuous map 
$f: \; \R^n!p\Dold \;->\; \R^n, \ (x, \, d) \,|->
f(x, \, d)$ which is locally Lipschitz on $x$ uniformly on $d$;
\item
a continuous \function 
$@a: \,  !O!p@W->\R$ and continuous \functions
$@m, \, @n: \, !O ->\R_{> 0}$
\ei
s.t. $\forall$ $d\in\Dold$,
\be{e-a-200}
L_{f_d}@F(@x) \le @a(@x,d)\;\;  \mbox{  almost everywhere} \ @x\in\cal O\,,
\ee
where $f_d$ is the vector field defined by 
$f_d = f(!.,\, d)$
(recall that $\nabla @F$ is defined almost everywhere, since $@F$ is locally Lipschitz
by Rademacher's theorem).
$\Therefore$ $\exists$  smooth \function $@Y: !O->\R$ s.t.
\[
\abs{@F(@x) - @Y(@x)} \,< \, @m(@x)\; \ \ 
\forall\, @x\in\cal O\,
\]
and $\forall$ $d\in\Dold$,
\[
L_{f_d}@Y(@x) \, \le\, @a(@x,d) \,+ \,@n(@x)\; \ \
\forall\,@x\in\cal O\,. %\eqno{\Box}
\]
\ells


The next result immediately follows by Lemma~\ref{lipschitzimpliessmooth}.

\bc{contimpliessmooth}
Under the assumptions of Lemma~$\ref{contimplieslipschitz}$ there also exists a 
smooth \function $\hat V$ on $\GO$, 
satisfying inequalities
\be{hat_V_close_to_V}
\abs{ V(x) -\hat V(x)} !<\upso(x) \;\;\forall x !e\GO
\ee
and
\be{hat_V_dissip}
L_{f_d}\hat V(x) !< \xddissipx + \upst(x)
\;\; \mbox{ $\forall$}d \in
@W, 
\mbox{ $\forall$} x !e\GO.
\ee  
\ecs

\bpr
Suppose that $\GO$, $V$,  $\upso$, $\upst$ are given and the 
assumptions of
Lemma~\ref{contimplieslipschitz} hold. Replacing $\upso$ by $\upso/2$ and
$\upst$ by $\upst/2$, and applying Lemma~\ref{contimplieslipschitz},
we can find a locally Lipschitz \function $\tilde V$, defined on $\GO$
and satisfying~\rref{tilde_V_close_to_V} and~\rref{hat_V_dissip} 
with  $\upso/2$ and  $\upst/2$ instead of $\upso$ and $\upst$. 
Next, Lemma~\ref{lipschitzimpliessmooth}, applied 
with the same $\cal O$ and with
$@F:= \tilde V$, $@a:= \xddissip + \upst/2$,  $@m:= \upso/2$, and 
$@n:= \upst/4$, furnishes a smooth \function $\hat V := @Y$ as needed.
\epr

In section~\ref{dV0} we have constructed {\afunction} $V_0$, 
satisfying inequalities~\rref{e-V0bd-2sd} on $\E_1$ and~\rref{V0_dissip3}
along all trajectories of $@S$, contained in $\E_1$. 
$\therefore$ by Lemma~\ref{ipsdissip}, the proximal 
inequality~\rref{pro_alpha_xi_dissip} holds for $V_0$ at any interior 
point of $\E_1$. $\Therefore$ Corollary~\ref{contimpliessmooth}, applied with
$\GO:= {\rm int }\ \E_1$, $\upso(!.) := \floor(\abs !.)/2$,
 $\upst(!.) := \precomp(\abs !.) /2$, 
$\xdissip (x,d) := \precomp(\abs x)$, 
provides a smooth \function
\[
V_1 : {\rm int }\ \E_1 \to \R_{>0}; \quad  V_1:= \hat V_0,
\]
 satisfying the following two
 conditions $\forall$ $@x!e{\rm int }\ \E_1$:
\[
\floor(\abs @x)/2 !<V_1(@x) !<\pot(\abs @x) +  \floor(\abs @x)/2
\]
 (we will replace $\floor(\abs !.)/2$ by $\floor(!.)$ and 
$\pot(\abs !.) +  \floor(\abs !.)/2$ by $\pot(\abs !.)$ from now 
on to avoid cluttering the notation),
\be{V0_dissip3smooth}
 L_{f_d} V_1(@x) !< - \precomp_1(\abs @x) 
\quad \forall d !e@W,
\ee
 where $\precomp_1(!.) ==\precomp(!.) /2$.


\newpage
\subsection{extending to the rest of $\states$ and smoothing at the origin} \label{dV2}
\renewcommand{\xdissip}{@a_3}
To construct a UOSS-Lyapunov-like function defined on the whole
 $\states$, we must ``patch'' $V_1$ with some smooth, proper, and
 positive definite function s.t. the dissipation inequality still
 holds.

\bll{extending}
Suppose $@S$ is a system of type~\rref{dsystem}, and $@r!e\ki$.
Define  $\intEO:=\cbrace {x !e\states\,:\, \abs{x} > 2@r(\abs{h(x)})}$ 
and suppose  that $V_1\,:\, \intEO \to \R_{ !>0}$ is a 
smooth function satisfying, with some suitable $\ki$ functions,
the inequality
\be{floortop}
\floor(\abs @x) !<V_1(@x) !<\pot(\abs @x) \quad \quad 
\ee
and 
inequality~\rref{V0_dissip3smooth} on $\intEO$. $\Therefore$ $\exists$   
Lyapunov-like function $V_2$ for $@S$, smooth away from the origin,
 and a class $\ki$ \function $@F$, 
s.t. 
\[
V_2(x) = @F!oV_1(x) \quad \forall \; x \mbox{ s.t. }
  \abs{x} > 3@r(\abs{h(x)}),
\]  
and the following dissipation inequality holds 
with some $\check \xdissip !e\ki$, $\check \ydissip !e\kk$: 
\be{check-dissip}
\nabla V_2(@x) !.f(@x, d) !< - \check \xdissip(\abs{@x})
                         + \check \ydissip(\abs{h(@x)}) 
\quad\quad \forall @x\neq 0, \quad \forall d !e@W.
\ee
\ells

\bpr Let
 \[
\E_2:=\left\{@x :  \abs{@x} > 3@r(\abs{h(@x)})\right\}.
\]
Since the sets $\cbrace{@x\,:\,  \abs{@x} \geq
3@r(\abs{h(@x)})}$ and $\cbrace{@x\,:\, \abs{@x} \leq
2@r(\abs{h(@x)})}$ are disjoint and closed in the topology
 of $\states \setminus \cbrace 0$, one can find a smooth \function
$@f$ : $\states \setminus \left\{0\right\} \to
[0, 1]$ with the property that 
%
\[@f(@x) = \twoif {1} {\; \mbox{if} \; \abs{@x} \geq
3@r(\abs{h(@x)}),}
{0} {\; \mbox{if} \; \abs{@x} \leq
2@r(\abs{h(@x)})}
\] 
and $@f$ is nonzero elsewhere. It is easy to see that
$\abs{\nabla@f(x)}$ is bounded above by $@n_2!e\kk$ of $\abs{x}$
outside the unit ball centered at 0. One can also find a smooth,
strictly increasing \function $@n_1 \;\colon [0,1] \to \R_{!>0}$, s.t.
$@n_1(0) =0$ and $\abs{\nabla@f(x)} !<\frac{1}{@n_1(\abs{x})}$
$\forall$ $x$ s.t. $0<\abs x !<1$.

%
Let $@n_3!e\kk$ s.t. $\max_{d !e@W}
%11/27                                                            ^^^^^
\abs{f(@x, d)} !<@n_3(\abs @x)$. Take any smooth
\function $@p_1 \colon [0, \pot(1)] \to \R_{!>0}$ with 
$@p_1'(s) > 0$ $\forall$ $s !e(0, \pot(1))$, s.t.
\[
\frac {@p_1(\pot(r))}{@n_1(r)} < s(r), \quad 
\frac {@p_1(r^2)}{@n_1(r)} < s(r) 
\]
for some $s!e\kk$ and $\forall$ $0 < r !<1$

Let $@p_2!e\kk$ s.t. $@p_2(r) \leq
@p_1'(r)$ $\forall$ nonnegative  $r !<1$.

Let $@F(r)= \int_0^r@p_2(r_1)dr_1$. $\Therefore$ $@F(r) !<@p_1(r)$ for
all $r !<1$, s.t. $\frac {@F(\pot(r))}{@n_1(r)} < s(r)$ and
$\frac{@F(r^2)}{@n_1(r)} < s(r)$ $\forall$ $r !e(0, 1]$.
%

Now let 
\be{def:V2}
V_2(@x) = @f(@x)@F(V_1(@x)) + (1-@f(@x))@F(\abs@x^2).
\ee

If $\abs{@x} > 3 @r(\abs{h(@x)})$, $\therefore$ $V_2 ==@F!oV_1$
 in a neighborhood of $@x$,  s.t., $\forall$ $d !e@W$,
\be{xi>3rho}
\nabla V_2(@x) !.f(@x, d) 
=      @F'(V_1(@x))[\nabla V_1(@x) !.f(@x, d)] 
!< - @p_2(\floor(\abs @x)) \precompone (\abs @x).
\ee

On the other hand,  if  
$ \abs{@x} !< 3@r(\abs{h(@x)})$, $\therefore$ 
\begin{eqnarray*}
&&{\nabla V_2(@x) !.f(@x, d)}
= [\nabla @f(@x) !.f(@x, d)]\, @F(V_1(@x)) +  
@f(@x)\,@F'(V_1(@x))\, [\nabla V_1(@x) !.f(@x, d)] \\
& &\;\qquad \quad - [\nabla @f(@x) !.f(@x, d)]\,@F({\abs @x}^2)
+ (1 -@f(@x))\,2@F'({\abs @x}^2)\,[@x!.f(@x, d)]\\
&\leq& \abs{\nabla @f(@x)}\abs{f(@x, d)}\, @F(V_1(@x)) 
 + \abs{\nabla @f(@x)} \abs{ f(@x, d)}\,@F({\abs @x}^2) \\
& &\qquad\quad
 {}+ 2\abs{(1 -@f(@x))}\,@F'({\abs @x}^2)\,\abs{@x} \abs{f(@x, d)}\\
&\leq& \abs{\nabla @f(@x)}\abs{f(@x, d)}\, @F(\pot(\abs @x))
 + \abs{\nabla @f(@x)} \abs{ f(@x, d)}\,@F({\abs @x}^2)\\
& &\,\qquad \quad + 2@F'({\abs @x}^2)\,\abs{@x} \abs{f(@x, d)},
\end{eqnarray*}
where the first inequality follows from the fact that 
$@f(@x)\,@F'(V_1(@x))\, [\nabla V_1(@x) !.f(@x, d)] !<0$.
Next, the definition of $@F$ provides the following bounds for the three
terms in the right-hand side of the last inequality:
\[
@F(\pot(\abs @x))
\abs{\nabla@f(@x)} \abs{f(@x, d)} \leq
\max \left\{s(\abs @x), @F(\pot (\abs @x))@n_2(\abs @x)
\right\}@n_3(\abs @x),
\]
\[
2@F'(\abs @x^2)[\abs{@x}\abs{f(@x, d)}] \leq
 2@p_2(\abs @x^2) \abs@x@n_3(\abs @x),
\]
\[
@F(\abs{@x}^2) \abs{\nabla @f(@x)} \abs{f(@x, d)}
!<\max 
\left\{s(\abs @x), @F(\abs @x^2)@n_2(\abs @x)
\right\}@n_3(\abs @x).
\]
Define $\check \xdissip$,  $\check \ydissip$, $\check @a_1$, 
and $\check @a_2$ by
\[
\check \xdissip(r) := @p_2(\floor(r)) \precompone (r),
\]
\begin{eqnarray*}
\check \ydissip(r) &:=&  2@p_2((3@r(r))^2)  3@r(r) @n_3(3@r(r)) \\
 && + \max \left\{s(3@r(r)), @F((3@r(r))^2)@n_2(3@r(r))
\right\}@n_3(3@r(r)) + \check \xdissip(3@r(r)),
\end{eqnarray*}
\[
\check @a_1(r) = \min\left\{@F(r^2), @F!o\floor (r)\right\},
\]
 and
\[
\check @a_2(r) = \max\left\{@F(r^2), @F!o\pot(r) \right\}.
\]
$\Therefore$ inequalities
\be{check-V_proper}
\check @a_1(\abs{x}) !<V_2(x) !<\check @a_2(\abs x) \quad
\forall \,x \neq 0 
\ee 
and~\rref{check-dissip} hold for $V_2$. 
Define $V_2(0):=0$. 
 $\Therefore$ inequality $@a_1(\abs{@x}) !<V(@x) !<@a_2 (\abs{@x})$
 holds for $V_2$ with 
$@a_1:=\check @a_1$, $@a_2:=\check @a_2$ on 
$\states$, and in particular implies that $V_2$ is continuous as $0$.
$\Therefore$ $V_2$ is a UOSS-Lyapunov-like
function for $@S$, smooth away from the origin.  
\epr 

Recall that $V_2(@x) ==@F(V_1(@x))$ $\forall$ $@x$ with 
$\abs @x> 3@r(\abs{h(@x)})$. $\therefore$  
inequality~\rref{V0_dissip3smooth} implies that 
\be{y_term_d.n.m_1}
\nabla V_2 !.f(@x, d) !< -\xdissip(\abs{@x})
 \mbox{ $\forall$} \abs @x> 3@r(\abs{h(@x)}), \mbox{ $\forall$} d \in
 @W.   
\ee


\bp{smooth_V_constr}
Suppose that a system $@S$ of type~\rref{dsystem} admits a 
continuous UOSS-Lyapunov-like function $V_2$, smooth away from $0$ 
and satisfying inequalities~\rref{check-V_proper}, \rref{check-dissip}, 
and~\rref{y_term_d.n.m_1}. $\Therefore$ $@S$ admits a UOSS-Lyapunov function.   
\ep

 The basic idea used to obtain a
Lyapunov function, smooth on the whole $\states$, is  composing
 $V_2$ with some appropriately chosen 
 $@b!e\ki$.  This technique was previously utilized
 in Lin-Sontag-Wang

We will need the
 following generalization of Lemma 4.3 from there, where this function 
 $@b$ is constructed. In our setting we also need 
 the derivative of $@b$ to $!e\ki$, which was not required
 in Lin-Sontag-Wang.   

The proof is only
 a slight modification of the proof of the mentioned lemma.

\bll{from_LSW}
  Assume that
$
V : \R^n \longrightarrow \R_{\ge 0}
$
is $C^0$, positive definite, and the restriction
$V|_{\R^n\setminus\cbrace 0}$ is $C^\infty$.

$\Therefore$ given any $m !e\N$ $\exists$  $@b_m!e\ki$,
smooth on $(0,\infty)$, satisfying the following conditions:
\begin{itemize}
\item $@b_m^{(i)} (t) \ra 0$ as $t\ra 0^+$ $\forall$ $i = 0, 1, !...$,

\item $@b_m^{(i)} !e\ki \; \any i !<m$,

\item $W_m \deq  @b!oV$
is a $C^\infty$ function on all of $\R^n$.
\end{itemize}
\ells
We now return to the proof of Proposition~\ref{smooth_V_constr}. 

\bpr
Take the function $V_2$, and 
apply Lemma~\ref{from_LSW} to get {\afunction} $@b_1$,
with derivative in class $\kk$, s.t. 
\[
V_3:= @b_1 !oV_2
\] 
is smooth on $\states$. 
%REV1
%($\A = \cbrace 0$ in our case)
 Let $@a_1(!.) := @b_1 (\check @a_1(!.)/2)$,
$@a_2(!.) := 
@b_1 (\check @a_2 (!.)+ \check @a_1(!.)/2)$. 
$\Therefore$ $ @a_i !e\ki$ and 
\[
 @a_1(\abs x) !<V_3(x) !<  @a_2(\abs x).
\]

Furthermore, $\forall$ $x !e\states \setminus \cbrace 0$ we have
\begin{eqnarray*}
L_{f_d}V_3(x) &=& @b_1'(V_2(x))\sqb{L_{f_d}V_2}(x) \\
&\leq& @b_1'(V_2(x))\rb{- \xdissip(\abs x)/2 
            + \check \ydissip(\abs{h(x)})} \\
&\leq&  -@b_1'(\check @a_1(\abs x)/2) \check \xdissip(\abs x)/2 + 
@b_1'(\check @a_2(\abs x) + \check @a_1(\abs x)/2)\check \ydissip(\abs{h(x)}).
\end{eqnarray*}
Recall that, because of~\rref{y_term_d.n.m_1},
 the $\check \ydissip$ term in the last estimate can be dropped if 
$\abs x > 3@r(\abs{h(x)})$, s.t. we can write
\[
L_{f_d}V_3(x) !< -@b_1'(\check @a_1(\abs x)/2) \check 
\xdissip(\abs x)/2 + 
@b_1'(\check @a_2(3@r(\abs {h(x)})) + \check @a_1(3@r(\abs {h(x)}))/2)
 \check \ydissip(\abs{h(x)}).
\]

Thus $V_3$ is a smooth UOSS-Lyapunov function, satisfying the
dissipation inequality with 
\[
\xdissip(!.): =
 @b_1'( \check @a_1(!.)/2) \, \check \xdissip(!.)/2
\] 
and 
\[
\ydissip(!.): =
 @b_1'(\check @a_2 (3@r(!.))+ \check @a_1(3@r(!.)))/2) \,
\check \ydissip(!.).
\]
This completes the construction.
\epr

\newpage
\subsection{Proof of expo decay Lf}

\bpr
Assume that system
$@.{x}(t)=f(x(t), \uc(t),\ddfun(t)), \; y(t)=h(x(t))$
admits a UIOSS-Lyapunov function with
 $@a_i$
($i=1,2$) as in $@a_1(\abs{@x}) !<V(@x) !<@a_2 (\abs{@x})$
 and with $@a, \ugdissip, \ygdissip$ as
in~\rref{pqUIOSS-dissip}. 
 Replacing $@a$ by $@a\circ@a_2^{-1}$, we
have
\be{DVioss1}
 \frac{d}{dt}V(x(t, @x,\uc))\,\le\, -@a(V(x(t, @x,\uc))) +
 \ugdissip(\abs{\uc(t)}) + \ygdissip(\abs{y(t,@x,\uc)})
\ee
for almost all $t \in[0, \; \tmax(@x, u))$.
According to Lemma~12 in Praly-Wang, $\exists$ 
\function $@r\in\ki$ which can be extended as a $C^1$-function to a
neighborhood of $[0, \, \infty)$ s.t. $@r'(r)\frac{@a(r)}{2}
\ge @r(r)$ $\forall$ $r\ge 0$.  Consider $W(@x) :=
@r(V(@x))$.  Observe that $W$ is again proper and positive
definite.
Along any trajectory $x(t):= x(t, @x, \uc)$ (with $y(t):=y(t,@x, \uc)$),
at any point where~\rref{pqUIOSS-dissip} holds, one has
that $\frac{d}{dt}W(x(t))=@r'(V(x(t)))\frac{d}{dt}V(x(t))$ is upper bounded by
\[
-@r'(V(x(t)))\frac{@a(V(x(t)))}{2}+@r'(V(x(t))) !.
\]
\[
\left(-\frac{@a(V(x(t)))}{2}+ \ugdissip(\abs{\uc(t)}) + \ygdissip(\abs{y(t)})\right),
\]
which in turn is bounded by
\begin{equation}
\label{e-51}
-@r(V(x(t)) +
@r'(V(x(t)))\left(-\frac{@a(V(x(t)))}{2}
+ \ugdissip(\abs{\uc(t)}) + \ygdissip(\abs{y(t)})\right) \,.
\end{equation}
Observe that when
$V(x(t)) \ge @a^{-1}(2\ugdissip(\abs{\uc(t)}) + 2\ygdissip(\abs{y(t)}))$
it holds that
\be{e-52}
@r'(V(x(t))) \left(-\frac{@a(V(x(t)))}{2 }
+ \ugdissip(\abs{\uc(t)}) + \ygdissip(\abs{y(t)})\right)\le 0,
\ee
while if instead
$V(x(t))\le @a^{-1}(2\ugdissip(\abs{\uc(t)}) + 2\ygdissip(\abs{y(t)}))$, 
$\therefore$
\be{e-53}
@r'(V(x(t)))\left(
\ugdissip(\abs{\uc(t)}) + \ygdissip(\abs{y(t)})\right)
\;!<\; \hat\ugdissip(\abs{\uc(t)}) + \hat\ygdissip(\abs{y(t)})
\ee
for some $\ki$-functions $\hat\ugdissip$ and $\hat\ygdissip$
(using here the fact that $@r'(s)$ is a continuous function).
Combining~\rref{e-52}
and~\rref{e-53}, one concludes
{}from the estimate~\rref{e-51} on $\frac{d}{dt}W(x(t))$ that
\[
\frac{d}{dt}W(x(t)) \,\le\,
-W(x(t)) + \hat\ugdissip(\abs{\uc(t)}) +\hat\ygdissip(\abs{y(t)}) 
\]
for almost all $t!e[0, \tmax)$.
\epr


}\edo

%\bibliographystyle{siam}
%\bibliography{ksw}
\begin{thebibliography}{10}

\bibitem{ALS}
{\sc F. Albertini and E. D. Sontag}, {\em Continuous control-Lyapunov functions
  for asymptotically controllable time-varying systems}, Internat. J.
 Control, 72 (1999), pp. 1630--1641.

\bibitem{intrinsic}
{\sc D. Angeli}, {\em Intrinsic robustness of global asymptotic stability},
  Systems Control Lett., 38 (1999), pp. 297--307.

\bibitem{AS}
{\sc D. Angeli and E. D. Sontag}, {\em Forward completeness, unboundedness
  observability, and their Lyapunov characterizations}, Systems
  Control Lett., 38 (1999), pp. 209--217.

\bibitem{ASW}
{\sc D. Angeli, E. D. Sontag, and Y. Wang}, {\em A characterization of integral
  input to state stability}, IEEE Trans. Automat. Control, 45 (2000),
  pp. 1082--1097.

\bibitem{Artstein}
{\sc Z. Artstein}, {\em Relaxed controls and the dynamics of control systems},
  SIAM J. Control Optim., 16 (1978), pp. 689--701.

\bibitem{boothby}
{\sc W. M. Boothby}, {\em An Introduction to Differentiable Manifolds and
  Riemannian Geometry}, 2nd ed., Academic Press, Orlando, FL, 1986.

\bibitem{Teel}
{\sc P. D. Christofides and A. Teel}, {\em Singular perturbations and
  input-to-state stability}, IEEE Trans. Automat. Control, 41 (1996),
  pp. 1645--1650.

\bibitem{CLSW}
{\sc F. Clarke, Y. Ledyaev, R. Stern, and P. Wolensky}, {\em Nonsmooth Analysis
  and Control Theory}, Springer-Verlag, New York, 1998.

\bibitem{CLSS}
{\sc F. Clarke, Y. S. Ledyaev, E. Sontag, and A. Subbotin}, {\em Asymptotic
  controllability implies feedback stabilization}, IEEE Trans.
  Automat. Control, 42 (1997), pp. 1394--1407.

\bibitem{HJ}
{\sc J. W. Helton and M. R. James}, {\em Extending $H_\infty$ Control to Nonlinear
  Systems}, SIAM, Philadelphia, 1999.

\bibitem{morse-iss}
{\sc J. Hespanha and A. Morse}, {\em Certainty equivalence implies
  detectability}, Systems Control Lett., 36 (1999), pp. 1--13.

\bibitem{hill-moylan}
{\sc D. Hill and P. Moylan}, {\em Dissipative dynamical systems: Basic
  input-output and state properties}, J. Franklin Institute, 5 (1980),
  pp. 327--357.

\bibitem{hill}
{\sc D. J. Hill and P. Moylan}, {\em Dissipative nonlinear systems: Basic
  properties and stability analysis}, in Proc. CDC 1992, pp. 3259--3264.

\bibitem{hu}
{\sc X. M. Hu}, {\em On state observers for nonlinear systems}, Systems
Control Lett., 17 (1991), pp. 645--473.

\bibitem{isi-iss}
{\sc A. Isidori}, {\em Global almost disturbance decoupling with stability for
  non-minimum-phase single-input single-output nonlinear systems}, Systems
 Control Lett., 28 (1996), pp. 115--122.

\bibitem{isi-newbook}
{\sc A. Isidori}, {\em Nonlinear Control
  Systems {\rm II}}, Springer-Verlag, London, 1999.

\bibitem{James}
{\sc M. James}, {\em A partial differential inequality for dissipative
  nonlinear systems}, Systems Control Lett., 21 (1993),
  pp. 315--320.

\bibitem{Jiang-Praly1}
{\sc Z.-P. Jiang and L. Praly}, {\em Preliminary results about robust
  Lagrange stability in adaptive nonlinear regulation}, Internat. J.
 Control, 6 (1992), pp. 285--307.

\bibitem{JTP}
{\sc Z.-P. Jiang, A. Teel, and L. Praly}, {\em Small-gain theorem for ISS
  systems and applications}, Math. Control Signals Systems, 7
  (1994), pp. 95--120.

\bibitem{Khalil95}
{\sc H. K. Khalil}, {\em Nonlinear Systems}, 2nd ed., Prentice-Hall, Upper Saddle River,
  NJ, 1996.

\bibitem{ksw-cdc98}
{\sc M. Krichman and E. D. Sontag}, {\em A version of a converse Lyapunov
  theorem for input-output to state stability}, in Proceedings of the 37th IEEE
  Conference on Decision and Control, Tampa, FL, 1998, IEEE, Piscataway, NJ, 1998,
  pp. 4121--4126.

\bibitem{krstic-deng}
{\sc M. Krsti\'{c} and H. Deng}, {\em Stabilization of Uncertain Nonlinear
  Systems}, Springer-Verlag, London, 1998.

\bibitem{krstic-book}
{\sc M. Krsti\'{c}, I. Kanellakopoulos, and P. V. Kokotovi\'{c}}, {\em
  Nonlinear and Adaptive Control Design}, John Wiley and Sons, New York, 1995.

\bibitem{liberzon}
{\sc D. Liberzon, E. D. Sontag, and Y. Wang}, {\em On integral-input-to-state
  stabilization}, in Proceedings of the 1999 American Control Conference, San Diego, 
  1999, pp. 1598--1602.

\bibitem{LSW}
{\sc Y. Lin, E. D. Sontag, and Y. Wang}, {\em A smooth converse Lyapunov
  theorem for robust stability}, SIAM J. Control Optim., 34
  (1996), pp. 124--160.

\bibitem{lu1}
{\sc W. M. Lu}, {\em A class of globally stabilizing controllers for nonlinear
  systems}, Systems Control Lett., 25 (1995), pp. 13--19.

\bibitem{lu2}
{\sc W. M. Lu}, {\em A state-space
  approach to parameterization of stabilizing controllers for nonlinear
  systems}, IEEE Trans. Automat. Control, 40 (1995), pp. 1576--1588.

\bibitem{Marino}
{\sc R. Marino and P. Tomei}, {\em Nonlinear output feedback tracking with
  almost disturbance decoupling}, IEEE Trans. Automat. Control, 44
  (1999), pp. 18--28.

\bibitem{MPD}
{\sc F. Mazenc, L. Praly, and W. Dayawansa}, {\em Global stabilization by
  output feedback: Examples and counterexamples}, Systems Control
  Lett., 23 (1994), pp. 119--125.

\bibitem{Morse}
{\sc A. S. Morse}, {\em Control using logic-based switching}, in Trends in
  Control: A European Perspective, A. Isidori, ed., Springer-Verlag, London, 1995,
  pp. 69--114.

\bibitem{Pan}
{\sc D. J. Pan, Z. Z. Han, and Z. J. Zhang}, {\em Bounded-input-bounded-output
  stabilization of nonlinear systems using state detectors}, Systems
  Control Lett., 21 (1993), pp. 189--198.

\bibitem{PW}
{\sc L. Praly and Y. Wang}, {\em Stabilization in spite of matched unmodelled
  dynamics and an equivalent definition of input-to-state stability},
  Math. Control Signals Systems, 9 (1996), pp. 1--33.

\bibitem{Sepulchre}
{\sc R. Sepulchre, M. Jankovic, and P. V. Kokotovi\'{c}}, {\em Integrator
  forwarding: A new recursive nonlinear robust design}, Automatica J. IFAC, 33 (1997),
  pp. 979--984.

\bibitem{EDS-old}
{\sc E. D. Sontag}, {\em Conditions for abstract nonlinear regulation},
  Inform. Control, 51 (1981), pp. 105--127.

\bibitem{coprime}
{\sc E. D. Sontag}, {\em Smooth stabilization
  implies coprime factorization}, IEEE Trans. Automat. Control, 34
  (1989), pp. 435--443.

\bibitem{coprime-cdc}
{\sc E. D. Sontag}, {\em Some connections
  between stabilization and factorization}, in Proceedings of the  IEEE Conference on Decision and
  Control, Tampa, FL, 1989, IEEE, Piscataway, NJ, 1989, pp. 990--995.

\bibitem{sontag-ejc}
{\sc E. D. Sontag}, {\em On the
  input-to-state stability property}, Euro. J. Control, 1 (1995),
  pp. 24--36.

\bibitem{iiss-eds}
{\sc E. D. Sontag}, {\em Comments on integral
  variants of input-to-state stability}, Systems Control Lett., 34
  (1998), pp. 93--100.

\bibitem{mct}
{\sc E. D. Sontag}, {\em Mathematical Control
  Theory: Deterministic Finite Dimensional Systems}, 2nd ed., Springer-Verlag, New York,
  1998.

\bibitem{NFSR}
{\sc E. D. Sontag}, {\em Stability and
  stabilization: Discontinuities and the effect of disturbances}, in Proceedings of
  the NAYO ASI Nonlinear Analysis, Differential Equations and Control,
  Montreal, 1998, F. Clarke and R. Stern, eds., Kluwer, Dordrecht, The Netherlands, 1999,
  pp. 551--598.

\bibitem{00paris}
{\sc E. D. Sontag}, {\em The ISS philosophy  
as a unifying framework for stability-like behavior}, in Nonlinear Control in
  the Year 2000, A. Isidori, F. Lamnabhi-Lagarrigue, and W. Respondek, eds.,
  Lecture Notes in Control and Inform. Sci. 258, Springer-Verlag, Berlin,
  2000.

\bibitem{ywscl}
{\sc E. D. Sontag and Y. Wang}, {\em On characterizations of the input-to-state
  stability property}, Systems Control Lett., 24 (1995),
  pp. 351--359.

\bibitem{ref8}
{\sc E. D. Sontag and Y. Wang}, {\em Detectability of
  nonlinear systems}, in Proceedings of the Conference on Information Science and Systems (CISS
  96), Princeton, NJ, 1996, pp. 1031--1036.

\bibitem{iss-new}
{\sc E. D. Sontag and Y. Wang}, {\em New
  characterizations of the input to state stability property}, IEEE
  Trans. Automat. Control, 41 (1996), pp. 1283--1294.

\bibitem{oss-scl}
{\sc E. D. Sontag and Y. Wang}, {\em Output-to-state
  stability and detectability of nonlinear systems}, Systems Control
  Lett., 29 (1997), pp. 279--290.

\bibitem{ios}
{\sc E. D. Sontag and Y. Wang}, {\em Notions of input to
  output stability}, Systems Control Lett., 38 (1999), pp. 351--359.

\bibitem{lios}
{\sc E. D. Sontag and Y. Wang}, {\em Lyapunov
  characterizations of input to output stability}, SIAM J. Control
  Optim., 39 (2000), pp. 226--249.

\bibitem{tsinias-scl}
{\sc J. Tsinias}, {\em Sontag's ``input to state stability condition'' and
  global stabilization using state detection}, Systems Control
  Lett., 20 (1993), pp. 219--226.

\bibitem{Tsinias}
{\sc J. Tsinias}, {\em Input to state
  stability properties of nonlinear systems and applications to bounded
  feedback stabilization using saturation}, ESAIM Control Optim. Calc. Var.,
  2 (1997), pp. 57--85.

\bibitem{vdS}
{\sc A. van der Schaft}, {\em ${L}^2$-gain analysis of nonlinear systems and
  nonlinear state feedback ${H}^\infty$-control}, IEEE Trans.
  Automat. Control, 37 (1992), pp. 770--784.

\bibitem{Warga}
{\sc J. Warga}, {\em Optimal Control of Differential and Functional Equations},
  Academic Press, New York, 1972.

\bibitem{willems-dissipation}
{\sc J. C. Willems}, {\em Mechanisms for the stability and instability in
  feedback systems}, Proc.  IEEE, 64 (1976), pp. 24--35.

\end{thebibliography}

\end{document}

{udsystem}
 $@.{x}(t)=f(x(t), \uc(t),\ddfun(t)), \; y(t)=h(x(t))$

{eq:UIOSS}
$\abs{x(t, @x, \uc, \ddfun)} !<\max \left\{ @b(\abs{@x}, t), 
\ulevel\left(\norma{\uc}{[0, t]}\right),
\ylevel\left(\norma{y}{[0, t]}\right)\right\}$

{ROSSsystem}
$@.{x}(t)= g(x(t),\dfun(t)) = f(x(t), \dufun(t)\sm(\abs{x(t)}),\ddfun(t))$

{UOSS-dissip-smooth}
$\nabla V (x) !.f(x, d) !< -\xdissip(\abs x) + \ydissip(\abs{h(x)})$

{dsystem}
$@.{x}=f(x(t), \dfun(t)), \, y(t)=h(x(t))$

{eq:UOSS}
$\abs{x(t, @x, \dfun)} !<\max \left\{ @b(\abs{@x}, t), \ylevel\left(\norma{y}{[0, t]}\right)\right\}$

{prop:UO}
UO iff $\exists$
$@r_1, @c_1, @c_2 !e\kk$ and $c!>0$ s.t.:
% the following implication holds: 
\beqn
\abso{h(x(t, @x, \dfun))} &\leq& @r_1(\abs{x(t, @x, \dfun)}) \quad 
\forall\,t !e[0, T]\nonumber\\[3mm]
& =>&  \abs{x(t, @x, \dfun)} !<@c_1(t) + @c_2(\abs @x) + c
 \quad \forall t !e[0, T]
\eeqn
$\forall$ $@x!e\states$, $\dfun !e\M_@W$, $T !e[0, \tmax(@x, \dfun))$

{eq:UO}
\beqn
\abso{h(x(t, @x, \dfun))} &\leq& @r_1(\abs{x(t, @x, \dfun)}) \quad 
\forall\,t !e[0, T]\nonumber\\[3mm]
& =>&  \abs{x(t, @x, \dfun)} !<@c_1(t) + @c_2(\abs @x) + c
 \quad \forall t !e[0, T]
\eeqn

{iiUOSS}
$\exists$ \functions $\iy$, $\iini$ $!e\kk$, $\ix !e \ki$ s.t.
%
{eq:iiUOSS}
$\int_0^t \ix(\abs{x(s, @x, \dfun)})\,ds \psp !<\psp \iini(\abs@x) + 
          \int_0^t \iy (\abs{h(x(s, @x, \dfun))})\,ds$
$\forall$ initial state $@x$, disturbance $\dfun !e\M_@W$, and
time $t !e[0, \tmax(@x, \dfun))$

{prop:newprop}
$\ve$-$@d$ characterization of \newprop\

{main_thm:UOSS}
\item[$1$.]  $@S$ is UOSS. \label{item1:UOSS}
\item[$2$.]  $@S$ is \newprop. \label{item2:UOSS}
\item[$3$.]  $@S$ is iiUOSS and UO. \label{item3:UOSS}
\item[$4$.]  $@S$ admits a UOSS-Lyapunov function. \label{item4:UOSS}

$\io{@x, \dfun}= \inf\{t !e[0, \tmax(@x, \dfun)): 
\abs{x(t,@x, \dfun )} !<@r(\abs{y(t,@x, \dfun)})\}$

\newprop:
$\exists$ $@r!e\ki, \;\klfun!e\kl$ s.t.:
$\forall$   $@x!e\states$, $\dfun !e\M_@W$, and any $T < \tmax(@x, \dfun)$,
if 
\[
\abs{x(t, @x, \dfun)} !>@r(\abs{h(x(t, @x, \dfun))}) 
\,(\forall \; 0 !<t !<T)
\;=>\;
\abs{x(t, @x, \dfun)} !<\klfun (\abs{@x}, t) \; (\forall\; 0 !<t!<T)
\]
